\oldpage[IV-3]{116}
\section{Propriétés topologiques des morphismes plats de présentation finie. Critères de platitude}
\label{section:IV.11-fr}

Alors que dans le \S~2 nous avons considéré les énoncés concernant la platitude
qui ne dépendent d'aucune hypothèse de finitude, et que le \S~6 étudie la notion
de platitude dans le cadre des préschémas localement noethériens (mais sans
hypothèse de finitude sur les morphismes), le présent paragraphe est consacré à
la notion de $f$-platitude dans le cas où le morphisme $f:X\to Y$ est
\emph{localement de présentation finie}. L'intérêt de la notion de morphisme
plat de présentation finie tient au fait que c'est elle qui semble exprimer
techniquement de la façon la plus adéquate la notion intuitive de «~famille de
préschémas algébriques paramétrée par un schéma $Y$~», dont l'étude est un des
objets principaux de la Géométrie algébrique. D'ailleurs, alors même qu'on ne
s'intéresserait au départ qu'au cas d'un schéma de base noethérien, il est
indispensable, pour certaines raisons techniques (par exemple pour certaines
applications de la théorie de la «~descente~», qui conduit à introduire des
schémas non nécessairement noethériens) de ne pas se limiter à ce cas, dès
qu'il s'agit de problèmes de nature essentiellement relative liés aux
morphismes localement de présentation finie. Nous suivrons systématiquement ce
principe, déjà étayé par les résultats des \S\S~8 et 9, dans toute la suite de ce
Chapitre, et même dans la suite de notre Traité, quitte à lui sacrifier à
l'occasion la simplicité de certaines démonstrations, que des hypothèses
noethériennes permettent parfois d'alléger\footnote{Ce principe s'inspire
également de la nécessité de donner droit de cité, comme «~espaces de
paramètres~» pour les familles de schémas algébriques, à des espaces annelés
(et même des «~topos~» annelés) quelconques, pour lesquels il ne peut plus être
question en général d'hypothèses noethériennes. Il semble assez clair que l'on
ne pourra plus éluder longtemps cette nouvelle extension de la Géométrie
algébrique, et il convient dès à présent de développer les notions et
techniques de nature «~relative~» de la théorie des schémas de sorte qu'elles
puissent s'adapter pratiquement telles quelles à ce cadre plus général.}. Dans
le présent paragraphe, cela nous conduit à reprendre, dans le contexte de la
«~présentation finie~» (notamment au n$^\circ$~3) certains énoncés de platitude,
déjà obtenus dans le contexte noethérien. L'outil technique essentiel pour
faire la réduction au cas noethérien est le théorème de compatibilité de la
platitude avec les limites projectives de préschémas (11.2.6), complétant les
résultats généraux du \S~8. Nous prouvons aussi en passant (11.3.1) un résultat
souvent utilisé

\oldpage[IV-3]{117}
par la suite, impliquant que l'ensemble des points de platitude d'un morphisme
localement de présentation finie est ouvert.

Dans les n$^\mathrm{os}$~4 à 8, nous étudions la question de la «~descente~» de
la platitude, consistant à trouver des conditions utiles sur un morphisme de
changement de base $Y'\to Y$ (non plat en général) pour pouvoir conclure que si
$X\times_YY'$ est plat sur $Y'$, $X$ est plat sur $Y$. Ces résultats, plus
techniques que ceux des n$^\mathrm{os}$~1 à 3, sont d'une utilisation moins
fréquente dans la suite~; ils joueront cependant un rôle important dans les
techniques de construction non projectives, au chapitre suivant. Le seul
résultat des n$^\mathrm{os}$~4 à 8 utilisé dans la suite du chap.~IV est le
critère valuatif de platitude (11.8), qu'on appliquera dans (15.2).

Enfin, les n$^\mathrm{os}$~9 et 10 sont consacrés à l'étude d'une notion qui
précise, en théorie des schémas, celle de densité au sens topologique, savoir la
notion de famille de sous-préschémas \emph{schématiquement dense} dans un
préschéma donné, et notamment l'étude du comportement de cette notion par
changement de base (plat ou quelconque). Cette notion est surtout utilisée,
pour l'instant, dans l'étude des schémas en groupes.

\subsection{Ensembles de platitude (cas noethérien)}
\label{subsection:IV.11.1-fr}

\begin{theorem}[11.1.1]
\label{IV.11.1.1-fr}
Soient $Y$ un préschéma localement noethérien, $f:X\to Y$ un morphisme
localement de type fini, $\mathcal F$ un $\mathcal O_X$-Module cohérent. Alors
l'ensemble $U$ des $x\in X$ tels que $\mathcal F$ soit $f$-plat au point $x$
est ouvert dans $X$.
\end{theorem}

La question étant évidemment locale sur $X$ et $Y$, on peut supposer
$Y=\operatorname{Spec}(A)$, $X=\operatorname{Spec}(B)$, $A$ étant noethérien et
$B$ une $A$-algèbre de type fini. On a alors $\mathcal F=\widetilde M$, où $M$
est un $B$-module de type fini. Appliquons le critère
($\mathbf 0_{\mathrm{III}}$, 9.2.6)~: il suffit donc de démontrer l'assertion
suivante~:

\begin{env}[11.1.1.1]
\label{IV.11.1.1.1-fr}
\emph{Soient $A$ un anneau noethérien, $B$ une $A$-algèbre de type fini, $M$ un
$B$-module de type fini, $\mathfrak q$ un idéal premier de $B$, $\mathfrak p$
son image réciproque dans $A$. On suppose que $M_{\mathfrak q}$ soit un
$A_{\mathfrak p}$-module plat. Alors il existe $g\in B-\mathfrak q$ tel que,
pour tout idéal premier $\mathfrak q'\supset\mathfrak q$ de $B$ tel que
$g\notin\mathfrak q'$, $M_{\mathfrak q'}$ soit un
$A_{\mathfrak p'}$-module plat, en désignant par $\mathfrak p'$ l'image
réciproque de $\mathfrak q'$ dans $A$ (il revient d'ailleurs au même
($\mathbf 0_{\mathrm I}$, 6.3.1) de dire que $M_{\mathfrak q'}$ est un
$A$-module plat).}
\end{env}

Considérons pour cela $B_{\mathfrak q'}$ comme une $A$-algèbre~; on a
évidemment
$\mathfrak pB_{\mathfrak q'}\subset\mathfrak q'B_{\mathfrak q'}$. On sait
alors ($\mathbf 0_{\mathrm{III}}$, 10.2.2) que pour que $M_{\mathfrak q'}$ soit
un $A$-module plat, il faut et il suffit que
$M_{\mathfrak q'}/\mathfrak pM_{\mathfrak q'}$ soit un
$(A/\mathfrak p)$-module plat et que l'on ait
$\operatorname{Tor}_1^A(M_{\mathfrak q'},A/\mathfrak p)=0$. Or, on a
$M_{\mathfrak q'}=M\otimes_BB_{\mathfrak q'}$~; $B_{\mathfrak q'}$ étant plat
sur $B$, on a
\[
  M_{\mathfrak q'}/\mathfrak pM_{\mathfrak q'}
  =(M/\mathfrak pM)_{\mathfrak q'}
\]
et
\[
  \operatorname{Tor}_1^A(M_{\mathfrak q'},A/\mathfrak p)
  =(\operatorname{Tor}_1^A(M,A/\mathfrak p))_{\mathfrak q'}
\]
(en définissant le Tor à l'aide d'une résolution projective de
$A/\mathfrak p$)~; pour la même raison, comme on doit avoir
$g\notin\mathfrak q'$, $B_{\mathfrak q'}$ est plat sur $B_g$, donc
\[
  (M/\mathfrak pM)_{\mathfrak q'}
  =((M/\mathfrak pM)_g)_{\mathfrak q'B_g}
\]
et
\[
  (\operatorname{Tor}_1^A(M,A/\mathfrak p))_{\mathfrak q'}
  =((\operatorname{Tor}_1^A(M,A/\mathfrak p))_g)_{\mathfrak q'B_g},
\]
où dans ces formules, $M/\mathfrak pM$ et
$\operatorname{Tor}_1^A(M,A/\mathfrak p)$ sont considérés comme des
$B$-modules. Compte tenu de ($\mathbf 0_{\mathrm I}$, 6.3.2), on voit qu'on
est ramené à prouver le

\begin{lemma}[11.1.1.2]
\label{IV.11.1.1.2-fr}
Sous les conditions de (11.1.1.1), il existe $g\in B-\mathfrak q$ tel que~:
\begin{enumerate}[label=(\roman*)]
  \item $(M/\mathfrak pM)_g$ soit un $(A/\mathfrak p)$-module plat~;
  \item $(\operatorname{Tor}_1^A(M,A/\mathfrak p))_g=0$.
\end{enumerate}
\end{lemma}

\oldpage[IV-3]{118}

En vertu du théorème de platitude générique (6.9.1) appliqué à l'anneau
intègre $A/\mathfrak p$, à la $(A/\mathfrak p)$-algèbre de type fini
$B/\mathfrak pB$ et au $(B/\mathfrak pB)$-module de type fini
$M/\mathfrak pM$, il existe un $h\in A-\mathfrak p$ tel que, si $\bar h$ est
son image canonique dans $A/\mathfrak p$,
$(M/\mathfrak pM)_{\bar h}$ soit un $(A/\mathfrak p)$-module plat. D'autre
part, comme $M_{\mathfrak q}$ est un $A_{\mathfrak p}$-module plat, et par
suite un $A$-module plat ($\mathbf 0_{\mathrm I}$, 6.3.1), on a
$\operatorname{Tor}_1^A(M_{\mathfrak q},A/\mathfrak p)=0$, ce qui s'écrit
aussi
$(\operatorname{Tor}_1^A(M,A/\mathfrak p))_{\mathfrak q}=0$. Mais comme $A$
et $B$ sont noethériens, $\operatorname{Tor}_1^A(M,A/\mathfrak p)$ est un
$B$-module de type fini, donc ($\mathbf 0_{\mathrm I}$, 5.2.2) il existe
$g\in B-\mathfrak q$ tel que
$(\operatorname{Tor}_1^A(M,A/\mathfrak p))_g=0$. D'ailleurs, on a
$(M/\mathfrak pM)_{\bar h}=(M/\mathfrak pM)_h$ ($M/\mathfrak pM$ étant
considéré dans le second membre comme un $A$-module)~; de plus,
$(M/\mathfrak pM)_h$ étant un $B$-module,
$(M/\mathfrak pM)_{hg}$ est encore un $(A/\mathfrak p)$-module plat, car il
s'écrit $(M/\mathfrak pM)_{\overline{hg}}$, où $\overline{hg}$ est l'image
canonique de $hg$ dans $B/\mathfrak pB$, et il suffit d'appliquer
($\mathbf 0_{\mathrm I}$, 6.3.2)~; enfin, on a
$(\operatorname{Tor}_1^A(M,A/\mathfrak p))_{hg}=0$ et
$hg\in B-\mathfrak q$ puisque $h\in A-\mathfrak p$. C.Q.F.D.

\begin{corollary}[11.1.2]
\label{IV.11.1.2-fr}
Soient $Y$ un préschéma localement noethérien, $f:X\to Y$ un morphisme
localement de type fini, $\mathcal F$, $\mathcal F'$ deux
$\mathcal O_X$-Modules cohérents,
$u:\mathcal F'\to\mathcal F$ un homomorphisme de $\mathcal O_X$-Modules. On
suppose que $\mathcal F$ est $f$-plat. Alors l'ensemble $U$ des $x\in X$ tels
que, en posant $y=f(x)$, l'homomorphisme
\[
  u_x\otimes1:\mathcal F'_x\otimes_{\mathcal O_y}k(y)\longrightarrow
  \mathcal F_x\otimes_{\mathcal O_y}k(y)
\]
soit injectif, est ouvert dans $X$.
\end{corollary}

En effet, soit $\mathcal N$ (resp. $\mathcal P$) le noyau (resp. le conoyau) de
$u$~; appliquons ($\mathbf 0_{\mathrm{III}}$, 10.2.4)\footnote{Voici une
démonstration de ($\mathbf 0_{\mathrm{III}}$, 10.2.4), qui n'a pas été publiée
dans l'\emph{Algèbre commutative} de N.~Bourbaki. Compte tenu de
($\mathbf 0_{\mathrm I}$, 6.1.2), il suffit de voir que b) entraîne a). Posons
$P=\operatorname{Im}(u)$, $Q=\operatorname{Coker}(u)$,
$R=\operatorname{Ker}(u)$. Le composé
$M\otimes k\to P\otimes k\to N\otimes k$ est injectif, et
$M\otimes k\to P\otimes k$ est surjectif, donc
$P\otimes k\to N\otimes k$ est injectif et
$M\otimes k\to P\otimes k$ est bijectif. La suite exacte
$0\to P\to N\to Q\to0$ donne la suite exacte
\[
0=\operatorname{Tor}_1^A(N,k)\to\operatorname{Tor}_1^A(Q,k)
\to P\otimes k\to N\otimes k\to Q\otimes k\to0
\]
et puisque $P\otimes k\to N\otimes k$ est injectif, on a
$\operatorname{Tor}_1^A(Q,k)=0$~; comme $Q$ est un $B$-module de type fini,
($\mathbf 0_{\mathrm{III}}$, 10.2.2) montre que $Q$ est un $A$-module plat~;
alors $P$ est aussi un $A$-module plat par ($\mathbf 0_{\mathrm I}$, 6.1.2).
La suite $0\to R\to M\to P\to0$ étant exacte, il en est de même de
$0\to R\otimes k\to M\otimes k\to P\otimes k\to0$ par
($\mathbf 0_{\mathrm I}$, 6.1.2)~; puisque $M\otimes k\to P\otimes k$ est
bijectif, on a $R\otimes k=0$~; mais comme $B$ est noethérien, $R$ est un
$B$-module de type fini, donc on a $R=0$ en vertu du lemme de Nakayama.} aux
anneaux locaux $\mathcal O_y$ et $\mathcal O_x$ et aux
$\mathcal O_x$-modules $\mathcal F'_x$ et $\mathcal F_x$~: dire que
$u_x\otimes1$ est injectif équivaut à dire que $\mathcal N_x=0$ et que
$\mathcal P$ est $f$-plat au point $x$. Or comme $\mathcal N$ et $\mathcal P$
sont cohérents ($\mathbf 0_{\mathrm I}$, 5.3.4), l'ensemble des $x$ où
$\mathcal N_x=0$ est ouvert ($\mathbf 0_{\mathrm I}$, 5.2.2), et l'ensemble
des $x$ où $\mathcal P$ est $f$-plat est ouvert par (11.1.1)~; d'où la
conclusion.

En particulier~:

\begin{corollary}[11.1.3]
\label{IV.11.1.3-fr}
Soient $Y$ un préschéma localement noethérien, $f:X\to Y$ un morphisme plat et
localement de type fini, $g$ une section de $\mathcal O_X$ au-dessus de $X$.
L'ensemble des $x\in X$ tels que $g_x\otimes1$ soit non diviseur de zéro dans
$\mathcal O_x\otimes_{\mathcal O_{f(x)}}k(f(x))$ est ouvert dans $X$.
\end{corollary}

Il suffit d'appliquer (11.1.2) à l'endomorphisme de $\mathcal O_X$ défini par
$g$.

\begin{corollary}[11.1.4]
\label{IV.11.1.4-fr}
Soient $Y$ un préschéma localement noethérien, $f:X\to Y$ un morphisme
localement de type fini, $\mathcal F$ un $\mathcal O_X$-Module cohérent et
$f$-plat. Soit $(g_i)_{1\leq i\leq n}$ une suite de sections de $\mathcal O_X$
au-dessus de $X$. Alors l'ensemble $U$ des $x\in X$ tels que la suite
$((g_i)_x\otimes1)$ soit
$(\mathcal F_x\otimes_{\mathcal O_{f(x)}}k(f(x)))$-régulière est ouvert dans
$X$.
\end{corollary}

\oldpage[IV-3]{119}

Comme $\mathcal F$ est $f$-plat, il résulte de ($\mathbf 0$, 15.1.16) que $U$
est aussi l'ensemble des $x\in X$ tels que la suite $((g_i)_x)$ soit
$\mathcal F_x$-régulière et que le $\mathcal O_x$-module $\mathcal G_x$ (où
$\mathcal G=\mathcal F/(\sum_{i=1}^n g_i\mathcal F)$) soit $f$-plat. Mais
$\mathcal F$ et $\mathcal G$ sont cohérents, donc le corollaire résulte de
(11.1.1) et de ($\mathbf 0$, 15.2.4).

\begin{corollary}[11.1.5]
\label{IV.11.1.5-fr}
Soient $X$, $Y$, $Z$ trois préschémas localement noethériens,
$f:X\to Y$, $g:Y\to Z$ deux morphismes de type fini, $\mathcal F$ un
$\mathcal O_X$-Module cohérent. Alors l'ensemble $U$ des $y\in Y$ tels que,
pour toute générisation $y'$ de $y$, $\mathcal F$ soit $(g\circ f)$-plat en
tous les points de $f^{-1}(y')$ (i.e. tels que
$\mathcal F\otimes_Y\operatorname{Spec}(\mathcal O_{Y,y})$ soit plat
relativement au morphisme
$X\times_Y\operatorname{Spec}(\mathcal O_{Y,y})\to Z$) est ouvert dans $Y$.
\end{corollary}

Si $V$ est l'ensemble des $x\in X$ où $\mathcal F$ est $(g\circ f)$-plat, $U$
est l'ensemble des $y\in Y$ tels que pour toute générisation $y'$ de $y$, on
ait $f^{-1}(y')\subset V$. Or $V$ est ouvert (11.1.1) donc localement
constructible dans $X$, et l'ensemble $W$ des $y\in Y$ tels que
$f^{-1}(y)\subset V$ est égal à $Y-f(X-V)$, donc est aussi localement
constructible dans $Y$ en vertu du théorème de Chevalley (1.8.4). Mais il
résulte alors de ($\mathbf 0_{\mathrm{III}}$, 9.2.5) que les points de $U$ sont
les points intérieurs à $W$, d'où la conclusion.

\begin{corollary}[11.1.6]
\label{IV.11.1.6-fr}
Soient $A$ un anneau noethérien, $B$ une $A$-algèbre de type fini, $C$ une
$B$-algèbre de type fini, $M$ un $C$-module de type fini. Alors l'ensemble des
$\mathfrak q\in\operatorname{Spec}(B)$ tels que $M_{\mathfrak q}$ soit un
$A$-module plat est ouvert dans $\operatorname{Spec}(B)$.
\end{corollary}

Compte tenu de (2.1.2), c'est une conséquence de (11.1.5) appliqué à
$Z=\operatorname{Spec}(A)$, $Y=\operatorname{Spec}(B)$,
$X=\operatorname{Spec}(C)$, $\mathcal F=\widetilde M$.

Les résultats de ce numéro seront débarrassés des hypothèses noethériennes
dans (11.3).

\subsection{Platitude d'une limite projective de préschémas}
\label{subsection:IV.11.2-fr}

\begin{env}[11.2.1]
\label{IV.11.2.1-fr}
Soient $A$ un anneau, $M$, $N$ deux $A$-modules, $A'$ une $A$-algèbre~;
posons $M'=M\otimes_AA'$, $N'=N\otimes_AA'$. Rappelons
($\mathbf{III}$, 6.3.8) que l'on définit, pour tout $i$, un homomorphisme
canonique de $A$-modules
\begin{equation}
\psi_i:\operatorname{Tor}_i^A(M,N)\to\operatorname{Tor}_i^{A'}(M',N')
\label{IV.11.2.1.1-fr}
\tag{11.2.1.1}
\end{equation}
de la façon suivante~: on considère une résolution gauche de $M$ par des
$A$-modules libres
\begin{equation}
\cdots\longrightarrow L_{i+1}\xrightarrow{f_i}L_i\longrightarrow\cdots
\longrightarrow L_0\xrightarrow{\varepsilon}M\longrightarrow0
\label{IV.11.2.1.2-fr}
\tag{11.2.1.2}
\end{equation}
d'où l'on déduit par tensorisation avec $A'$ un complexe de $A'$-modules
\begin{equation}
\cdots\longrightarrow L'_{i+1}\xrightarrow{f'_i}L'_i\longrightarrow\cdots
\longrightarrow L'_0\xrightarrow{\varepsilon'}M'\longrightarrow0
\label{IV.11.2.1.3-fr}
\tag{11.2.1.3}
\end{equation}
où l'on a posé $L'_i=L_i\otimes_AA'$, $f'_i=f_i\otimes1_{A'}$,
$\varepsilon'=\varepsilon\otimes1_{A'}$. Considérons d'autre part une
résolution gauche de $M'$ par des $A'$-modules libres
\begin{equation}
\cdots\longrightarrow L''_{i+1}\xrightarrow{f''_i}L''_i\longrightarrow\cdots
\longrightarrow L''_0\xrightarrow{\varepsilon''}M'\longrightarrow0.
\label{IV.11.2.1.4-fr}
\tag{11.2.1.4}
\end{equation}
\end{env}

\oldpage[IV-3]{120}

Comme les $L'_i$ sont des $A'$-modules libres, on sait (M, V, 1.1) qu'il y a
des $A'$-homomorphismes $u_i$ formant un diagramme commutatif
\[
\xymatrix{
\cdots\ar[r]&L'_{i+1}\ar[r]^{f'_i}\ar[d]_{u_{i+1}}&
L'_i\ar[r]\ar[d]_{u_i}&\cdots\ar[r]&
L'_0\ar[r]^{\varepsilon'}\ar[d]_{u_0}&M'\ar[r]\ar[d]_{1_{M'}}&0\\
\cdots\ar[r]&L''_{i+1}\ar[r]_{f''_i}&L''_i\ar[r]&\cdots\ar[r]&
L''_0\ar[r]_{\varepsilon''}&M'\ar[r]&0.
}
\]
Si l'on compose l'homomorphisme $u_\bullet:L'_\bullet\to L''_\bullet$ de
complexes ainsi défini avec l'homomorphisme canonique
$L_\bullet\to L'_\bullet$, on obtient un homomorphisme de complexes de
$A$-modules $L_\bullet\to L''_\bullet$~; en remarquant que l'on a
$(L_i\otimes_AN)\otimes_AA'=L'_i\otimes_{A'}N'$, on en déduit un
homomorphisme de complexes de $A$-modules
$L_\bullet\otimes_AN\to L''_\bullet\otimes_{A'}N'$, d'où, en passant à
l'homologie, les homomorphismes canoniques (11.2.1.1). Comme les $u_i$ sont
bien déterminés à homotopie près (M, V, 1.1), les homomorphismes (11.2.1.1)
ne dépendent pas du choix des $u_i$ ni du choix des résolutions libres
$L_\bullet$ et $L''_\bullet$.

\begin{sloppypar}
Comme les $\operatorname{Tor}_i^{A'}(M',N')$ sont des $A'$-modules, on déduit
canoniquement de (11.2.1.1) des $A'$-homomorphismes
\begin{equation}
\begin{gathered}
\varphi_i:\operatorname{Tor}_i^A(M,N)\otimes_AA'\\[-2pt]
\longrightarrow\operatorname{Tor}_i^{A'}(M',N').
\end{gathered}
\label{IV.11.2.1.5-fr}
\tag{11.2.1.5}
\end{equation}
\end{sloppypar}

Considérons maintenant deux homomorphismes d'anneaux
\[
\rho:A\to A^{(1)},\qquad \sigma:A^{(1)}\to A^{(2)},
\]
et leur composé $\sigma\circ\rho:A\to A^{(2)}$~; posons
$M^{(j)}=M\otimes_AA^{(j)}$, $N^{(j)}=N\otimes_AA^{(j)}$ pour $j=1,2$.
Alors l'homomorphisme canonique composé
\[
\operatorname{Tor}_i^A(M,N)\to
\operatorname{Tor}_i^{A^{(1)}}(M^{(1)},N^{(1)})\to
\operatorname{Tor}_i^{A^{(2)}}(M^{(2)},N^{(2)})
\]
est le même que l'homomorphisme canonique déduit de $\sigma\circ\rho$~;
cela résulte de ce que, si $L_\bullet^{(j)}$ est une résolution libre de
$M^{(j)}$, le diagramme
\[
\xymatrix{
L_\bullet\ar[r]&L_\bullet\otimes_AA^{(1)}\ar[r]\ar[d]&
L_\bullet^{(1)}\ar[d]\\
&L_\bullet\otimes_AA^{(2)}\ar[r]&
L_\bullet^{(1)}\otimes_{A^{(1)}}A^{(2)}\ar[r]&L_\bullet^{(2)}
}
\]
est commutatif.

\begin{env}[11.2.2]
\label{IV.11.2.2-fr}
Les notations étant celles de (11.2.1), considérons maintenant un
\emph{système inductif filtrant} de $A$-algèbres
$(A_\alpha,g_{\beta\alpha})$, et pour tout indice $\alpha$, posons
$M_\alpha=M\otimes_AA_\alpha$, $N_\alpha=N\otimes_AA_\alpha$~; il résulte
alors de (11.2.1) que pour chaque $i$,
$(\operatorname{Tor}_i^{A_\alpha}(M_\alpha,N_\alpha),
\psi_{i\beta\alpha})$, où $\psi_{i\beta\alpha}$ est l'homomorphisme
canonique (11.2.1.1) correspondant à $g_{\beta\alpha}$, est un \emph{système
inductif} de $A_\alpha$-modules. Posons
$A'=\varinjlim A_\alpha$, $M'=\varinjlim M_\alpha=M\otimes_AA'$,
$N'=\varinjlim N_\alpha=N\otimes_AA'$~; si on désigne par
$g_\alpha:A_\alpha\to A'$ l'homomorphisme canonique, on en déduit des
homomorphismes canoniques (11.2.1.1)
$\psi_{i\alpha}:\operatorname{Tor}_i^{A_\alpha}(M_\alpha,N_\alpha)
\to\operatorname{Tor}_i^{A'}(M',N')$ qui (toujours en vertu de (11.2.1))
forment un système inductif d'homomorphismes~; nous nous proposons de
compléter le résultat de (M, V, 9.4$^*$) en montrant que les
\begin{equation}
\psi_i=\varinjlim_\alpha\psi_{i\alpha}:
\varinjlim_\alpha\operatorname{Tor}_i^{A_\alpha}(M_\alpha,N_\alpha)
\to\operatorname{Tor}_i^{A'}(M',N')
\label{IV.11.2.2.1-fr}
\tag{11.2.2.1}
\end{equation}
sont des \emph{isomorphismes de $A'$-modules}. Pour cela, nous procédons
comme dans (M, V, 9.5$^*$), en associant à chaque $M_\alpha$ sa
\emph{résolution libre canonique}. Tout revient (compte tenu de l'exactitude
du foncteur $\varinjlim$) à prouver le
\end{env}

\oldpage[IV-3]{121}

\begin{lemma}[11.2.2.2]
\label{IV.11.2.2.2-fr}
Soient $(A_\alpha,g_{\beta\alpha})$ un système inductif filtrant d'anneaux,
$(M_\alpha,h_{\beta\alpha})$ un système inductif d'ensembles,
$A'=\varinjlim A_\alpha$, $M'=\varinjlim M_\alpha$,
$g_\alpha:A_\alpha\to A'$, $h_\alpha:M_\alpha\to M'$ les applications
canoniques. Pour tout $\alpha$, soit $F(M_\alpha)$ le $A_\alpha$-module des
combinaisons linéaires formelles d'éléments de $M_\alpha$~; soit de même
$F(M')$ le $A'$-module des combinaisons linéaires formelles d'éléments de
$M'$~; si $h^0_{\beta\alpha}:F(M_\alpha)\to F(M_\beta)$ (pour
$\alpha\leq\beta$) et $h'_\alpha:F(M_\alpha)\to F(M')$ sont les
$A_\alpha$-homomorphismes déduits de $h_{\beta\alpha}$ et $h_\alpha$
respectivement, $(F(M_\alpha),h^0_{\beta\alpha})$ est un système inductif
de $A_\alpha$-modules et $(h'_\alpha)$ un système inductif
d'homomorphismes. Alors
\[
h''=\varinjlim h'_\alpha:\varinjlim F(M_\alpha)\to
F(M')=F(\varinjlim M_\alpha)
\]
est un isomorphisme.
\end{lemma}

Pour la démonstration, voir Bourbaki, \emph{Alg.}, chap.~II, 3$^e$ éd.,
\S~6, n$^\circ$~6, cor. de la prop.~10.

\begin{env}[11.2.3]
\label{IV.11.2.3-fr}
Reprenons les notations de (11.2.1) et considérons particulièrement le cas
$i=1$~; posons $T=\operatorname{Tor}_1^A(M,N)$, $T'=T\otimes_AA'$,
$T''=\operatorname{Tor}_1^{A'}(M',N')$~; alors $\psi_1$ est
l'homomorphisme
\[
\operatorname{Ker}(f_0\otimes1_N)/\operatorname{Im}(f_1\otimes1_N)
\to
\operatorname{Ker}(f''_0\otimes1_{N'})/
\operatorname{Im}(f''_1\otimes1_{N'})
\]
qui se déduit par passage aux quotients de la restriction
$\operatorname{Ker}(f_0\otimes1_N)\to
\operatorname{Ker}(f''_0\otimes1_{N'})$ de
\[
u_1\otimes1:L_0\otimes_AN\to L''_0\otimes_{A'}N'=L''_0\otimes_AN.
\]
Posons $R=\operatorname{Ker}(\varepsilon)$,
$R''=\operatorname{Ker}(\varepsilon'')$, de sorte que l'on a les suites
exactes
\[
0\to R\xrightarrow{j}L_0\xrightarrow{\varepsilon}M\to0
\qquad\text{et}\qquad
0\to R''\xrightarrow{j''}L''_0\xrightarrow{\varepsilon''}M'\to0,
\]
d'où on déduit les suites exactes d'homologie
\begin{equation}
0=\operatorname{Tor}_1^A(L_0,N)\to T\xrightarrow{\partial}R\otimes_AN
\xrightarrow{j\otimes1}L_0\otimes_AN\xrightarrow{\varepsilon\otimes1}
M\otimes_AN\to0
\label{IV.11.2.3.1-fr}
\tag{11.2.3.1}
\end{equation}
\begin{equation}
0=\operatorname{Tor}_1^{A'}(L''_0,N')\to T''\xrightarrow{\partial''}
R''\otimes_{A'}N'\xrightarrow{j''\otimes1}L''_0\otimes_{A'}N'
\xrightarrow{\varepsilon''\otimes1}M'\otimes_{A'}N'\to0.
\label{IV.11.2.3.2-fr}
\tag{11.2.3.2}
\end{equation}
On a d'autre part un homomorphisme de $A'$-modules
\[
v:R'=\operatorname{Ker}(\varepsilon)\otimes_AA'
\to\operatorname{Ker}(\varepsilon')\xrightarrow{u_0}
\operatorname{Ker}(\varepsilon'')=R''.
\]
\end{env}

\oldpage[IV-3]{122}

Montrons que le diagramme
\begin{equation}
\xymatrix{
&T'\ar[r]^{\partial\otimes1}\ar[d]_{\varphi_1}&
R'\otimes_{A'}N'\ar[r]\ar[d]_{v\otimes1}&
L'_0\otimes_{A'}N'\ar[r]\ar[d]_{u_0\otimes1}&
M'\otimes_{A'}N'\ar[r]\ar@{=}[d]&0\\
0\ar[r]&T''\ar[r]_{\partial''}&R''\otimes_{A'}N'\ar[r]&
L''_0\otimes_{A'}N'\ar[r]&M'\otimes_{A'}N'\ar[r]&0
}
\label{IV.11.2.3.3-fr}
\tag{11.2.3.3}
\end{equation}
est \emph{commutatif}. Pour cela, on vérifie aussitôt (M, IV, 1) que
l'homomorphisme $\partial:T\to R\otimes_AN$ provient (dans le cas actuel)
par passage au quotient de l'homomorphisme
$w:\operatorname{Ker}(f_0\otimes1_N)\to R\otimes_AN$, restriction de
l'homomorphisme $g_0\otimes1_N:L_1\otimes_AN\to R\otimes_AN$, où
$g_0:L_1\to R$ est tel que $f_0=j\circ g_0$~; de même pour $\partial''$.
Il suffit donc de voir que le diagramme
\[
\xymatrix{
\operatorname{Ker}(f_0\otimes1_N)\ar[r]^w\ar[d]&R\otimes_AN\ar[r]&
R'\otimes_{A'}N'=(R\otimes_AN)\otimes_AA'\ar[d]_{v\otimes1}\\
\operatorname{Ker}(f''_0\otimes1_{N'})\ar[rr]&&R''\otimes_{A'}N'
}
\]
est commutatif, ce qui résulte de la commutativité du diagramme
\[
\xymatrix{
L_1\ar[r]\ar[d]&R\ar[d]\\
L''_1\ar[r]&R''.
}
\]

\begin{lemma}[11.2.4]
\label{IV.11.2.4-fr}
Soient $A$ un anneau, $C$ une $A$-algèbre, $M$ un $A$-module, $N$ un
$C$-module, $A'$ une $A$-algèbre. Posons $C'=C\otimes_AA'$,
$M'=M\otimes_AA'$, $N'=N\otimes_AA'=N\otimes_CC'$. Supposons que
$M\otimes_AN$ soit un $C$-module plat. Alors l'homomorphisme canonique
\[
\varphi_1:\operatorname{Tor}_1^A(M,N)\otimes_AA'
\to\operatorname{Tor}_1^{A'}(M',N')
\]
(cf. (11.2.1.5)) est surjectif.
\end{lemma}

Gardons les notations de (11.2.3)~; l'exactitude à droite du produit
tensoriel montre que la suite
$L'_1\xrightarrow{f'_0}L'_0\xrightarrow{\varepsilon'}M'\to0$ est exacte~;
comme $L'_0$ et $L'_1$ sont des $A'$-modules libres, \emph{on peut supposer
que l'on a pris} $L'_0=L''_0$, $L'_1=L''_1$, $u_0$ et $u_1$ étant les
\emph{applications identiques} et $f''_0=f'_0$. Comme
$R=\operatorname{Im}(f_0)$ et $R''=\operatorname{Im}(f''_0)=
\operatorname{Im}(f'_0)$, l'homomorphisme $v$ est surjectif, et il en est
donc de même de $v\otimes1$. La démonstration sera achevée si l'on prouve que
la première ligne de (11.2.3.3) est \emph{exacte}, $v\otimes1$ étant
surjective et $u_0\otimes1$ injectif (Bourbaki, \emph{Alg. comm.},
chap.~I$^{\mathrm{er}}$, \S~1, n$^\circ$~4, cor.~2 de la prop.~2).
Utilisons pour cela

\oldpage[IV-3]{123}

l'hypothèse que $M\otimes_AN$ est un $C$-module plat. Posant
$Q=\operatorname{Ker}(\varepsilon\otimes1)=\operatorname{Im}(j\otimes1)$
dans la suite exacte (11.2.3.1), on a les deux suites exactes
$0\to Q\to L_0\otimes_AN\to M\otimes_AN\to0$ et
$T\to R\otimes_AN\to Q\to0$, où les homomorphismes sont des homomorphismes
de $C$-modules~; utilisant l'hypothèse de platitude, on en déduit
($\mathbf 0_{\mathrm I}$, 6.1.2) la suite exacte
\[
0\to Q\otimes_CC'\to(L_0\otimes_AN)\otimes_CC'
\to(M\otimes_AN)\otimes_CC'\to0
\]
et d'autre part, le produit tensoriel étant exact à droite, on a la suite
exacte
\[
T\otimes_CC'\to(R\otimes_AN)\otimes_CC'\to Q\otimes_CC'\to0,
\]
d'où finalement la suite exacte
\[
T\otimes_CC'\to(R\otimes_AN)\otimes_CC'\to
(L_0\otimes_AN)\otimes_CC'\to(M\otimes_AN)\otimes_CC'\to0.
\]
Mais par définition, pour tout $C$-module $P$,
$P\otimes_CC'=P\otimes_AA'$, d'où la conclusion.

\begin{lemma}[11.2.5]
\label{IV.11.2.5-fr}
Soient $A$ un anneau, $\mathfrak J$ un idéal de $A$, $B$ une $A$-algèbre,
$M$ un $B$-module, $A\to A'$ un homomorphisme d'anneaux. On pose
$\mathfrak J'=\mathfrak JA'$, $B'=B\otimes_AA'$,
$M'=M\otimes_AA'=M\otimes_BB'$. Soit $\mathfrak p'$ un idéal premier de
$B'$ contenant $\mathfrak J'B'$. On suppose vérifiée l'une des hypothèses
suivantes~:
\begin{enumerate}[label=\emph{\alph*)}]
\item $\mathfrak J$ est nilpotent.
\item $A'$ est noethérien, $B'$ est une $A'$-algèbre de type fini, $M'$ un
$B'$-module de type fini.
\end{enumerate}
Dans ces conditions, supposons vérifiées les deux propriétés suivantes~:
\begin{enumerate}[label=\emph{(\roman*)}]
\item $M\otimes_A(A/\mathfrak J)$ est un $(A/\mathfrak J)$-module plat.
\item L'homomorphisme canonique composé
\[
\operatorname{Tor}_1^A(M,A/\mathfrak J)\xrightarrow{\psi_1}
\operatorname{Tor}_1^{A'}(M',A'/\mathfrak J')\xrightarrow{\theta}
(\operatorname{Tor}_1^{A'}(M',A'/\mathfrak J'))_{\mathfrak p'}
\]
(où $\psi_1$ est l'homomorphisme (11.2.1.1) et $\theta$ l'homomorphisme
canonique d'un $B'$-module dans son localisé en $\mathfrak p'$) est nul.
\end{enumerate}
Alors $M'_{\mathfrak p'}$ est un $A'$-module plat.
\end{lemma}

Notons que dans l'hypothèse b) $M'_{\mathfrak p'}$ est un
$B'_{\mathfrak p'}$-module de type fini, $B'_{\mathfrak p'}$ une
$A'$-algèbre noethérienne, et que $\mathfrak J'B'_{\mathfrak p'}$ est
contenu dans le radical $\mathfrak p'B'_{\mathfrak p'}$ de
$B'_{\mathfrak p'}$~; dans l'hypothèse a), $\mathfrak J'$ est nilpotent~;
on va donc pouvoir appliquer le critère de platitude
($\mathbf 0_{\mathrm{III}}$, 10.2.1) ou ($\mathbf 0_{\mathrm{III}}$, 10.2.2)
suivant que a) ou b) est vérifiée. En premier lieu, on a
\[
M'_{\mathfrak p'}\otimes_{A'}(A'/\mathfrak J')
=(M'\otimes_{A'}(A'/\mathfrak J'))_{\mathfrak p'}
=((M\otimes_A(A/\mathfrak J))\otimes_{A/\mathfrak J}
(A'/\mathfrak J'))_{\mathfrak p'}~;
\]
l'hypothèse (i) entraîne donc que
$M'_{\mathfrak p'}\otimes_{A'}(A'/\mathfrak J')$ est un
$(A'/\mathfrak J')$-module plat, compte tenu de
($\mathbf 0_{\mathrm I}$, 6.2.1 et 6.3.2). Reste donc à voir que
$\operatorname{Tor}_1^{A'}(M'_{\mathfrak p'},A'/\mathfrak J')=0$~; or ce
$B'_{\mathfrak p'}$-module est égal à
$(\operatorname{Tor}_1^{A'}(M',A'/\mathfrak J'))_{\mathfrak p'}$, en vertu
de la platitude de $B'_{\mathfrak p'}$ sur $B'$. Mais en vertu de
l'hypothèse (ii), l'homomorphisme composé
$\theta\circ\varphi_1:\operatorname{Tor}_1^A(M,A/\mathfrak J)\otimes_AA'
\to(\operatorname{Tor}_1^{A'}(M',A'/\mathfrak J'))_{\mathfrak p'}$ est
nul~; par ailleurs, (11.2.4) appliqué à $C=A/\mathfrak J$ et $N=C$ montre
(compte tenu de l'hypothèse (i)) que $\varphi_1$ est \emph{surjectif} (car
$A'/\mathfrak J'=C\otimes_AA'$)~; donc l'homomorphisme $\theta$ est
\emph{nul} et comme l'image par $\theta$ de
$\operatorname{Tor}_1^{A'}(M',A'/\mathfrak J')$ engendre le
$B'_{\mathfrak p'}$-module
$(\operatorname{Tor}_1^{A'}(M',A'/\mathfrak J'))_{\mathfrak p'}$, ce
dernier est nul. C.Q.F.D.

\begin{theorem}[11.2.6]
\label{IV.11.2.6-fr}
Les notations étant celles de (8.5.1) et (8.8.1), on suppose $S_\alpha$
quasi-compact, $X_\alpha$ et $Y_\alpha$ de présentation finie sur
$S_\alpha$~; soient $f_\alpha:X_\alpha\to Y_\alpha$ un
$S_\alpha$-morphisme, $\mathcal F_\alpha$ un
$\mathcal O_{X_\alpha}$-Module quasi-cohérent de présentation finie.

\oldpage[IV-3]{124}

\begin{enumerate}[label=\emph{(\roman*)}]
\item Soient $x$ un point de $X$, $x_\lambda$ sa projection canonique dans
$X_\lambda$. Pour que $\mathcal F$ soit $f$-plat au point $x$, il faut et il
suffit qu'il existe $\lambda\geq\alpha$ tel que $\mathcal F_\lambda$ soit
$f_\lambda$-plat au point $x_\lambda$.
\item Pour que $\mathcal F$ soit $f$-plat, il faut et il suffit qu'il existe
$\lambda\geq\alpha$ tel que $\mathcal F_\lambda$ soit $f_\lambda$-plat.
\end{enumerate}
\end{theorem}

On peut supposer que $S_\alpha=S_0$~; comme $Y_0$ est de présentation finie
sur $S_0$, $Y_0$ est quasi-compact, et $f_0:X_0\to Y_0$ est un morphisme de
présentation finie (1.6.2, (v)), donc on peut aussi se borner au cas où
$S_0=Y_0$. En outre, en vertu de la quasi-compacité de $S_0$ et du fait que
l'ensemble d'indices $L$ est filtrant, on peut se borner au cas où
$S_0=\operatorname{Spec}(A_0)$ est affine. De plus $X_0$ est quasi-compact,
donc le même raisonnement montre qu'on peut aussi supposer
$X_0=\operatorname{Spec}(B_0)$ affine~; on a alors
$\mathcal F_0=\widetilde M_0$, où $M_0$ est un $B_0$-module de présentation
finie, et l'énoncé (11.2.6) est dans ce cas équivalent au suivant (compte tenu
de ($\mathbf 0_{\mathrm I}$, 6.3.1))~:

\begin{corollary}[11.2.6.1]
\label{IV.11.2.6.1-fr}
Soient $A_0$ un anneau, $(A_\lambda)_{\lambda\in L}$ un système inductif
filtrant de $A_0$-algèbres, $B_0$ une $A_0$-algèbre de présentation finie,
$M_0$ un $B_0$-module de présentation finie. On pose
\[
\begin{gathered}
B_\lambda=B_0\otimes_{A_0}A_\lambda,\qquad
M_\lambda=M_0\otimes_{A_0}A_\lambda=M_0\otimes_{B_0}B_\lambda,\\
A=\varinjlim A_\lambda,\qquad
B=\varinjlim B_\lambda=B_0\otimes_{A_0}A,\\
M=\varinjlim M_\lambda=M_0\otimes_{A_0}A=M_0\otimes_{B_0}B.
\end{gathered}
\]
\begin{enumerate}[label=\emph{(\roman*)}]
\item Soit $\mathfrak p$ un idéal premier de $B$, et pour tout $\lambda$,
soit $\mathfrak p_\lambda$ son image réciproque dans $B_\lambda$. Pour que
$M_{\mathfrak p}$ soit un $A$-module plat, il faut et il suffit qu'il existe
$\lambda$ tel que $(M_\lambda)_{\mathfrak p_\lambda}$ soit un
$A_\lambda$-module plat.
\item Pour que $M$ soit un $A$-module plat, il faut et il suffit qu'il existe
$\lambda$ tel que $M_\lambda$ soit un $A_\lambda$-module plat.
\end{enumerate}
\end{corollary}

On doit seulement démontrer que les conditions sont nécessaires (2.1.4).
Nous procéderons en plusieurs étapes.

\medskip
\noindent\textbf{I)} \emph{Réduction au cas où les $A_\lambda$ sont
noethériens.} En vertu de (8.9.1), il existe un sous-anneau $A'_0$ de $A_0$,
qui est une $\mathbf Z$-algèbre de type fini, une $A'_0$-algèbre de type fini
$B'_0$ et un $B'_0$-module de type fini $M'_0$ tels que l'on ait
$B_0=B'_0\otimes_{A'_0}A_0$ et $M_0=M'_0\otimes_{A'_0}A_0$~; comme on a
$B_\lambda=B'_0\otimes_{A'_0}A_\lambda$ et
$M_\lambda=M'_0\otimes_{A'_0}A_\lambda$, on peut remplacer $A_0$, $B_0$,
$M_0$ par $A'_0$, $B'_0$, $M'_0$ dans l'énoncé de (11.2.6.1), en considérant
les $A_\lambda$ comme des $A'_0$-algèbres~; on peut donc déjà supposer que
$A_0$ est noethérien. Soit $H$ l'ensemble des couples
$(\lambda,C_\lambda)$, où $C_\lambda$ est une sous-$A_0$-algèbre de
\emph{type fini} de $A_\lambda$~; ordonnons $H$ en posant
$(\lambda,C_\lambda)\leq(\mu,D_\mu)$ si $\lambda\leq\mu$ et si
l'homomorphisme $\varphi_{\mu\lambda}:A_\lambda\to A_\mu$ est tel que
$\varphi_{\mu\lambda}(C_\lambda)\subset D_\mu$~; pour cet ordre $H$ est
filtrant croissant, car si $(\lambda,C_\lambda)$ et $(\mu,D_\mu)$ sont deux
éléments quelconques de $H$, on les majorera par un $(\nu,E_\nu)$ en prenant
$\nu\geq\lambda$, $\nu\geq\mu$ dans $L$, puis $E_\nu$ égal à la
sous-$A_0$-algèbre de $A_\nu$ engendrée par
$\varphi_{\nu\lambda}(C_\lambda)$ et $\varphi_{\nu\mu}(D_\mu)$. Pour un
élément $\xi=(\lambda,C_\lambda)$ de $H$, on posera $A_\xi=C_\lambda$, et
pour $\xi\leq\eta=(\mu,D_\mu)$ (donc $\lambda\leq\mu$ et
$\varphi_{\mu\lambda}(C_\lambda)\subset D_\mu$),
$\varphi_{\eta\xi}:A_\xi\to A_\eta$ sera la restriction à $C_\lambda$ de
$\varphi_{\mu\lambda}$, considéré comme homomorphisme dans $D_\mu$~; il est
clair que l'on obtient ainsi un système inductif filtrant de $A_0$-algèbres.
On pose $B_\xi=B_0\otimes_{A_0}A_\xi$,
$M_\xi=M_0\otimes_{A_0}A_\xi$~; cette fois les $A_\xi$ sont
noethériens~; en outre la formule de la double limite inductive (Bourbaki,
\emph{Alg.}, chap.~II, 3$^e$ éd., \S~6, n$^\circ$~4, prop.~7) prouve que
l'on a encore
$\varinjlim_HA_\xi=A$, $\varinjlim_HB_\xi=B$,
$\varinjlim_HM_\xi=M$. Supposons (11.2.6.1) prouvé pour le système inductif
$(A_\xi)_{\xi\in H}$~; soit $\mathfrak p$ un idéal premier de $B$,

\oldpage[IV-3]{125}

tel que $M_{\mathfrak p}$ soit un $A$-module plat~; il existe alors
$\xi\in H$ tel que, si $\mathfrak p_\xi$ est l'image réciproque de
$\mathfrak p$ dans $B_\xi$, $(M_\xi)_{\mathfrak p_\xi}$ soit un
$A_\xi$-module plat. Soit $\xi=(\lambda,C_\lambda)$, de sorte que
l'injection $A_\xi=C_\lambda\to A_\lambda$ donne un homomorphisme
$B_\xi=B_0\otimes_{A_0}C_\lambda\to B_\lambda=B_0\otimes_{A_0}A_\lambda$,
et si $\mathfrak p_\lambda$ est l'image réciproque de $\mathfrak p$ dans
$B_\lambda$, $\mathfrak p_\xi$ est l'image réciproque de
$\mathfrak p_\lambda$ dans $B_\xi$~; on a par suite
\[
(M_\lambda)_{\mathfrak p_\lambda}
=(M_\xi\otimes_{C_\lambda}A_\lambda)_{\mathfrak p_\lambda}
=((M_\xi)_{\mathfrak p_\xi}\otimes_{C_\lambda}A_\lambda)_{\mathfrak p_\lambda},
\]
donc $(M_\lambda)_{\mathfrak p_\lambda}$ est un $A_\lambda$-module plat
($\mathbf 0_{\mathrm I}$, 6.2.1 et 6.3.2). On traite de même le cas (ii) de
l'énoncé. Nous pouvons donc par la suite supposer les $A_\lambda$
noethériens pour $\lambda\in L$ (mais non nécessairement $A$ lui-même).

\medskip
\noindent\textbf{II)} \emph{Réduction de l'énoncé global (ii) à l'énoncé
ponctuel (i).} Supposons que $\mathcal F$ soit $f$-plat. Pour tout $\lambda$,
soit $U_\lambda$ l'ensemble des $x_\lambda\in X_\lambda$ tels que
$\mathcal F_\lambda$ soit $f_\lambda$-plat au point $x_\lambda$~; on sait
que $U_\lambda$ est \emph{ouvert} dans $X_\lambda$ puisque $S_\lambda$ est
noethérien et $f_\lambda$ de type fini (11.1.1)~; soit
$V_\lambda=v_\lambda^{-1}(U_\lambda)$ son image réciproque dans $X$. Comme,
par hypothèse, pour tout $x\in X$, il y a un $\lambda$ tel que
$\mathcal F_\lambda$ soit $f_\lambda$-plat au point $x_\lambda$, projection
de $x$ dans $X_\lambda$, cela signifie que $x\in V_\lambda$ pour un
$\lambda$~; autrement dit, $X$ est réunion des $V_\lambda$. D'ailleurs
(2.1.4), pour $\lambda\leq\mu$, on a $V_\lambda\subset V_\mu$, donc, comme
$X$ est quasi-compact, il existe un indice $\mu$ tel que $X=V_\mu$. Comme
les $X_\lambda$ sont quasi-compacts, il résulte de (8.3.4) qu'il existe un
indice $\nu\geq\mu$ tel que $v_{\mu\nu}^{-1}(U_\mu)=X_\nu$~; mais par
(2.1.4), cela entraîne que $\mathcal F_\nu$ est $f_\nu$-plat.

\medskip
\noindent\textbf{III)} \emph{Fin de la démonstration.} Il reste à prouver
(i) en supposant $S_0$ affine et noethérien~; si $y_0$ est la projection de
$x$ dans $S_0$, on peut en outre, en vertu de (2.1.4) et de
($\mathbf I$, 3.6.5), remplacer $S_0$ par
$\operatorname{Spec}(\mathcal O_{y_0})$, autrement dit on peut se borner au
cas où $A_0$ est un anneau \emph{local} noethérien, d'idéal maximal
$\mathfrak m_0=\mathfrak j_{y_0}$~; par définition, $\mathfrak m_0$ est
l'image réciproque de l'idéal premier $\mathfrak p=\mathfrak j_x$ de $B$,
et $M_{\mathfrak p}$ est supposé être un $A$-module plat~; on a donc en
particulier $\operatorname{Tor}_1^A(A/\mathfrak m_0A,M_{\mathfrak p})=0$,
ce qui s'écrit aussi, puisque les $\operatorname{Tor}_1^A(A/\mathfrak m_0A,M)$
sont des $B$-modules et que $B_{\mathfrak p}$ est plat sur $B$,
$(\operatorname{Tor}_1^A(A/\mathfrak m_0A,M))_{\mathfrak p}=0$.
Notons maintenant que $\operatorname{Tor}_1^{A_0}(A_0/\mathfrak m_0,M_0)$
est un $B_0$-module \emph{de type fini}, car on peut le définir en prenant
une résolution de $A_0/\mathfrak m_0$ par des $A_0$-modules libres
\emph{de type fini} ($A_0$ étant noethérien) et en tensorisant par $M_0$, ce
qui donne des $B_0$-modules de type fini~; comme $B_0$ est noethérien,
l'homologie du complexe ainsi obtenu est bien formée de $B_0$-modules de
type fini. Soit $(t_i^0)_{1\leq i\leq m}$ un système de générateurs du
$B_0$-module $\operatorname{Tor}_1^{A_0}(A_0/\mathfrak m_0,M_0)$ et soit
$t_i$ l'image canonique (11.2.1.1) de $t_i^0$ dans
$\operatorname{Tor}_1^A(A/\mathfrak m_0A,M)$. L'hypothèse entraîne qu'il
existe un $h\in B-\mathfrak p$ tel que $ht_i=0$ pour $1\leq i\leq m$. Or,
on a (11.2.2.1)
\[
\operatorname{Tor}_1^A(A/\mathfrak m_0A,M)=
\varinjlim\operatorname{Tor}_1^{A_\lambda}
(A_\lambda/\mathfrak m_0A_\lambda,M_\lambda)~;
\]
il existe donc un $\lambda$ tel que, si les $t_i^\lambda$ sont les images
des $t_i^0$ dans
$\operatorname{Tor}_1^{A_\lambda}(A_\lambda/\mathfrak m_0A_\lambda,
M_\lambda)$, il existe $h_\lambda\in B_\lambda$ d'image $h$, tel que
$h_\lambda t_i^\lambda=0$ pour $1\leq i\leq m$. Soit
$\mathfrak p_\lambda=\mathfrak j_{x_\lambda}$ l'idéal premier de
$B_\lambda$ image réciproque de $\mathfrak p$~; on a
$h_\lambda\in B_\lambda-\mathfrak p_\lambda$, donc les images canoniques
des $t_i^0$ dans
$(\operatorname{Tor}_1^{A_\lambda}(A_\lambda/\mathfrak m_0A_\lambda,
M_\lambda))_{\mathfrak p_\lambda}$ sont nulles et par suite
l'homomorphisme
$\operatorname{Tor}_1^{A_0}(A_0/\mathfrak m_0,M_0)\to
(\operatorname{Tor}_1^{A_\lambda}(A_\lambda/\mathfrak m_0A_\lambda,
M_\lambda))_{\mathfrak p_\lambda}$ est nul. Les conditions du lemme
(11.2.5) sont donc remplies ($A_0/\mathfrak m_0$ étant un corps), et
$(M_\lambda)_{\mathfrak p_\lambda}$ est un $A_\lambda$-module plat, ce qui
achève de démontrer le théorème.

\begin{corollary}[11.2.7]
\label{IV.11.2.7-fr}
Soient $S=\operatorname{Spec}(A)$ un schéma affine, $f:X\to S$ un morphisme,
$\mathcal F$ un $\mathcal O_X$-Module quasi-cohérent, $x$ un point de $X$.
Les conditions suivantes sont équivalentes~:
\begin{enumerate}[label=\emph{\alph*)}]
\item $f$ est un morphisme de présentation finie, $\mathcal F$ est un
$\mathcal O_X$-Module de présentation finie et $\mathcal F$ est $f$-plat au
point $x$ (resp. $f$-plat).

\oldpage[IV-3]{126}

\item Il existe un schéma affine noethérien $S_0=\operatorname{Spec}(A_0)$,
un morphisme de type fini $f_0:X_0\to S_0$, un
$\mathcal O_{X_0}$-Module cohérent $\mathcal F_0$, un morphisme $S\to S_0$
tels que le $S$-préschéma $X_0\otimes_{S_0}S$ soit $S$-isomorphe à $X$ et
que, si on identifie $X$ à $X_0\otimes_{S_0}S$, $\mathcal F$ soit isomorphe
à $\mathcal F_0\otimes_{\mathcal O_{S_0}}\mathcal O_S$, et $\mathcal F_0$
soit $f_0$-plat au point $x_0$ projection de $x$ dans $X_0$ (resp.
$f_0$-plat).
\item Les conditions de b) sont vérifiées et en outre $A_0$ est une
sous-$\mathbf Z$-algèbre de type fini de $A$, le morphisme $S\to S_0$
correspondant à l'injection canonique $A_0\to A$.
\end{enumerate}
\end{corollary}

Il est clair que c) implique b) et b) implique a) en vertu de (2.1.4).
D'autre part, on peut considérer $A$ comme limite inductive de ses
sous-$\mathbf Z$-algèbres de type fini, et on sait par (8.9.1) qu'il y a une
telle sous-algèbre $A_0$ et un morphisme de type fini $f_0:X_0\to S_0$ tels
que $X$ soit $S$-isomorphe à $X_0\times_{S_0}S$ et $\mathcal F$ isomorphe à
$\mathcal F_0\otimes_{\mathcal O_{S_0}}\mathcal O_S$~; on peut donc écrire
$X=\varprojlim X_\lambda$, où $X_\lambda=X_0\otimes_{A_0}A_\lambda$,
$A_\lambda$ parcourant l'ensemble des sous-$\mathbf Z$-algèbres de type
fini de $A$ contenant $A_0$, et $\mathcal F=v_\lambda^*(\mathcal F_\lambda)$,
où $v_\lambda:X\to X_\lambda$ est la projection canonique et
$\mathcal F_\lambda=\mathcal F_0\otimes_{A_0}A_\lambda$. Il résulte de
(11.2.6) qu'il existe $\lambda$ tel que $\mathcal F_\lambda$ soit
$f_\lambda$-plat au point $x_\lambda=v_\lambda(x)$ (resp.
$f_\lambda$-plat)~; alors le sous-anneau $A_\lambda$ de $A$ vérifie les
conditions de c).

\begin{proposition}[11.2.8]
\label{IV.11.2.8-fr}
Soient $g:X\to S$, $h:Y\to S$ deux morphismes de présentation finie,
$f:X\to Y$ un $S$-morphisme, $\mathcal F$ un
$\mathcal O_X$-Module quasi-cohérent de présentation finie~; pour tout
$s\in S$, soient $X_s=g^{-1}(s)=X\times_S\operatorname{Spec}(k(s))$,
$Y_s=h^{-1}(s)=Y\times_S\operatorname{Spec}(k(s))$, $f_s$ le morphisme
$f\times1:X_s\to Y_s$, $\mathcal F_s=\mathcal F\otimes_{\mathcal O_S}k(s)$.
Alors l'ensemble $E$ des $s\in S$ tels que $\mathcal F_s$ soit $f_s$-plat
est localement constructible.
\end{proposition}

La propriété dont nous voulons démontrer qu'elle est constructible vérifie la
condition (9.2.1, (i)), en vertu de (2.2.13) et (2.5.1). Compte tenu de
(9.2.3), on peut donc se limiter au cas où $S$ est affine, noethérien et
intègre de point générique $\eta$ et à prouver que $E$ ou $S-E$ est un
voisinage de $\eta$ dans $S$. Si $\eta\in S-E$, cela résulte aussitôt du
lemme (9.4.7.1). Si au contraire $\eta\in E$, il résulte tout d'abord de
(11.2.6), appliqué suivant la méthode de (8.1.2, a)), qu'il y a un voisinage
ouvert $U$ de $\eta$ dans $S$ tel que
$\mathcal F|g^{-1}(U)$ soit $h^{-1}(U)$-plat~; \emph{a fortiori} (2.1.4)
$\mathcal F_s$ est $f_s$-plat pour tout $s\in U$, ce qui achève la
démonstration.

Le théorème suivant généralise (11.2.6.1, (ii))~:

\begin{theorem}[11.2.9]
\label{IV.11.2.9-fr}
(Raynaud). --- Soient $A_0$ un anneau, $(A_\lambda)_{\lambda\in L}$ un système
inductif filtrant de $A_0$-algèbres, $B_0$ une $A_0$-algèbre de présentation
finie, $\mathfrak J_0$ un idéal de type fini de $B_0$, $M_0$ un $B_0$-module
de présentation finie. On pose
$B_\lambda=B_0\otimes_{A_0}A_\lambda$,
$\mathfrak J_\lambda=\mathfrak J_0B_\lambda$,
$M_\lambda=M_0\otimes_{A_0}A_\lambda=M_0\otimes_{B_0}B_\lambda$,
$A=\varinjlim A_\lambda$,
$B=\varinjlim B_\lambda=B_0\otimes_{A_0}A$,
$\mathfrak J=\mathfrak J_0B$,
$M=\varinjlim M_\lambda=M_0\otimes_{A_0}A=M_0\otimes_{B_0}B$. Pour que
$\operatorname{gr}_{\mathfrak J}^*(M)$ soit un $A$-module plat, il faut et il
suffit qu'il existe $\lambda$ tel que
$\operatorname{gr}_{\mathfrak J_\lambda}^*(M_\lambda)$ soit un
$A_\lambda$-module plat~; les homomorphismes canoniques
\begin{equation}
\left\{
\begin{aligned}
\operatorname{gr}_{\mathfrak J_\lambda}^*(M_\lambda)
  \otimes_{A_\lambda}A&\longrightarrow
  \operatorname{gr}_{\mathfrak J}^*(M),\\
\operatorname{gr}_{\mathfrak J_\lambda}^*(M_\lambda)
  \otimes_{A_\lambda}A_\mu&\longrightarrow
  \operatorname{gr}_{\mathfrak J_\mu}^*(M_\mu)
\end{aligned}
\right.
\qquad\text{pour }\lambda\leq\mu
\label{IV.11.2.9.1-fr}
\tag{11.2.9.1}
\end{equation}
sont alors bijectifs et
$\operatorname{gr}_{\mathfrak J_\mu}^*(M_\mu)$ est un $A_\lambda$-module plat
pour $\lambda\leq\mu$.
\end{theorem}

\begin{proof}
Montrons d'abord que les conditions sont suffisantes, ce qui revient à prouver
la bijectivité de (11.2.9.1). Cela résulte du lemme suivant~:

\begin{lemma}[11.2.9.2]
\label{IV.11.2.9.2-fr}
Soient $A$ un anneau, $B$ une $A$-algèbre, $M$ un $B$-module, $\mathfrak J$ un
idéal de $B$. Soient $A'$ une $A$-algèbre~; posons
$B'=B\otimes_AA'$, $M'=M\otimes_AA'=M\otimes_BB'$,
$\mathfrak J'=\mathfrak JB'$. Si
$\operatorname{gr}_{\mathfrak J}^*(M)$ est un $A$-module plat,
l'homomorphisme canonique
\[
(\operatorname{gr}_{\mathfrak J}^*(M))\otimes_AA'
\longrightarrow \operatorname{gr}_{\mathfrak J'}^*(M')
\]
est bijectif.
\end{lemma}

\oldpage[IV-3]{127}

En effet, par récurrence sur $k$, l'hypothèse que les
$\mathfrak J^kM/\mathfrak J^{k+1}M$ sont des $A$-modules plats pour $k\geq0$
entraîne d'abord, par ($\mathbf 0_{\mathrm I}$, 6.1.2), que
$M/\mathfrak J^{k+1}M$ est un $A$-module plat~; en outre
($\mathbf 0_{\mathrm I}$, 6.1.2), la suite
\[
0\longrightarrow(\mathfrak J^{k+1}M)\otimes_AA'
\longrightarrow M\otimes_AA'
\longrightarrow(M/\mathfrak J^{k+1}M)\otimes_AA'\longrightarrow0
\]
est alors exacte, autrement dit $(\mathfrak J^{k+1}M)\otimes_AA'$ s'identifie
à son image canonique $\mathfrak J'^{,k+1}M'$ dans $M'=M\otimes_AA'$.
D'autre part, toujours en vertu de ($\mathbf 0_{\mathrm I}$, 6.1.2), la suite
\[
0\longrightarrow(\mathfrak J^{k+1}M)\otimes_AA'
\longrightarrow(\mathfrak J^kM)\otimes_AA'
\longrightarrow(\mathfrak J^kM/\mathfrak J^{k+1}M)\otimes_AA'
\longrightarrow0
\]
est exacte, ce qui prouve le lemme.

Pour démontrer la nécessité des conditions de (11.2.9), nous procéderons en
plusieurs étapes.

\medskip
\noindent\textbf{I)} \emph{Réduction au cas où les $A_\lambda$ sont
noethériens.} --- On procède comme dans la réduction I) de (11.2.6.1), dont
nous gardons les notations~; il faut simplement commencer par remplacer
$A'_0$ par une sous-$A'_0$-algèbre de type fini $A''_0$ de $A_0$ telle que,
si l'on pose $B''_0=B'_0\otimes_{A'_0}A''_0$, il y ait dans $B''_0$ un idéal
de type fini $\mathfrak J''_0$ tel que $\mathfrak J''_0B_0=\mathfrak J_0$.
On considère pour cela les sous-$A'_0$-algèbres $A'_\alpha$ de type fini de
$A_0$, qui forment une famille filtrante, et l'on a
$B_0=\varinjlim B'_\alpha$, où
$B'_\alpha=B'_0\otimes_{A'_0}A'_\alpha$~; il y a donc un indice $\beta$ tel
qu'un système fini de générateurs de $\mathfrak J_0$ soit formé d'images dans
$B_0$ d'éléments de $B'_\beta$~; on prendra alors $A''_0=A'_\beta$,
$B''_0=B'_\beta$ et pour $\mathfrak J''_0$ l'idéal engendré par ces éléments.
On peut donc supposer que $A_0$ (donc aussi $B_0$) est noethérien. On définit
ensuite comme \emph{loc. cit.} l'ensemble filtrant $H$ et les $A_\xi$,
$B_\xi$, $M_\xi$ pour $\xi\in H$~; on posera aussi
$\mathfrak J_\xi=\mathfrak J_0B_\xi$ pour tout $\xi\in H$. Supposons alors
que l'on ait démontré qu'il existe un $\xi=(\lambda,C_\lambda)$ tel que
$\operatorname{gr}_{\mathfrak J_\xi}^*(M_\xi)$ soit un $C_\lambda$-module
plat~; comme $M_\lambda=M_\xi\otimes_{C_\lambda}A_\lambda$, il résulte de
(11.2.9.2) que $\operatorname{gr}_{\mathfrak J_\lambda}^*(M_\lambda)$ est un
$A_\lambda$-module plat.

\medskip
\noindent\textbf{II)} \emph{Lemmes préliminaires.}

\begin{lemma}[11.2.9.3]
\label{IV.11.2.9.3-fr}
Soient $A$ un anneau,
$B=\bigoplus_{i\geq0}B_i$ une $A$-algèbre graduée à degrés positifs,
$M=\bigoplus_{i\in\mathbf Z}M_i$ un $B$-module gradué. On suppose que $B_0$
est un anneau local, et que chacun des $M_i$ est un $B_0$-module de type fini.
Pour que $M$ soit un $B$-module de type fini, il faut et il suffit que, si
$\mathfrak q$ est l'image réciproque dans $A$ de l'idéal maximal $\mathfrak m$
de $B_0$, $M\otimes_Ak(\mathfrak q)$ soit un
$(B\otimes_Ak(\mathfrak q))$-module de type fini.
\end{lemma}

La condition étant évidemment nécessaire, prouvons qu'elle est suffisante.
Comme $B_0$ (et par suite $B$) est une $A_{\mathfrak q}$-algèbre, on peut
remplacer $A$ par $A_{\mathfrak q}$, autrement dit supposer que $A$ est aussi
un anneau local dont $\mathfrak q$ est l'idéal maximal. Par hypothèse, il existe
un entier $N$ tel que l'homomorphisme canonique
\[
\left(\bigoplus_{-N\leq i\leq N}(M_i\otimes_{B_0}B)\right)
\otimes_Ak(\mathfrak q)\longrightarrow M\otimes_Ak(\mathfrak q)
\]

\oldpage[IV-3]{128}

soit surjectif~; cela signifie aussi que, pour tout entier $n$,
l'homomorphisme canonique de $k(\mathfrak q)$-modules
\[
\left(\bigoplus_{-N\leq i\leq N}(M_i\otimes_{B_0}B_{n-i})\right)
\otimes_Ak(\mathfrak q)\longrightarrow M_n\otimes_Ak(\mathfrak q)
\]
est surjectif. Or, $M_n$ est un $B_0$-module de type fini, $B_0$ un anneau
local et $\mathfrak qB_0\subset\mathfrak m$~; donc le lemme de Nakayama prouve
que chacun des homomorphismes canoniques
\[
\bigoplus_{-N\leq i\leq N}(M_i\otimes_{B_0}B_{n-i})\longrightarrow M_n
\]
est surjectif, d'où la conclusion.

\begin{corollary}[11.2.9.4]
\label{IV.11.2.9.4-fr}
Sous les hypothèses de (11.2.9.3), on suppose en outre que chacun des $B_i$ et
chacun des $M_i$ soit un $B_0$-module de présentation finie, et que $M$ soit
un $A$-module plat. Pour que $M$ soit un $B$-module de présentation finie, il
faut et il suffit que $M\otimes_Ak(\mathfrak q)$ soit un
$(B\otimes_Ak(\mathfrak q))$-module de présentation finie.
\end{corollary}

En vertu de (11.2.9.3), si la condition de l'énoncé est vérifiée, $M$ est un
$B$-module de type fini, donc il existe un $B$-module libre \emph{gradué} de
type fini $L$ (ayant donc une base finie formée d'éléments homogènes) et un
homomorphisme gradué surjectif de degré 0, $u:L\to M$. Soit $R=\ker(u)$, qui
est un $B$-module gradué~; il y a donc un nombre fini d'entiers $m_j$
($1\leq j\leq r$) tels que pour chaque entier $i$, $R_i$ soit le noyau d'un
homomorphisme surjectif
\[
\bigoplus_{1\leq j\leq r}B_{i+m_j}\longrightarrow M_i~;
\]
on conclut alors de l'hypothèse sur les $M_i$ et les $B_i$ que $R_i$ est un
$B_0$-module de type fini (Bourbaki, \emph{Alg. comm.}, chap.~I, \S~2,
n$^\circ$~8, lemme 9). Pour prouver que $R$ est un $B$-module de type fini,
notons qu'en vertu de l'hypothèse de platitude et de
($\mathbf 0_{\mathrm I}$, 6.1.2) la suite
\[
0\longrightarrow R\otimes_Ak(\mathfrak q)
\longrightarrow L\otimes_Ak(\mathfrak q)
\longrightarrow M\otimes_Ak(\mathfrak q)\longrightarrow0
\]
est exacte, et l'hypothèse entraîne donc (Bourbaki, \emph{loc. cit.}) que
$R\otimes_Ak(\mathfrak q)$ est un $(B\otimes_Ak(\mathfrak q))$-module de type
fini~; il suffit donc d'appliquer à $R$ le lemme (11.2.9.3).

\medskip
\noindent\textbf{III)} \emph{Réduction au cas où les homomorphismes de
transitions $\varphi_{\mu\lambda}:A_\lambda\to A_\mu$ (pour
$\lambda\leq\mu$) sont injectifs.} --- Soit $A'_\lambda$ l'image de
$A_\lambda$ par l'homomorphisme canonique $\varphi_\lambda:A_\lambda\to A$~;
il est clair que les $A'_\lambda$ forment un système inductif de sous-anneaux
noethériens de $A$, dont la limite inductive est $A$, et les homomorphismes de
transition $A'_\lambda\to A'_\mu$ (pour $\lambda\leq\mu$) sont injectifs.
Posons $B'_\lambda=B_0\otimes_{A_0}A'_\lambda$,
$\mathfrak J'_\lambda=\mathfrak J_0B'_\lambda$,
$M'_\lambda=M_0\otimes_{A_0}A'_\lambda=M_\lambda\otimes_{A_\lambda}A'_\lambda$~;
on a encore $B=\varinjlim B'_\lambda$, $M=\varinjlim M'_\lambda$. Supposons
que l'on ait prouvé qu'il existe un $\lambda$ tel que, pour $\mu\geq\lambda$,
$\operatorname{gr}_{\mathfrak J'_\mu}^*(M'_\mu)$ soit un $A'_\mu$-module
plat. Soit $\mathfrak a_\lambda$ le noyau de l'homomorphisme $\varphi_\lambda$,
qui est donc un idéal de type fini de l'anneau noethérien $A_\lambda$. Il
résulte de la définition de la limite inductive qu'il existe un indice
$\mu\geq\lambda$ tel que $\varphi_{\mu\lambda}(\mathfrak a_\lambda)=0$,
autrement dit l'homomorphisme $\varphi_{\mu\lambda}$ se factorise en
$\varphi_{\mu\lambda}:A_\lambda\to A'_\lambda\to A_\mu$, et l'on peut écrire
$B_\mu=B'_\lambda\otimes_{A'_\lambda}A_\mu$,
$\mathfrak J_\mu=\mathfrak J'_\lambda B_\mu$, et
$M_\mu=M'_\lambda\otimes_{A'_\lambda}A_\mu$~; on déduit donc de (11.2.9.2)
que $\operatorname{gr}_{\mathfrak J_\mu}^*(M_\mu)$ est un $A_\mu$-module
plat.

\medskip
\noindent\textbf{IV)} \emph{Réduction au cas où
$B_0=A_0[T_1,\ldots,T_n]$ et
$\operatorname{gr}_{\mathfrak J_0}^*(B_0)=B_0$.} --- En vertu de l'hypothèse,
il y a un système $(t_i)_{1\leq i\leq n}$ de générateurs de la $A_0$-algèbre
de type fini $B_0$, tel que les $t_i$ pour $i\leq m$ engendrent l'idéal
$\mathfrak J_0$ de $B_0$. Posons $B'_0=A_0[T_1,\ldots,T_n]$,

\oldpage[IV-3]{129}

algèbre de polynômes (donc noethérienne), et soit $\mathfrak J'_0$ l'idéal de
$B'_0$ engendré par les $T_i$ d'indice $i\leq m$~; on a alors un
$A_0$-homomorphisme surjectif $u_0:B'_0\to B_0$ transformant chaque $T_i$ en
$t_i$ ($1\leq i\leq n$), donc tel que $u_0(\mathfrak J'_0)=\mathfrak J_0$~;
cet homomorphisme permet de considérer $M_0$ comme un $B'_0$-module de type
fini. On pose alors
\[
B'_\lambda=B'_0\otimes_{A_0}A_\lambda
=A_\lambda[T_1,\ldots,T_n],\qquad
\mathfrak J'_\lambda=\mathfrak J'_0B'_\lambda,
\]
\[
B'=B'_0\otimes_{A_0}A=A[T_1,\ldots,T_n],\qquad
\mathfrak J'=\mathfrak J'_0B',
\]
de sorte que $\mathfrak J'$ est l'idéal de $B'$ engendré par les $T_i$
d'indice $i\leq m$~; en outre, il est clair que pour tout entier $k\geq0$,
on a $\mathfrak J'^kM=\mathfrak J^kM$ et
$\mathfrak J'_\lambda{}^kM_\lambda=\mathfrak J_\lambda^kM_\lambda$ pour tout
$\lambda$~; donc
$\operatorname{gr}_{\mathfrak J'}^*(M)=\operatorname{gr}_{\mathfrak J}^*(M)$
en tant que $A$-module et
$\operatorname{gr}_{\mathfrak J'_\lambda}^*(M_\lambda)
=\operatorname{gr}_{\mathfrak J_\lambda}^*(M_\lambda)$ en tant que
$A_\lambda$-module~; on peut donc substituer $B'$, $B'_\lambda$,
$\mathfrak J'$, $\mathfrak J'_\lambda$ à $B$, $B_\lambda$, $\mathfrak J$,
$\mathfrak J_\lambda$ respectivement dans la démonstration~; on notera en
outre que par construction $\operatorname{gr}_{\mathfrak J'}^*(B')$
s'identifie à $B'$ et $\operatorname{gr}_{\mathfrak J'_\lambda}^*(B'_\lambda)$
à $B'_\lambda$.

\medskip
\noindent\textbf{V)} \emph{Preuve du fait que
$\operatorname{gr}_{\mathfrak J}^*(M)$ est un
$\operatorname{gr}_{\mathfrak J}^*(B)$-module de présentation finie.} --- Nous
supposons donc désormais $A_0$ et les $A_\lambda$ noethériens, les
homomorphismes de transition $A_\lambda\to A_\mu$ injectifs,
$B_0=A_0[T_1,\ldots,T_n]$, $\mathfrak J_0$ étant engendré par les $T_i$
d'indice $i\leq m$~; l'anneau
$C_0=\operatorname{gr}_{\mathfrak J_0}^*(B_0)$ s'identifie donc à $B_0$ et
$C=\operatorname{gr}_{\mathfrak J}^*(B)$ s'identifie à $B$~; nous utiliserons
seulement tout d'abord dans ce qui suit le fait que
$C\cong C_0\otimes_{A_0}A$.

Nous aurons besoin de la variante suivante de (6.9.3)~:

\begin{lemma}[11.2.9.5]
\label{IV.11.2.9.5-fr}
Soient $A_0$ un anneau noethérien, $B_0$ une $A_0$-algèbre de type fini,
$\mathfrak J_0$ un idéal de $B_0$, $M_0$ un $B_0$-module de type fini. Il
existe alors une suite $(S_{0i})_{1\leq i\leq m}$ de sous-schémas de
$S_0=\operatorname{Spec}(A_0)$ ayant les propriétés suivantes~:
\begin{enumerate}[label=\arabic{enumi}$^\circ$]
\item Les espaces sous-jacents aux $S_{0i}$ sont deux à deux disjoints et
forment un recouvrement de $S_0$.
\item Pour chaque $i$, l'ensemble $S_{0i}$ est ouvert dans
$\bigcup_{j\geq i}S_{0j}$.
\item Chaque schéma $S_{0i}$ est affine.
\item Si $A_{0i}$ est l'anneau de $S_{0i}$ et si l'on pose
$B_{0i}=B_0\otimes_{A_0}A_{0i}$, $\mathfrak J_{0i}=\mathfrak J_0B_{0i}$,
alors
$\operatorname{gr}_{\mathfrak J_{0i}}^*(M_0\otimes_{A_0}A_{0i})$ est un
$A_{0i}$-module plat.
\end{enumerate}
\end{lemma}

On procède comme dans (6.9.3) par récurrence noethérienne, en supposant le
lemme vrai lorsqu'on y remplace $A_0$ par $A'_0=A_0/\mathfrak a_0$, où
l'idéal $\mathfrak a_0$ est tel que $V(\mathfrak a_0)\neq S_0$, $B_0$ par
$B'_0=B_0\otimes_{A_0}A'_0$, $\mathfrak J_0$ par
$\mathfrak J'_0=\mathfrak J_0B'_0$ et $M_0$ par $M_0\otimes_{A_0}A'_0$. On
considère le nilradical $\mathfrak N_0$ de $A_0$, et il suffit évidemment de
prouver le lemme en remplaçant $A_0$ par $A_0/\mathfrak N_0$ et $B_0$,
$\mathfrak J_0$, $M_0$ par les objets correspondants comme ci-dessus.
Autrement dit, on peut supposer que $A_0$ est \emph{réduit}~; d'autre part,
$\operatorname{gr}_{\mathfrak J_0}^*(B_0)$ est une $A_0$-algèbre de type fini
et $\operatorname{gr}_{\mathfrak J_0}^*(M_0)$ un
$\operatorname{gr}_{\mathfrak J_0}^*(B_0)$-module de type fini~; en vertu du
théorème de platitude générique (6.9.1), il existe un ouvert
$D(t_0)\subset S_0$ tel que si l'on pose $A_{01}=(A_0)_{t_0}$,
$\operatorname{gr}_{\mathfrak J_0}^*(M_0)\otimes_{A_0}A_{01}$ soit un
$A_{01}$-module plat~; mais puisque $A_{01}$ est un $A_0$-module plat,
$\operatorname{gr}_{\mathfrak J_0}^*(M_0)\otimes_{A_0}A_{01}$ s'identifie à
$\operatorname{gr}_{\mathfrak J_{01}}^*(M_0\otimes_{A_0}A_{01})$, où
$B_{01}=B_0\otimes_{A_0}A_{01}$ et $\mathfrak J_{01}=\mathfrak J_0B_{01}$.
Le complémentaire de $D(t_0)$ dans $S_0$ est alors de la forme
$V(\mathfrak a_0)$, et on conclut en appliquant l'hypothèse de récurrence
noethérienne.

Appliquons ce lemme à la situation actuelle, en gardant les mêmes notations~;
posons $U_{0i}=\bigcup_{j\leq i}S_{0j}$, qui est un ouvert quasi-compact de
$S_0$. Il y a donc une famille
$(f_{ik})_{1\leq i\leq m,\,1\leq k\leq n}$ d'éléments de $A_0$ tels que pour
chaque $i$, $U_{0i}$ soit réunion des $D(f_{ik})$ ($1\leq k\leq n$).

\oldpage[IV-3]{130}

Le $C_0$-module $\operatorname{gr}_{\mathfrak J_0}^*(M_0)$ est engendré par
un nombre fini d'éléments homogènes de degré 1, de sorte qu'il y a un
épimorphisme de $C_0$-modules gradués
\[
u_0:C_0^r\longrightarrow\operatorname{gr}_{\mathfrak J_0}^*(M_0).
\]
Comme $C=C_0\otimes_{A_0}A$ et que l'homomorphisme canonique
$\operatorname{gr}_{\mathfrak J_0}^*(M_0)\otimes_{A_0}A
\to\operatorname{gr}_{\mathfrak J}^*(M)$ est surjectif, on déduit donc de
$u_0$ un épimorphisme de $C$-modules gradués
\[
u:C^r\xrightarrow{u_0\otimes1_A}
\operatorname{gr}_{\mathfrak J_0}^*(M_0)\otimes_{A_0}A
\longrightarrow\operatorname{gr}_{\mathfrak J}^*(M),
\]
et tout revient à voir que le $C$-module gradué $Q=\ker(u)$ est \emph{de type
fini}.

\begin{lemma}[11.2.9.6]
\label{IV.11.2.9.6-fr}
Sous les hypothèses précédentes, soit $\mathfrak R_0$ un idéal de $A_0$~;
posons
\[
A'_0=A_0/\mathfrak R_0,\qquad
B'_0=B_0\otimes_{A_0}A'_0=B_0/\mathfrak R_0B_0,\qquad
\mathfrak J'_0=\mathfrak J_0B'_0,
\]
et $\mathfrak R=\mathfrak R_0A$, $A'=A/\mathfrak R$, et supposons en outre
que $\operatorname{gr}_{\mathfrak J'_0}^*(M_0\otimes_{A_0}A'_0)$ soit un
$A'_0$-module plat. Alors il existe un nombre fini d'éléments $q_h$
($1\leq h\leq s$) de $Q$ tels que, pour tout idéal premier
$\mathfrak p\supset\mathfrak RB$ de $B$, les images canoniques des $q_h$ dans
$Q_\mathfrak p$ engendrent le $C_\mathfrak p$-module $Q_\mathfrak p$.
\end{lemma}

Supposons ce lemme démontré, et notons qu'on peut encore l'appliquer en
remplaçant $A_0$ par $(A_0)_{t_0}$, $B_0$, $\mathfrak J_0$, $M_0$,
$A_\lambda$ et $A$ par les objets correspondants $(B_0)_{t_0}$,
$(\mathfrak J_0)_{t_0}$, $(M_0)_{t_0}$, $(A_\lambda)_{t_0}$ et $A_{t_0}$ pour
tout $t_0\in A_0$, puisque $(A_0)_{t_0}$ est un $A_0$-module plat. On
appliquera alors ce lemme en remplaçant successivement $A_0$ par chacun des
anneaux $(A_0)_{f_{ik}}$, $\mathfrak R_0$ étant remplacé par l'idéal
$\mathfrak R_{ik}$ définissant le sous-schéma fermé de $D(f_{ik})$ induit par
$S_{0i}$ sur l'ouvert $S_{0i}\cap D(f_{ik})$ de $S_{0i}$. Comme l'hypothèse
de platitude de (11.2.9.6) est vérifiée dans chacun de ces cas en raison de
(11.2.9.5), on obtient ainsi pour chaque couple $(i,k)$ une famille finie
$(a_{ikh})_{1\leq h\leq l}$ d'éléments de $Q_{f_{ik}}$ dont les images dans
$Q_{\mathfrak p_{ik}}$ engendrent ce $C_{\mathfrak p_{ik}}$-module, pour tout
idéal premier $\mathfrak p_{ik}$ de $B_{f_{ik}}$ contenant
$\mathfrak R_{ik}B_{f_{ik}}$. On peut écrire, pour un entier $d>0$ convenable,
$a_{ikh}=b_{ikh}/f_{ik}^d$ pour tous les indices $i$, $k$, $h$, avec
$b_{ikh}\in Q$. Puisque les $S_{0i}$ recouvrent $S_0$, tout idéal premier
$\mathfrak p$ de $B$ est tel que son image dans $S_0$ appartienne à un
ensemble $S_{0i}\cap D(f_{ik})$, donc l'image $\mathfrak p_{ik}$ de
$\mathfrak p$ dans $B_{f_{ik}}$ contient $\mathfrak R_{ik}B_{f_{ik}}$. Comme
les images des $b_{ikh}$ ($1\leq h\leq l$) dans
$Q_{\mathfrak p_{ik}}=Q_\mathfrak p$ engendrent ce $C_\mathfrak p$-module,
on en conclut que les $b_{ikh}$
($1\leq i\leq m$, $1\leq k\leq n$, $1\leq h\leq l$) engendrent le
$C$-module $Q$ (Bourbaki, \emph{Alg. comm.}, chap.~II, \S~3, n$^\circ$~3,
th.~1).

Reste à démontrer le lemme (11.2.9.6). Posons
\[
C'=C\otimes_AA'=C/\mathfrak RC,\qquad
B'=B\otimes_AA',\qquad \mathfrak J'=\mathfrak JB',\qquad
M'=M\otimes_AA'.
\]
Par hypothèse $\operatorname{gr}_{\mathfrak J}^*(M)$ est un $A$-module plat,
donc $\operatorname{gr}_{\mathfrak J'}^*(M')$ s'identifie à
$\operatorname{gr}_{\mathfrak J}^*(M)\otimes_AA'$ (11.2.9.2), et on a la
suite exacte ($\mathbf 0_{\mathrm I}$, 6.1.2)
\begin{equation}
0\longrightarrow Q\otimes_AA'\longrightarrow C'^{,r}
\xrightarrow{u\otimes1_{A'}}\operatorname{gr}_{\mathfrak J'}^*(M')
\longrightarrow0.
\label{IV.11.2.9.7-fr}
\tag{11.2.9.7}
\end{equation}
Posons $C'_0=C_0\otimes_{A_0}A'_0=C_0/\mathfrak R_0C_0$ et considérons
l'épimorphisme de $C'_0$-modules gradués
\[
u'_0:C'_0{}^{,r}\xrightarrow{u_0\otimes1_{A'_0}}
\operatorname{gr}_{\mathfrak J_0}^*(M_0)\otimes_{A_0}A'_0
\longrightarrow
\operatorname{gr}_{\mathfrak J'_0}^*(M_0\otimes_{A_0}A'_0).
\]
Soit $Q_0=\ker(u'_0)$~; utilisant le fait que
$\operatorname{gr}_{\mathfrak J'_0}^*(M_0\otimes_{A_0}A'_0)$ est un
$A'_0$-module plat, on voit que son produit tensoriel sur $A'_0$ par $A'$
s'identifie à $\operatorname{gr}_{\mathfrak J'}^*(M')$ (11.2.9.2) et on a la
suite exacte ($\mathbf 0_{\mathrm I}$, 6.1.2)
\begin{equation}
0\longrightarrow Q_0\otimes_{A'_0}A'\longrightarrow C'^{,r}
\xrightarrow{u\otimes1_{A'}}\operatorname{gr}_{\mathfrak J'}^*(M')
\longrightarrow0.
\label{IV.11.2.9.8-fr}
\tag{11.2.9.8}
\end{equation}

\oldpage[IV-3]{131}

d'où, par comparaison avec (11.2.9.7), un isomorphisme de $C'$-modules
gradués
\[
Q\otimes_AA'\xrightarrow{\sim}Q_0\otimes_{A'_0}A'.
\]
Comme $C'_0$ est noethérien, $Q_0$ est un $C'_0$-module de type fini, donc
$Q_0\otimes_{A'_0}A'$ est un $C'$-module gradué de type fini, et il en est par
suite de même de $Q\otimes_AA'$~; soient $q_h$ ($1\leq h\leq s$) des éléments
homogènes de $Q$ dont les images dans $Q\otimes_AA'$ engendrent ce
$C'$-module.

Considérons maintenant un idéal premier quelconque
$\mathfrak p\supset\mathfrak RB$ de $B$~; on a
$C_\mathfrak p=\operatorname{gr}_{\mathfrak J_\mathfrak p}^*(B_\mathfrak p)$
par platitude, et
$\operatorname{gr}_{\mathfrak J_\mathfrak p}^0(B_\mathfrak p)
=B_\mathfrak p/\mathfrak J_\mathfrak p$ est réduit à 0 ou est un anneau
local~; pour prouver le lemme, on peut évidemment se borner au second cas. Il
est clair alors que chacun des $(B_\mathfrak p/\mathfrak J_\mathfrak p)$-modules
$\operatorname{gr}_{\mathfrak J_\mathfrak p}^i(B_\mathfrak p)$ est de
présentation finie. D'autre part, pour tout indice $i$,
$M/\mathfrak J^iM$ s'identifie à
$(M_0/\mathfrak J_0^iM_0)\otimes_{A_0}A$, et est donc un $B$-module de
présentation finie. Par récurrence sur $i$, l'hypothèse que
$\operatorname{gr}_{\mathfrak J}^*(M)$ est un $A$-module plat entraîne, en
vertu de ($\mathbf 0_{\mathrm I}$, 6.1.2), que $M/\mathfrak J^iM$ est un
$A$-module plat. L'application de (11.2.6.1) où l'on remplace $M_0$ par
$M_0/\mathfrak J_0^kM_0$ pour $k\leq i$ montre qu'il existe un indice
$\lambda_i$ tel que, pour $\mu\geq\lambda_i$, chacun des
$M_\mu/\mathfrak J_\mu^{k+1}M_\mu$ soit un $A_\mu$-module plat, et par suite
aussi ($\mathbf 0_{\mathrm I}$, 6.1.2) chacun des
$\operatorname{gr}_{\mathfrak J_\mu}^k(M_\mu)$ est un $A_\mu$-module plat
pour $k\leq i$. On déduit par suite de (11.2.9.2) que chacun des
$(B/\mathfrak J)$-modules $\operatorname{gr}_{\mathfrak J}^i(M)$ est de
présentation finie, donc
$\operatorname{gr}_{\mathfrak J_\mathfrak p}^i(M_\mathfrak p)$ est un
$(B_\mathfrak p/\mathfrak J_\mathfrak p)$-module de présentation finie. Par
ailleurs les images des $q_h$ dans
\[
Q_\mathfrak p\otimes_{B_\mathfrak p}k(\mathfrak p)
=(Q\otimes_AA')_\mathfrak p\otimes_{A'}k(\mathfrak p)
\]
engendrent ce module sur
$C_\mathfrak p\otimes_{B_\mathfrak p}k(\mathfrak p)$, puisque
\[
C_\mathfrak p\otimes_{B_\mathfrak p}k(\mathfrak p)
=C'_\mathfrak p\otimes_{A'}k(\mathfrak p)~;
\]
comme on a la suite exacte
\[
Q_\mathfrak p\otimes_{B_\mathfrak p}k(\mathfrak p)
\longrightarrow C_\mathfrak p^{,r}\otimes_{B_\mathfrak p}k(\mathfrak p)
\longrightarrow
\operatorname{gr}_{\mathfrak J_\mathfrak p}^*(M_\mathfrak p)
\otimes_{B_\mathfrak p}k(\mathfrak p)\longrightarrow0,
\]
on en conclut que
$\operatorname{gr}_{\mathfrak J_\mathfrak p}^*(M_\mathfrak p)
\otimes_{B_\mathfrak p}k(\mathfrak p)$ est un
$(C_\mathfrak p\otimes_{B_\mathfrak p}k(\mathfrak p))$-module de présentation
finie. Appliquant le lemme (11.2.9.4) où $A$, $B$, $M$ sont remplacés
respectivement par $B_\mathfrak p$, $C_\mathfrak p$ et
$\operatorname{gr}_{\mathfrak J_\mathfrak p}^*(M_\mathfrak p)$, on en conclut
que $Q_\mathfrak p$ est un $C_\mathfrak p$-module de type fini. Utilisant
maintenant le lemme de Nakayama (et le fait que $C=B$) on en déduit que les
images des $q_h$ dans $Q_\mathfrak p$ engendrent ce $C_\mathfrak p$-module.
Ceci achève de prouver le lemme (11.2.9.6) et le fait que
$\operatorname{gr}_{\mathfrak J}^*(M)$ est un $C$-module de présentation
finie.

\medskip
\noindent\textbf{VI)} \emph{Fin de la démonstration.} --- Posons pour abréger
$C_\lambda=C_0\otimes_{A_0}A_\lambda$ (égal en fait à $B_\lambda$),
$N_\lambda=\operatorname{gr}_{\mathfrak J_\lambda}^*(M_\lambda)$ et
$N=\operatorname{gr}_{\mathfrak J}^*(M)$. Notons d'abord que pour chaque
entier $k$, $\mathfrak J^kM$ s'identifie à
$\varinjlim\mathfrak J_\lambda^kM_\lambda$ par exactitude du foncteur
$\varinjlim$~; l'image de $\mathfrak J_\lambda^kM_\lambda$ dans $M_\mu$ pour
$\lambda\leq\mu$ (resp. dans $M$) engendrant $\mathfrak J_\mu^kM_\mu$
(resp. $\mathfrak J^kM$)~; utilisant encore l'exactitude de $\varinjlim$, on
en conclut que $N$ s'identifie canoniquement à $\varinjlim N_\lambda$. Faisant
cette identification, nous allons d'abord prouver que~:
\begin{equation}
\text{\emph{Pour $\lambda$ assez grand, l'homomorphisme canonique
$v_\lambda:N_\lambda\otimes_{A_\lambda}A\to N$ est bijectif.}}
\label{IV.11.2.9.9-fr}
\tag{11.2.9.9}
\end{equation}
Comme les $C_\lambda$ sont noethériens et $N_\lambda$ un $C_\lambda$-module
de type fini, les $C$-modules $N_\lambda\otimes_{A_\lambda}A$ sont de
présentation finie et forment un système inductif filtrant, dont la limite
inductive s'identifie canoniquement à $N$ en vertu du fait que $\varinjlim$
commute aux produits tensoriels. En outre, les homomorphismes de transition
$v_{\mu\lambda}:N_\lambda\otimes_{A_\lambda}A
\to N_\mu\otimes_{A_\mu}A$

\oldpage[IV-3]{132}

pour $\lambda\leq\mu$ et les homomorphismes $v_\lambda$ sont \emph{surjectifs}.
Pour un $\lambda$ fixé, considérons les sous-$C$-modules
$\ker(v_{\mu\lambda})=N'_\mu$ de $N_\lambda\otimes_{A_\lambda}A$~; par
définition de la limite inductive, ils forment une famille filtrante croissante
de sous-$C$-modules de $\ker(v_\lambda)$, dont la réunion est
$\ker(v_\lambda)$~; mais on a vu dans V) que $N$ est un $C$-module de
présentation finie, donc (Bourbaki, \emph{Alg. comm.}, chap.~I, \S~2,
n$^\circ$~8, lemme 9) $\ker(v_\lambda)$ est un $C$-module de type fini~; il
existe par suite un indice $\mu\geq\lambda$ tel que
$N'_\mu=\ker(v_\lambda)$, ce qui signifie (vu le fait que $v_{\mu\lambda}$ est
surjectif) que $v_\mu$ est \emph{injectif}~; comme il est surjectif, cela
prouve (11.2.9.9).

Quitte à remplacer $A_0$ et $M_0$ par $A_\lambda$ et $M_\lambda$ pour
$\lambda$ assez grand, on peut donc supposer que l'homomorphisme canonique
$v_\lambda$ est bijectif pour \emph{tout} $\lambda$. Posons alors
$P_\lambda=N_0\otimes_{C_0}C_\lambda=N_0\otimes_{A_0}A_\lambda$, de sorte que
$N=\varinjlim P_\lambda$. Comme $C_0$ est une $A_0$-algèbre de présentation
finie et $P_0$ un $C_0$-module de présentation finie, on peut appliquer à
$C$ et $N$ le résultat de (11.2.6.1) et on voit donc qu'il existe un indice
$\lambda$ tel que, pour $\mu\geq\lambda$, $P_\mu$ soit un $A_\mu$-module
plat. Or, pour $\mu\geq\lambda$, on a un diagramme commutatif
\[
\xymatrix{
P_\mu \ar[r] \ar[d]^{w_\mu} & N \ar@{=}[d] \\
N_\mu \ar[r] & N
}
\]
où il résulte des définitions que $w_\mu$ est \emph{surjectif}. Comme, en
vertu de III), les homomorphismes $A_\mu\to A$ sont injectifs et que $P_\mu$
est un $A_\mu$-module plat, l'homomorphisme canonique
$P_\mu\to N=P_\mu\otimes_{A_\mu}A$ est \emph{injectif}~; on conclut donc du
diagramme commutatif précédent que $w_\mu$ est aussi injectif, donc bijectif,
et par suite $N_\mu$ est un $A_\mu$-module plat pour $\mu\geq\lambda$, ce qui
achève de prouver (11.2.9).
\end{proof}

\begin{remark}[11.2.10]
\label{IV.11.2.10-fr}
Nous ignorons si la généralisation de (11.2.6, (i)) analogue au th. de Raynaud,
est valable.
\end{remark}

\subsection{Application à l'élimination d'hypothèses noethériennes}
\label{subsection:IV.11.3-fr}
\label{IV.11.3-fr}

\begin{theorem}[11.3.1]
\label{IV.11.3.1-fr}
Soient $f:X\to Y$ un morphisme localement de présentation finie,
$\mathcal F$ un $\mathcal O_X$-Module quasi-cohérent de présentation finie.
Alors l'ensemble $U$ des points $x\in X$ tels que $\mathcal F$ soit $f$-plat
au point $x$ est ouvert dans $X$. Si de plus $\operatorname{Supp}(\mathcal F)=X$,
$f|U$ est un morphisme ouvert.
\end{theorem}

\begin{proof}
La seconde assertion a déjà été démontrée (2.4.6) et n'a été insérée que pour
mémoire.

La question étant locale sur $X$ et $Y$, on peut supposer que $X$ et $Y$ sont
affines, et que $f$ est un morphisme de présentation finie. Soit $x\in X$ un
point tel que $\mathcal F$ soit $f$-plat au point $x$. Appliquant (11.2.7), on
peut supposer que
\[
X=X_0\times_{Y_0}Y,\qquad f=f_0\times1,\qquad
\mathcal F=\mathcal F_0\otimes_{\mathcal O_{Y_0}}\mathcal O_Y,
\]
où $Y_0$ est noethérien, $f_0:X_0\to Y_0$ un morphisme de type fini,
$\mathcal F_0$ un $\mathcal O_{X_0}$-Module cohérent~; en outre, si $x_0$ est la
projection canonique de $x$ dans $X_0$, on peut supposer que $\mathcal F_0$ est
$f_0$-plat au point $x_0$. Alors, en vertu de (11.1.1), l'ensemble $U_0$ des
points de $X_0$ où $\mathcal F_0$ est $f_0$-plat est un voisinage de $x_0$~;
donc $\mathcal F$ est $f$-plat aux points de l'image réciproque de $U_0$ dans
$X$ (2.1.4), et cela prouve que $U$ est un voisinage de $x$.
\end{proof}

\oldpage[IV-3]{133}

\begin{corollary}[11.3.2]
\label{IV.11.3.2-fr}
Soient $f:X\to Y$, $g:Y\to Z$ deux morphismes de présentation finie,
$\mathcal F$ un $\mathcal O_X$-Module quasi-cohérent de présentation finie.
Alors l'ensemble $U$ des $y\in Y$ tels que, pour toute générisation $y'$ de
$y$, $\mathcal F$ soit $(g\circ f)$-plat en tous les points de $f^{-1}(y')$
(i.e. tels que $\mathcal F\otimes_Y\operatorname{Spec}(\mathcal O_{Y,y})$
soit plat relativement au morphisme
$X\times_Y\operatorname{Spec}(\mathcal O_{Y,y})\to Z$) est ouvert dans $Y$.
\end{corollary}

Le même raisonnement que dans (11.1.5) montre que si $V$ est l'ensemble des
$x\in X$ où $\mathcal F$ est $(g\circ f)$-plat, $U$ est l'ensemble des $y\in Y$
dont toute générisation appartient à $E=Y-f(X-V)$. Comme $g\circ f$ est de
présentation finie, $V$ est ouvert dans $X$ en vertu de (11.3.1), donc
ind-constructible (1.9.6), et par suite $X-V$ est pro-constructible~; mais comme
$f$ est quasi-compact, $f(X-V)$ est pro-constructible dans $Y$ (1.9.5, (vii)),
et par suite $E$ est ind-constructible dans $Y$. Mais il résulte alors de
(1.10.1) que $U$ est l'intérieur de $E$, donc ouvert dans $Y$.

\begin{corollary}[11.3.3]
\label{IV.11.3.3-fr}
Soient $A$ un anneau, $B$ une $A$-algèbre de présentation finie, $C$ une
$B$-algèbre de présentation finie, $M$ un $C$-module de présentation finie.
Alors l'ensemble des $\mathfrak q\in\operatorname{Spec}(B)$ tels que
$M_\mathfrak q$ soit un $A$-module plat est ouvert dans
$\operatorname{Spec}(B)$.
\end{corollary}

\begin{proposition}[11.3.4]
\label{IV.11.3.4-fr}
Soient $f:X\to S$ un morphisme localement de présentation finie, $\mathcal J$
un Idéal quasi-cohérent de type fini de $\mathcal O_X$, $Y$ le sous-préschéma
fermé de $X$ défini par $\mathcal J$, $\mathcal F$ un $\mathcal O_X$-Module de
présentation finie. On suppose que $\operatorname{gr}_{\mathcal J}^*(\mathcal F)$
est $f$-plat. Dans ces conditions~:
\begin{enumerate}[label=\emph{(\roman*)}]
\item $\mathcal F$ est $f$-plat aux points de $Y$.
\item Si $\operatorname{gr}_{\mathcal J}^*(\mathcal O_X)$ est $f$-plat,
$\operatorname{gr}_{\mathcal J}^*(\mathcal O_X)$ est une
$(\mathcal O_X/\mathcal J)$-Algèbre de présentation finie et
$\operatorname{gr}_{\mathcal J}^*(\mathcal F)$ est un
$(\operatorname{gr}_{\mathcal J}^*(\mathcal O_X))$-Module de présentation
finie.
\item Supposons $\operatorname{gr}_{\mathcal J}^*(\mathcal O_X)$ $f$-plat
(ce qui entraîne que $\mathcal O_X/\mathcal J$ est $f$-plat). Alors l'ensemble
$W$ des $y\in Y$ tels que $(\operatorname{gr}_{\mathcal J}^*(\mathcal F))_y$
soit un $(\mathcal O_X/\mathcal J)_y$-module plat est ouvert dans $Y$.
\item Supposons $\operatorname{gr}_{\mathcal J}^*(\mathcal O_X)$ $f$-plat.
Soit $S'\to S$ un morphisme quelconque~; soient $X'=X\times_SS'$,
$p:X'\to X$ la projection canonique, $Y'=p^{-1}(Y)=Y\times_SS'$ le
sous-préschéma fermé de $X'$, défini par
$\mathcal J'=p^*(\mathcal J)\mathcal O_{X'}$, $\mathcal F'=p^*(\mathcal F)
=\mathcal F\otimes_{\mathcal O_S}\mathcal O_{S'}$~; alors, si
$W'=p^{-1}(W)$, on a
$\operatorname{gr}_{\mathcal J'}^*(\mathcal F')|W'
\simeq p^*(\operatorname{gr}_{\mathcal J}^*(\mathcal F)|W)$, et
$\operatorname{gr}_{\mathcal J'}^*(\mathcal F')$ est un
$(\mathcal O_{X'}/\mathcal J')$-Module plat aux points de $W'$.
\end{enumerate}
\end{proposition}

\begin{proof}
Les questions étant locales sur $X$ et $S$, on peut supposer que
$S=\operatorname{Spec}(A)$ et $X=\operatorname{Spec}(B)$ sont affines, $B$ étant
une $A$-algèbre de présentation finie, et $\mathcal F=\widetilde M$, où $M$ est
un $B$-module de présentation finie, $\mathcal J=\widetilde{\mathfrak J}$, où
$\mathfrak J$ est un idéal de type fini de $B$~; en vertu de (8.9.1) et
(8.5.11), il existe un sous-anneau noethérien $A_0$ de $A$, une $A_0$-algèbre
de type fini $B_0$ tels que $B=B_0\otimes_{A_0}A$, un idéal $\mathfrak J_0$ de
$B_0$ tel que $\mathfrak J=\mathfrak J_0B$, un $B_0$-module de type fini $M_0$
tel que $M=M_0\otimes_{A_0}A=M_0\otimes_{B_0}B$. En outre, $A$ est limite
inductive de ses sous-$A_0$-algèbres de type fini $A_\lambda$~; on pose
$B_\lambda=B_0\otimes_{A_0}A_\lambda$,
$M_\lambda=M_0\otimes_{A_0}A_\lambda=M_0\otimes_{B_0}B_\lambda$,
$\mathfrak J_\lambda=\mathfrak J_0B_\lambda$, de sorte que
$B=\varinjlim B_\lambda$, $M=\varinjlim M_\lambda$. Cela étant, par hypothèse
$\operatorname{gr}_{\mathfrak J}^*(M)$ est un $A$-module plat~; donc il résulte
de (11.2.9) qu'il existe un $\lambda$ tel que
$\operatorname{gr}_{\mathfrak J_\lambda}^*(M_\lambda)$ soit un
$A_\lambda$-module plat. Pour prouver que $\mathcal F$ est $f$-plat aux points
de $Y$, on peut donc se borner au cas où $S$ est noethérien~; mais alors
$(\mathbf 0_{\mathrm I},6.1.2)$, appliqué par récurrence sur $n$, prouve que les
$\mathcal F/\mathcal J^{n+1}\mathcal F$ sont $f$-plats, et on conclut par
$(\mathbf 0_{\mathrm{III}},10.2.6)$.

\oldpage[IV-3]{134}

La démonstration de (ii) se ramène de la même façon au cas où $S$ est
noethérien, en utilisant (11.2.9)~; la conclusion est alors évidente,
$\operatorname{gr}_{\mathfrak J}^*(B)$ étant dans ce cas une
$(B/\mathfrak J)$-algèbre de type fini et
$\operatorname{gr}_{\mathfrak J}^*(M)$ un
$\operatorname{gr}_{\mathfrak J}^*(B)$-module de type fini.

Pour prouver (iii), on se ramène comme dans (i) au cas où $S$ et $X$ sont
affines~; avec les mêmes notations, on peut supposer, en vertu de (11.2.9) que
$\operatorname{gr}_{\mathfrak J_\lambda}^*(B_\lambda)$ et
$\operatorname{gr}_{\mathfrak J_\lambda}^*(M_\lambda)$ sont des
$A_\lambda$-modules plats et que l'on a
\[
\operatorname{gr}_{\mathfrak J}^*(B)
=\operatorname{gr}_{\mathfrak J_\lambda}^*(B_\lambda)\otimes_{A_\lambda}A
\quad\text{et}\quad
\operatorname{gr}_{\mathfrak J}^*(M)
=\operatorname{gr}_{\mathfrak J_\lambda}^*(M_\lambda)\otimes_{A_\lambda}A.
\]
Par suite $\operatorname{gr}_{\mathfrak J}^*(B)$ est une
$(B/\mathfrak J)$-algèbre de présentation finie et
$\operatorname{gr}_{\mathfrak J}^*(M)$ un
$\operatorname{gr}_{\mathfrak J}^*(B)$-module de présentation finie. Si $W^m$
est l'ensemble des $\mathfrak p\in\operatorname{Spec}(B/\mathfrak J)$ tels que
$\operatorname{gr}_{\mathfrak J}^m(M)_\mathfrak p$ soit un
$(B/\mathfrak J)$-module plat, $W^m$ est ouvert dans
$\operatorname{Spec}(B/\mathfrak J)$ (11.3.1) et l'on a
$W=\bigcap_mW^m$, donc $W$ est stable par générisation. L'assertion (iii)
résulte alors de (11.3.2) appliqué en y prenant
$X=\operatorname{Spec}(\operatorname{gr}_{\mathfrak J}^*(B))$,
$Y=Z=\operatorname{Spec}(B/\mathfrak J)$ et
$\mathcal F=(\operatorname{gr}_{\mathfrak J}^*(M))^{\sim}$.

Enfin, (iv) découle aussitôt de (11.2.9.2).

Généralisant la définition de (6.10.1), on dit que pour un préschéma $X$, un
sous-préschéma fermé $Y$ de $X$ défini par un Idéal quasi-cohérent $\mathcal J$
de $\mathcal O_X$, et un $\mathcal O_X$-Module quasi-cohérent $\mathcal F$,
$\mathcal F$ est \emph{normalement plat le long de $Y$ en un point $y$} si
$(\operatorname{gr}_{\mathcal J}^*(\mathcal F))_y$ est un
$\mathcal O_{Y,y}$-module plat. On dit que $\mathcal F$ est
\emph{normalement plat le long de $Y$} s'il l'est en tous les points de $Y$.
\end{proof}

\begin{corollary}[11.3.5]
\label{IV.11.3.5-fr}
Sous les hypothèses générales de (11.3.4), on suppose que
$\operatorname{gr}_{\mathcal J}^*(\mathcal O_X)$ et
$\operatorname{gr}_{\mathcal J}^*(\mathcal F)$ soient $f$-plats. Alors
l'ensemble $W$ des $y\in Y$ tels que $\mathcal F$ soit normalement plat le long
de $Y$ au point $y$ est ouvert dans $Y$, et (avec les notations de
(11.3.4, (iv))) $\mathcal F'$ est normalement plat le long de $Y'$ en tous les
points de $W'=p^{-1}(W)$.
\end{corollary}

\begin{proposition}[11.3.6]
\label{IV.11.3.6-fr}
Les notations étant celles de (8.5.1) et (8.8.1), on suppose $S_\alpha$
quasi-compact, $X_\alpha$ de présentation finie sur $S_\alpha$, $Y_\alpha$ un
sous-préschéma fermé de $X_\alpha$ défini par un Idéal quasi-cohérent
$\mathcal J_\alpha$ de type fini de $\mathcal O_{X_\alpha}$ tel que
$\operatorname{gr}_{\mathcal J_\alpha}^*(\mathcal O_{X_\alpha})$ soit plat sur
$S_\alpha$~; enfin on suppose que $\mathcal F_\alpha$ est un
$\mathcal O_{X_\alpha}$-Module de présentation finie. Pour que $\mathcal F$
soit normalement plat le long de $Y$, il faut et il suffit qu'il existe
$\lambda$ tel que $\mathcal F_\lambda$ soit normalement plat le long de
$Y_\lambda$.
\end{proposition}

\begin{proof}
Notons que $Y$ (resp. $Y_\lambda$) est le sous-préschéma fermé de $X$ (resp.
$X_\lambda$) défini par $\mathcal J=\mathcal J_\alpha\mathcal O_X$ (resp.
$\mathcal J_\lambda=\mathcal J_\alpha\mathcal O_{X_\lambda}$)~; en vertu de
l'hypothèse et de (11.2.9.2), on a
\[
\operatorname{gr}_{\mathcal J}^*(\mathcal O_X)
=\operatorname{gr}_{\mathcal J_\alpha}^*(\mathcal O_{X_\alpha})
 \otimes_{\mathcal O_{S_\alpha}}\mathcal O_S
\]
et
\[
\operatorname{gr}_{\mathcal J_\lambda}^*(\mathcal O_{X_\lambda})
=\operatorname{gr}_{\mathcal J_\alpha}^*(\mathcal O_{X_\alpha})
 \otimes_{\mathcal O_{S_\alpha}}\mathcal O_{S_\lambda}
\quad\text{pour }\lambda\geq\alpha,
\]
et $\operatorname{gr}_{\mathcal J}^*(\mathcal O_X)$ (resp.
$\operatorname{gr}_{\mathcal J_\lambda}^*(\mathcal O_{X_\lambda})$) est plat
sur $S$ (resp. $S_\lambda$), ce qui entraîne que $Y$ est plat sur $S$ (resp.
$Y_\lambda$ plat sur $S_\lambda$). Si $\mathcal F_\lambda$ est normalement plat
le long de $Y_\lambda$, $\operatorname{gr}_{\mathcal J_\lambda}^*
(\mathcal F_\lambda)|Y_\lambda$ est plat sur $Y_\lambda$, donc aussi sur
$S_\lambda$, et comme
$\operatorname{gr}_{\mathcal J_\lambda}^*(\mathcal F_\lambda)|x_\lambda=0$
pour $x_\lambda\notin Y_\lambda$,
$\operatorname{gr}_{\mathcal J_\lambda}^*(\mathcal F_\lambda)$ est plat sur
$S_\lambda$. On en conclut (11.2.9.2) que
\[
\operatorname{gr}_{\mathcal J}^*(\mathcal F)
=\operatorname{gr}_{\mathcal J_\lambda}^*(\mathcal F_\lambda)
 \otimes_{\mathcal O_{S_\lambda}}\mathcal O_S,
\]
donc $\mathcal F$ est normalement plat le long de $Y$. Pour prouver la
réciproque, on peut supposer que $S_\alpha$ et $X_\alpha$ sont affines et
adopter les notations de (11.2.9)~; puisque
$\operatorname{gr}_{\mathfrak J}^*(B)$ est un $A$-module plat, il en est de même
de $\operatorname{gr}_{\mathfrak J}^*(M)$, donc, en vertu de (11.2.9), il
existe $\lambda\geq\alpha$ tel que
$\operatorname{gr}_{\mathfrak J_\lambda}^*(M_\lambda)$ soit un
$A_\lambda$-module plat, d'où
$\operatorname{gr}_{\mathfrak J}^*(M)
=\operatorname{gr}_{\mathfrak J_\lambda}^*(M_\lambda)\otimes_{A_\lambda}A$.
En outre (11.3.4, (ii)),
$\operatorname{gr}_{\mathfrak J_\lambda}^*(M_\lambda)$ est un
$\operatorname{gr}_{\mathfrak J_\lambda}^*(B_\lambda)$-module de présentation
finie et $\operatorname{gr}_{\mathfrak J_\lambda}^*(B_\lambda)$ une
$(B_\lambda/\mathfrak J_\lambda)$-algèbre de présentation finie. La conclusion
résulte alors de (11.2.6).
\end{proof}

\oldpage[IV-3]{135}

\begin{proposition}[11.3.7]
\label{IV.11.3.7-fr}
Soient $f:X\to Y$ un morphisme localement de présentation finie,
$\mathcal F$, $\mathcal F'$ deux $\mathcal O_X$-Modules quasi-cohérents de
présentation finie, $u:\mathcal F'\to\mathcal F$ un
$\mathcal O_X$-homomorphisme, $x$ un point de $X$ et $y=f(x)$~; on suppose que
$\mathcal F$ est $f$-plat au point $x$. Les conditions suivantes sont alors
équivalentes~:
\begin{enumerate}[label=\emph{\alph*)}]
\item On a $(\ker u)_x=0$ et $\operatorname{Coker}u$ est un
$\mathcal O_X$-Module $f$-plat au point $x$.
\item L'homomorphisme
\[
(u\otimes1)_x:(\mathcal F'\otimes_{\mathcal O_Y}k(y))_x
\longrightarrow(\mathcal F\otimes_{\mathcal O_Y}k(y))_x
\]
est injectif.
\end{enumerate}
Si de plus $\mathcal F$ est $f$-plat, l'ensemble des points $x$ vérifiant les
conditions équivalentes précédentes est ouvert dans $X$.
\end{proposition}

\begin{proof}
La condition a) entraîne b) en vertu de $(\mathbf 0_{\mathrm I},6.1.2)$, sans
hypothèse sur $\mathcal F$. Pour démontrer la réciproque, on peut se borner au
cas où $X$ et $Y$ sont affines, puis, en vertu de (8.9.1) et (11.2.7), supposer
que l'on a
\[
X=X_0\times_{Y_0}Y,\quad f=f_0\times1,\quad
\mathcal F=\mathcal F_0\otimes_{\mathcal O_{Y_0}}\mathcal O_Y,
\]
\[
\mathcal F'=\mathcal F'_0\otimes_{\mathcal O_{Y_0}}\mathcal O_Y,
\qquad u=u_0\otimes1_Y,
\]
où $Y_0$ est noethérien, $f_0:X_0\to Y_0$ un morphisme de type fini,
$\mathcal F_0$, $\mathcal F'_0$ deux $\mathcal O_{X_0}$-Modules cohérents,
$u_0:\mathcal F'_0\to\mathcal F_0$ un homomorphisme~; en outre, si $x_0$ est la
projection de $x$ dans $X_0$, on peut supposer que $\mathcal F_0$ est
$f_0$-plat au point $x_0$. Posons $y_0=f_0(x_0)$~; en vertu de la transitivité
des fibres $(\mathbf I,3.6.4)$, le morphisme projection
$f^{-1}(y)\to f_0^{-1}(y_0)$ est fidèlement plat (2.2.13), et comme
\[
(u\otimes1)_x=((u_0\otimes1)_{x_0})\otimes1_{k(y)},
\]
l'hypothèse que $(u\otimes1)_x$ est injectif entraîne qu'il en est de même de
$(u_0\otimes1)_{x_0}$ (2.2.7). Or, cela entraîne, par
$(\mathbf 0_{\mathrm{III}},10.2.4)$ appliqué aux anneaux locaux noethériens
$\mathcal O_{Y_0,y_0}$ et $\mathcal O_{X_0,x_0}$ et aux
$\mathcal O_{X_0,x_0}$-modules $(\mathcal F'_0)_{x_0}$ et
$(\mathcal F_0)_{x_0}$ (dont le second est plat sur $\mathcal O_{Y_0,y_0}$),
que l'on a $(\ker u_0)_{x_0}=0$ et que $\operatorname{Coker}u_0$ est
$f_0$-plat au point $x_0$. On déduit d'abord de là que
$\operatorname{Coker}u$ est $f$-plat au point $x$ (2.1.4)~; en vertu de
l'exactitude à droite du produit tensoriel, on a d'ailleurs
\[
\operatorname{Coker}u=(\operatorname{Coker}u_0)
 \otimes_{\mathcal O_{Y_0}}\mathcal O_Y~;
\]
appliquant alors $(\mathbf 0_{\mathrm I},6.1.2)$ à la suite (exacte par
hypothèse)
\[
0\longrightarrow(\mathcal F'_0)_{x_0}\longrightarrow
(\mathcal F_0)_{x_0}\longrightarrow(\operatorname{Coker}u_0)_{x_0}
\longrightarrow0,
\]
on en déduit que $u_x$ est un homomorphisme injectif.

Enfin, il résulte de (11.1.2) que l'ensemble $U_0$ des points $z_0\in X_0$ tels
que le morphisme $(u_0\otimes1)_{z_0}$ soit injectif est ouvert~; par platitude
on en déduit que pour tout $z\in X$ au-dessus de $z_0\in U_0$ le morphisme
$(u\otimes1)_z$ est injectif, donc l'ensemble de ces points contient l'image
réciproque $U$ de $U_0$ dans $X$ et est par suite un voisinage de $x$, ce qui
achève la démonstration.
\end{proof}

\begin{theorem}[11.3.8]
\label{IV.11.3.8-fr}
Soient $f:X\to Y$ un morphisme localement de présentation finie,
$\mathcal F$ un $\mathcal O_X$-Module quasi-cohérent de présentation finie,
$(g_i)_{1\leq i\leq n}$ une suite de sections de $\mathcal O_X$ au-dessus de
$X$~; posons
\[
\mathcal F_i=\mathcal F\big/\sum_{j=1}^{i}g_j\mathcal F
\quad\text{pour }0\leq i\leq n
\quad(\text{avec }\mathcal F_0=\mathcal F).
\]
Soient $x$ un point de $X$, $y=f(x)$, et posons
$X_y=f^{-1}(y)=X\times_Y\operatorname{Spec}(k(y))$,
$\mathcal F_y=\mathcal F\otimes_{\mathcal O_Y}k(y)$, qui est un
$\mathcal O_{X_y}$-Module. On suppose que les $(g_i)_x$ appartiennent à l'idéal
maximal $\mathfrak m_x$. Les conditions suivantes sont équivalentes~:
\begin{enumerate}[label=\emph{\alph*)}]
\item La suite $((g_i)_x)$ est $\mathcal F_x$-régulière et les $\mathcal F_i$
$(0\leq i\leq n)$ sont $f$-plats au point $x$.
\item La suite $((g_i)_x)$ est $\mathcal F_x$-régulière et $\mathcal F_n$ est
$f$-plat au point $x$.
\item[\emph{b$'$)}] Il existe un voisinage $U$ de $x$ tel que la suite
$(g_i|U)$ soit $(\mathcal F|U)$-quasi-régulière, et $\mathcal F_n$ est
$f$-plat au point $x$.

\oldpage[IV-3]{136}

\item $\mathcal F$ est $f$-plat au point $x$, et la suite
$((g_i\otimes1)_x)$ d'éléments de $\mathcal O_{X_y,x}$ images des $(g_i)_x$ est
$(\mathcal F_y)_x$-régulière.
\item $\mathcal F$ est $f$-plat au point $x$, et pour tout morphisme
$Y'\to Y$ et tout point $x'\in X'=X\times_YY'$ au-dessus de $x$, si l'on pose
$\mathcal F'=\mathcal F\otimes_YY'$, la suite $(g_i\otimes1)_{x'}$
d'éléments de $\mathcal O_{X',x'}$ est $\mathcal F'_{x'}$-régulière.
\end{enumerate}
En outre l'ensemble des $x\in X$ vérifiant ces conditions est ouvert dans
l'ensemble $Z$ des $x$ tels que $g_i(x)=0$ pour tout $i$.
\end{theorem}

\begin{proof}
Le fait que a) entraîne d) résulte aussitôt de $(\mathbf 0,15.1.15)$, et c) est
un cas particulier de d)~; par ailleurs a) implique b) trivialement. Montrons
que b) ou c) entraîne a). Les $\mathcal F_i$ sont de présentation finie~; le
fait que c) implique a) résulte alors aussitôt de (11.3.7) par récurrence sur
$n$, et cela montre aussi que l'ensemble des $x\in X$ vérifiant c) est ouvert
dans $Z$. Pour montrer que b) entraîne a), on est aussitôt ramené, par
récurrence sur $n$, au cas $n=1$~; nous écrirons $g$ au lieu de $g_1$. La
question étant locale sur $X$ et $Y$, on peut supposer
$Y=\operatorname{Spec}(A)$, $X=\operatorname{Spec}(B)$ affines, $B$ étant une
$A$-algèbre de présentation finie, $\mathcal F=\widetilde M$, où $M$ est un
$B$-module de présentation finie. On peut alors (8.9.1) écrire
$B=B_0\otimes_{A_0}A$, $M=M_0\otimes_{A_0}A$, où $A_0$ est un sous-anneau
noethérien de $A$, $B_0$ une $A_0$-algèbre de type fini et $M_0$ un
$B_0$-module de type fini. En outre $A$ est limite inductive filtrante de ses
sous-anneaux $A_\lambda$ qui sont des $A_0$-algèbres de type fini (donc
noethériens), et si l'on pose
$B_\lambda=B_0\otimes_{A_0}A_\lambda$,
$M_\lambda=M_0\otimes_{A_0}A_\lambda$, $B_\lambda$ est une
$A_\lambda$-algèbre de type fini, $M_\lambda$ un $B_\lambda$-module de type
fini et l'on a $B=\varinjlim B_\lambda$, $M=\varinjlim M_\lambda$. Il existe
donc un indice $\lambda$ et un élément $g_\lambda\in B_\lambda$ tels que
$g=g_\lambda\otimes1$~; revenant aux notations géométriques et posant
$Y_\lambda=\operatorname{Spec}(A_\lambda)$,
$X_\lambda=\operatorname{Spec}(B_\lambda)$ et
$\mathcal F_\lambda=\widetilde{M_\lambda}$, il nous suffira de prouver qu'il y
a un $\mu\geq\lambda$ tel qu'au point $x_\mu\in X_\mu$ projection de $x$,
$(g_\mu)_{x_\mu}$ soit $(\mathcal F_\mu)_{x_\mu}$-régulier et
$\mathcal F_\mu/g_\mu\mathcal F_\mu$ plat sur $Y_\mu$ au point $x_\mu$. On en
déduira en effet, par $(\mathbf 0,15.1.16)$ que $\mathcal F_\mu$ est plat sur
$Y_\mu$ au point $x_\mu$, donc $\mathcal F$ plat sur $Y$ au point $x$.

Le fait que b) entraîne a) résultera donc de la proposition suivante~:

\begin{proposition}[11.3.9]
\label{IV.11.3.9-fr}
Les notations et hypothèses générales étant celles de (8.5.1) et (8.8.1),
supposons que $f_\alpha:X_\alpha\to Y_\alpha$ soit un morphisme localement de
présentation finie. Soient $\mathcal F_\alpha$, $\mathcal G_\alpha$ deux
$\mathcal O_{X_\alpha}$-Modules quasi-cohérents de présentation finie,
$h_\alpha:\mathcal F_\alpha\to\mathcal G_\alpha$ un
$\mathcal O_{X_\alpha}$-homomorphisme, $\mathcal H_\alpha$ son conoyau. Soient
enfin $x$ un point de $X$, $x_\lambda$ sa projection dans $X_\lambda$ pour
$\lambda\geq\alpha$. Pour que $h_x$ soit injectif et que
$\mathcal H=\operatorname{Coker}h$ soit $f$-plat au point $x$, il faut et il
suffit qu'il existe $\lambda\geq\alpha$ tel que $(h_\lambda)_{x_\lambda}$ soit
injectif et que $\mathcal H_\lambda=\operatorname{Coker}h_\lambda$ soit
$f_\lambda$-plat au point $x_\lambda$. En outre, l'ensemble des $x\in X$ ayant
les propriétés précédentes est ouvert dans $X$.
\end{proposition}

Rappelons qu'en vertu de l'exactitude à droite du produit tensoriel, on a
$\mathcal H_\mu=v_{\lambda\mu}^*(\mathcal H_\lambda)$ pour $\lambda\leq\mu$ et
$\mathcal H=v_\lambda^*(\mathcal H_\lambda)$, ce qui justifie les notations.

La suffisance de la condition provient de ce que, si la suite
\[
0\longrightarrow(\mathcal F_\lambda)_{x_\lambda}
\longrightarrow(\mathcal G_\lambda)_{x_\lambda}
\longrightarrow(\mathcal H_\lambda)_{x_\lambda}\longrightarrow0
\]
est exacte et $(\mathcal H_\lambda)_{x_\lambda}$ un
$\mathcal O_{f_\lambda(x_\lambda)}$-module plat, $\mathcal H_x$ est un
$\mathcal O_{f(x)}$-module plat par changement de base (2.1.4) et la suite
$0\to\mathcal F_x\to\mathcal G_x\to\mathcal H_x\to0$ est exacte (2.1.8).
Pour prouver que

\oldpage[IV-3]{137}

la condition est nécessaire, notons que $\mathcal H$ est de présentation
finie~; la question étant locale sur $X$, on peut supposer $X$ et $Y$ affines,
et, en vertu de (11.3.1), supposer que $\mathcal H$ est $f$-plat. Notons
maintenant le

\begin{lemma}[11.3.9.1]
\label{IV.11.3.9.1-fr}
Soient $f:X\to Y$ un morphisme localement de présentation finie,
$\mathcal G$, $\mathcal H$ deux $\mathcal O_X$-Modules quasi-cohérents de
présentation finie tels que $\mathcal H$ soit $f$-plat,
$p:\mathcal G\to\mathcal H$ un homomorphisme surjectif. Alors $\ker(p)$ est un
$\mathcal O_X$-Module de présentation finie.
\end{lemma}

On peut en effet supposer $X$, $Y$ affines, et qu'il existe un morphisme
$Y\to Y_0$ où $Y_0$ est noethérien, un morphisme $f_0:X_0\to Y_0$ de type fini
tels que $X=X_0\times_{Y_0}Y$, $f=(f_0)_{(Y)}$, deux
$\mathcal O_{X_0}$-Modules cohérents $\mathcal G_0$, $\mathcal H_0$ et un
homomorphisme $p_0:\mathcal G_0\to\mathcal H_0$ tels que $\mathcal G$,
$\mathcal H$ et $p$ se déduisent de $\mathcal G_0$, $\mathcal H_0$ et $p_0$ par
changement de base (8.9.1 et 8.5.2). En outre, on peut supposer $p_0$ surjectif
(8.5.7) et $\mathcal H_0$ $f_0$-plat (11.2.7). Alors, si
$\mathcal K_0=\ker(p_0)$, $\mathcal K_0$ est un
$\mathcal O_{X_0}$-Module cohérent $(\mathbf 0_{\mathrm I},5.3.4)$ et en vertu
de $(\mathbf 0_{\mathrm I},6.1.2)$, $\mathcal K=\ker(p)$ se déduit de
$\mathcal K_0$ par changement de base, donc est de présentation finie.

Ce lemme étant démontré, posons $\mathcal K=\ker(\mathcal G\to\mathcal H)$,
qui est donc de présentation finie. On a un homomorphisme canonique
$q:\mathcal F\to\mathcal K$, et par hypothèse
$q_x:\mathcal F_x\to\mathcal K_x$ est un isomorphisme~; par suite
$(\mathbf 0_{\mathrm I},5.2.7)$ il existe un voisinage $U$ de $x$ tel que
$q|U$ soit un isomorphisme, et en restreignant $X$, on peut supposer que la
suite
\[
0\longrightarrow\mathcal F\xrightarrow{h}\mathcal G
\longrightarrow\mathcal H\longrightarrow0
\]
est \emph{exacte}. Cela étant, il résulte de (11.2.6) que pour un $\lambda$
assez grand, $\mathcal H_\lambda$ est un $\mathcal O_{X_\lambda}$-Module plat~;
posons $\mathcal K_\lambda=\ker(\mathcal G_\lambda\to\mathcal H_\lambda)$ et
$\mathcal K_\mu=v_{\lambda\mu}^*(\mathcal K_\lambda)$ pour $\mu\geq\lambda$~;
il résulte de $(\mathbf 0_{\mathrm I},6.1.2)$ que l'on a
$\mathcal K_\mu=\ker(\mathcal G_\mu\to\mathcal H_\mu)$ et
$\mathcal K=v_\lambda^*(\mathcal K_\lambda)
=\ker(\mathcal G\to\mathcal H)=\mathcal F$. On a d'autre part pour tout
$\mu\geq\lambda$ un homomorphisme canonique
$q_\mu:\mathcal F_\mu\to\mathcal K_\mu$ avec
$q_\mu=v_{\lambda\mu}^*(q_\lambda)$ pour $\lambda\leq\mu$, et
$q=v_\mu^*(q_\mu)$. Comme $q$ est un isomorphisme, il résulte de (8.5.2.4) et
(11.3.9.1) qu'il existe $\mu\geq\lambda$ tel que $q_\mu$ soit un isomorphisme,
et par suite $u_\mu:\mathcal F_\mu\to\mathcal G_\mu$ un homomorphisme
injectif.

Revenons à la démonstration de (11.3.8).

Puisque l'ensemble des $x\in X$ vérifiant b) est ouvert dans $X$, il est clair
que b) entraîne b$'$. Montrons enfin que b$'$ entraîne c). En premier lieu,
$\mathcal F_n$ est $f$-plat dans un voisinage de $x$ (11.3.1) et on peut donc
se borner au cas où $U=X$ et où $\mathcal F_n$ est $f$-plat. Comme par
définition $\operatorname{gr}_{\mathcal J}^*(\mathcal F)$ est isomorphe à
$\mathcal F_n\otimes_{\mathcal O_Y}\mathcal O_Y[T_1,\ldots,T_n]$
$(\mathbf 0,15.2.2)$, il est $f$-plat, et on en conclut par (11.3.4, (i)) que
$\mathcal F$ lui-même est $f$-plat au voisinage de $x$. D'autre part, si
$(\mathcal J_y)_x$ est l'idéal de $\mathcal O_{X_y,x}$ engendré par les
$(g_i\otimes1)_x$, il résulte de (11.2.9.2) que, dans le diagramme
\[
\xymatrix{
(\operatorname{gr}_{\mathcal J_x}^0(\mathcal F_x)[T_1,\ldots,T_n])
 \otimes_{\mathcal O_y}k(y) \ar[r] \ar[d] &
\operatorname{gr}_{\mathcal J_x}^*(\mathcal F_x)
 \otimes_{\mathcal O_y}k(y) \ar[d] \\
(\operatorname{gr}_{(\mathcal J_y)_x}^0((\mathcal F_y)_x))
 [T_1,\ldots,T_n] \ar[r] &
\operatorname{gr}_{(\mathcal J_y)_x}^*((\mathcal F_y)_x)
}
\]

\oldpage[IV-3]{138}

les flèches verticales sont des isomorphismes~; comme la première ligne est un
isomorphisme, il en est de même de la seconde, donc la suite
$((g_i\otimes1)_x)$ est $(\mathcal F_y)_x$-quasi-régulière, et par suite aussi
$(\mathcal F_y)_x$-régulière, puisque $X_y$ est localement de type fini sur
$k(y)$. C.Q.F.D.
\end{proof}

\begin{theorem}[11.3.10]
\label{IV.11.3.10-fr}
(critère de platitude par fibres). --- Soient $S$ un préschéma,
$g:X\to S$, $h:Y\to S$ deux morphismes, $f:X\to Y$ un $S$-morphisme,
$\mathcal F$ un $\mathcal O_X$-Module quasi-cohérent, $x$ un point de $X$,
$y=f(x)$, $s=h(y)=g(x)$. On suppose vérifiée l'une des deux hypothèses
suivantes~:
\begin{enumerate}[label=\arabic{enumi}\textsuperscript{o}]
\item $S$, $X$ et $Y$ sont localement noethériens et $\mathcal F$ est cohérent.
\item $g$ et $h$ sont localement de présentation finie et $\mathcal F$ est de
présentation finie.
\end{enumerate}
Alors, avec les notations de (9.4.1), si $\mathcal F_x\neq0$, les conditions
suivantes sont équivalentes~:
\begin{enumerate}[label=\emph{\alph*)}]
\item $\mathcal F$ est $g$-plat au point $x$ et $\mathcal F_s$ est $f_s$-plat
au point $x$.
\item $h$ est plat au point $y$ et $\mathcal F$ est $f$-plat au point $x$.
\end{enumerate}
En outre, dans l'hypothèse 2\textsuperscript{o}, l'ensemble des $x\in X$
vérifiant la condition b) est ouvert dans $X$.
\end{theorem}

\begin{proof}
La dernière assertion résulte de (11.3.1) appliqué à $\mathcal O_Y$ et au
morphisme $h$ d'une part, à $\mathcal F$ et au morphisme $f$ (qui est
localement de présentation finie) de l'autre.

Comme $g=h\circ f$, il est clair que b) implique a) sans supposer
1\textsuperscript{o} ni 2\textsuperscript{o} (2.1.6 et 2.1.4). Pour prouver que
a) entraîne b), on peut se borner au cas où $S$, $X$ et $Y$ sont affines~;
sous l'hypothèse 2\textsuperscript{o}, en appliquant (11.2.7), on se ramène au
cas où de plus $S$ est noethérien, c'est-à-dire qu'on peut se borner à
considérer le cas où l'hypothèse 1\textsuperscript{o} est satisfaite. Alors
l'assertion équivaut au lemme suivant, qui améliore
$(\mathbf 0_{\mathrm{III}},10.2.5)$~:

\begin{lemma}[11.3.10.1]
\label{IV.11.3.10.1-fr}
Soient $\rho:A\to B$, $\sigma:B\to C$ des homomorphismes locaux d'anneaux
locaux noethériens, $k$ le corps résiduel de $A$, $M$ un $C$-module
$\neq0$ de type fini. Alors les conditions suivantes sont équivalentes~:
\begin{enumerate}[label=\emph{\alph*)}]
\item $M$ est un $A$-module plat et $M\otimes_Ak$ est un
$(B\otimes_Ak)$-module plat.
\item $B$ est un $A$-module plat et $M$ est un $B$-module plat.
\end{enumerate}
\end{lemma}

Nous établirons d'abord le lemme plus général suivant~:

\begin{lemma}[11.3.10.2]
\label{IV.11.3.10.2-fr}
Soient $A$ un anneau commutatif, $B$ une $A$-algèbre commutative,
$\mathfrak J$ un idéal de $A$, $M$ un $B$-module. On considère d'une part les
conditions suivantes~:
\begin{enumerate}[label=\emph{(\roman*)}]
\item $\mathfrak J$ est nilpotent.
\item $B$ est noethérien et $M$ est idéalement séparé pour la topologie
$\mathfrak JB$-préadiques $(\mathbf 0_{\mathrm{III}},10.2.1)$.
\item $\mathfrak JB$ est contenu dans le radical de $B$.
\end{enumerate}
On considère d'autre part les quatre propriétés~:
\begin{enumerate}[label=\emph{\alph*)}]
\item $M$ est un $B$-module plat.
\item $B$ est un $A$-module plat.
\item $M$ est un $A$-module plat et $M/\mathfrak JM$ un
$(B/\mathfrak JB)$-module plat.
\item $B/\mathfrak JB$ est un $(A/\mathfrak J)$-module plat et pour tout idéal
maximal $\mathfrak m\supset\mathfrak JB$ de $B$ on a $\mathfrak mM\neq M$.
\end{enumerate}
Alors~:
\begin{enumerate}[label=\arabic{enumi}\textsuperscript{o}]
\item Si l'une des conditions (i), (ii) est vérifiée, la conjonction de a) et
b) implique c), et c) implique a).

\oldpage[IV-3]{139}

\item Si la condition (i) ou la conjonction de (ii) et (iii) est vérifiée, la
conjonction de c) et d) implique la conjonction de a) et b).
\end{enumerate}
\end{lemma}

\noindent 1\textsuperscript{o} La première assertion est immédiate
$(\mathbf 0_{\mathrm I},6.2.1)$. Supposons donc c) vérifiée, et prouvons a).
Considérons les anneaux gradués $\operatorname{gr}_{\mathfrak J}^*(A)$,
$\operatorname{gr}_{\mathfrak J}^*(B)$ et le module gradué
$\operatorname{gr}_{\mathfrak J}^*(M)$ (à la fois sur
$\operatorname{gr}_{\mathfrak J}^*(A)$ et
$\operatorname{gr}_{\mathfrak J}^*(B)$) relatifs aux filtrations
$\mathfrak J$-préadique, ainsi que les applications canoniques surjectives
$(\mathbf 0_{\mathrm{III}},10.1.1.2)$
\begin{align*}
u&:\operatorname{gr}^0(B)\otimes_{\operatorname{gr}^0(A)}
   \operatorname{gr}^*(A)\longrightarrow\operatorname{gr}^*(B),\\
v&:\operatorname{gr}^0(M)\otimes_{\operatorname{gr}^0(A)}
   \operatorname{gr}^*(A)\longrightarrow\operatorname{gr}^*(M),\\
w&:\operatorname{gr}^0(M)\otimes_{\operatorname{gr}^0(B)}
   \operatorname{gr}^*(B)\longrightarrow\operatorname{gr}^*(M).
\end{align*}
Il est clair qu'on a un diagramme commutatif
\begin{equation}
\xymatrix{
\operatorname{gr}^0(M)\otimes_{\operatorname{gr}^0(A)}\operatorname{gr}^*(A)
 \ar[rr]^v \ar[dr]_{1\otimes u} && \operatorname{gr}^*(M) \\
& \operatorname{gr}^0(M)\otimes_{\operatorname{gr}^0(B)}
  \operatorname{gr}^*(B) \ar[ur]_w &
}
\label{IV.11.3.10.3-fr}
\tag{11.3.10.3}
\end{equation}
L'hypothèse c) entraîne que $v$ est bijective
$(\mathbf 0_{\mathrm{III}},10.2.1)$~; comme les deux autres applications du
diagramme sont surjectives, elles sont aussi bijectives. Mais comme en vertu de
l'hypothèse c), $M/\mathfrak JM$ est un $(B/\mathfrak JB)$-module plat, il
résulte de $(\mathbf 0_{\mathrm{III}},10.2.1)$ que $M$ est un $B$-module plat.

\noindent 2\textsuperscript{o} L'une ou l'autre des conditions (i), (iii)
implique que tout idéal maximal de $B$ contient $\mathfrak JB$. Il résulte donc
de 1\textsuperscript{o} et de la conjonction de c) et d) que $M$ est un
$B$-module \emph{fidèlement plat}, et par suite $\operatorname{gr}^0(M)$ un
$\operatorname{gr}^0(B)$-module \emph{fidèlement plat}
$(\mathbf 0_{\mathrm I},6.2.1)$. On a vu dans 1\textsuperscript{o} que
l'hypothèse c) entraîne que les trois applications du diagramme (11.3.10.3)
sont bijectives~; le fait que $\operatorname{gr}^0(M)$ soit un
$\operatorname{gr}^0(B)$-module fidèlement plat implique donc que $u$ est aussi
bijective $(\mathbf 0_{\mathrm I},6.4.1)$. D'autre part, les conditions (ii) et
(iii) impliquent que $B$ est un $A$-module idéalement séparé pour la filtration
$\mathfrak J$-préadique (Bourbaki, \emph{Alg. comm.}, chap.~III, \S~5,
n$^\circ$~4, prop.~2)~; on déduit donc encore de
$(\mathbf 0_{\mathrm{III}},10.2.1)$ que si la condition (i), ou la conjonction
de (ii) et (iii), est vérifiée, $B$ est un $A$-module plat.

\phantomsection\label{IV.11.3.10.4-fr}
\noindent\textbf{(11.3.10.4)} Le lemme (11.3.10.2) étant établi, on en déduit
(11.3.10.1) en prenant pour $\mathfrak J$ l'idéal maximal de $A$, et en
remarquant qu'alors les conditions (ii) et (iii) de (11.3.10.2) sont
satisfaites (Bourbaki, \emph{Alg. comm.}, chap.~III, \S~5, n$^\circ$~4,
prop.~2). Ceci achève donc aussi la démonstration de (11.3.7).
\end{proof}

\begin{corollary}[11.3.11]
\label{IV.11.3.11-fr}
Soient $g:X\to S$, $h:Y\to S$ deux morphismes localement de présentation
finie, $f:X\to Y$ un $S$-morphisme. Les conditions suivantes sont
équivalentes~:
\begin{enumerate}[label=\emph{\alph*)}]
\item $g$ est plat et $f_s:X_s\to Y_s$ est plat pour tout $s\in S$.
\item $h$ est plat en tous les points de $f(X)$ et $f$ est plat.
\end{enumerate}
\end{corollary}

Il suffit d'appliquer (11.3.10) pour $\mathcal F=\mathcal O_X$.

\begin{remark}[11.3.12]
\label{IV.11.3.12-fr}
Il serait intéressant de pouvoir donner de (11.3.2) et (11.3.10) des
démonstrations n'utilisant pas le passage à la limite inductive~; il suffirait
pour cela de démontrer le critère suivant~:

Soient $A$ un anneau, $B$ une $A$-algèbre de présentation finie, $M$ un
$B$-module de présentation finie, $\mathfrak J$ un

\oldpage[IV-3]{140}

idéal de $A$, $\mathfrak p$ un idéal premier de $B$ contenant $\mathfrak JB$.
Alors les deux conditions suivantes sont équivalentes~:
\begin{enumerate}[label=\emph{\alph*)}]
\item $M_\mathfrak p$ est un $A$-module plat.
\item $M_\mathfrak p/\mathfrak JM_\mathfrak p$ est un
$(A/\mathfrak J)$-module plat et
$\operatorname{Tor}_1^A(A/\mathfrak J,M_\mathfrak p)=0$.
\end{enumerate}
Lorsque $A$ est \emph{noethérien}, c'est une conséquence de
$(\mathbf 0_{\mathrm{III}},10.2.2)$ appliqué à la $A$-algèbre noethérienne
$B_\mathfrak p$~; peut-on en déduire l'énoncé général par un passage à la
limite inductive~?

Il convient de noter à ce propos qu'une telle généralisation n'est certainement
pas possible lorsqu'on remplace la condition b) ci-dessus par l'une des
conditions c), d) de $(\mathbf 0_{\mathrm{III}},10.2.1)$ où $M$ est remplacé
par $M_\mathfrak p$. Prenons par exemple pour $A$ un anneau local dont l'idéal
maximal $\mathfrak J$ est principal et tel que l'intersection
\[
\mathfrak R=\bigcap_{n\geq0}\mathfrak J^{n+1}
\]
ne soit pas réduite à $0$ (par exemple l'anneau des germes de fonctions
numériques indéfiniment différentiables au voisinage de $0$ dans $\mathbf R$).
Prenons $B=A$, $\mathfrak p=\mathfrak J$, et $M=A/\mathfrak N$, où
$\mathfrak N$ est un sous-module monogène $\neq0$ de $\mathfrak R$~; $M$ est
donc de présentation finie. Il est clair que $M$ n'est pas un $A$-module plat,
car étant de présentation finie, il serait libre (Bourbaki, \emph{Alg. comm.},
chap.~II, \S~3, n$^\circ$~2, cor.~2 de la prop.~5), ce qui est absurde puisque
$\mathfrak N\neq0$. Cependant,
\[
\mathfrak J^kM/\mathfrak J^{k+n}M\simeq A/\mathfrak J^n
\]
quels que soient $k$ et $n$ positifs, donc $M$ vérifie bien les conditions c)
et d) de $(\mathbf 0_{\mathrm{III}},10.2.1)$ puisque $A/\mathfrak J$ est un
corps.
\end{remark}
\begin{proposition}[11.3.13]
\label{IV.11.3.13-fr}
Soient $f:X\to Y$ un morphisme de préschémas, $x$ un point de $X$ tel que
$f$ soit plat au point $x$~; posons $y=f(x)$.
\begin{enumerate}[label=\emph{(\roman*)}]
\item Si $X$ est réduit (resp. normal) au point $x$, alors $Y$ est réduit
(resp. normal) au point $y$.
\item On suppose de plus que $f$ soit de présentation finie au point $x$.
Alors, si $Y$ est réduit (resp. normal) au point $y$ et si le morphisme $f$
est réduit (resp. normal) au point $x$ (6.8.1), $X$ est réduit (resp. normal)
au point $x$.
\end{enumerate}
\end{proposition}

\begin{proof}
La première assertion n'est mise que pour mémoire, ayant été démontrée dans
(2.1.13). Pour prouver (ii), on peut se borner au cas où $X=\operatorname{Spec}(B)$,
$Y=\operatorname{Spec}(A)$, où $A$ est un anneau local réduit (resp. intègre et
intégralement clos) et $B$ une $A$-algèbre de présentation finie. On peut alors
(8.9.1) écrire
\[
B=B_0\otimes_{A_0}A,
\]
où $A_0$ est une sous-$\mathbf Z$-algèbre de type fini de $A$ et $B_0$ une
$A_0$-algèbre de type fini. Soit $(C_\alpha)$ la famille filtrante croissante
des sous-$A_0$-algèbres de type fini de $A$, qui sont donc des
$\mathbf Z$-algèbres de type fini~; on a $A=\varinjlim C_\alpha$. Distinguons
maintenant les deux cas~:

\noindent 1\textsuperscript{o} Supposons $A$ réduit et $f$ réduit au point
$x$. Si $\mathfrak m$ est l'idéal maximal de $A$, soit
$\mathfrak p_\alpha$ l'idéal premier $\mathfrak m\cap C_\alpha$, et posons
$A'_\alpha=(C_\alpha)_{\mathfrak p_\alpha}$, de sorte que l'on a aussi
$A=\varinjlim A'_\alpha$ (5.13.3). Posons
$Y_\alpha=\operatorname{Spec}(A'_\alpha)$,
$X_\alpha=\operatorname{Spec}(B_0\otimes_{A_0}A'_\alpha)$ et soient
$f_\alpha:X_\alpha\to Y_\alpha$, $x_\alpha$ (resp. $y_\alpha$) la projection
de $x$ (resp. $y$) dans $X_\alpha$ (resp. $Y_\alpha$). Puisque $f$ est réduit
au point $x$, il existe $\alpha_0$ tel que $f_\alpha$ soit réduit au point
$x_\alpha$ pour $\alpha\geq\alpha_0$ ((6.7.8) et (11.2.6))~; comme
$Y_\alpha$ et $X_\alpha$ sont noethériens et que $A'_\alpha$ est réduit
(puisqu'il en est ainsi de $C_\alpha$), on déduit de (3.3.5) que $X_\alpha$
est réduit au point $x_\alpha$. Mais puisque
$B=\varinjlim(B_0\otimes_{A_0}A'_\alpha)$, on a aussi
$\mathcal O_{X,x}=\varinjlim\mathcal O_{X_\alpha,x_\alpha}$ (5.13.3) et par
suite $\mathcal O_{X,x}$ est réduit (5.13.2).

\noindent 2\textsuperscript{o} Supposons $A$ intégralement clos et $f$ normal
au point $x$. Comme $C_\alpha$ est universellement japonais (7.7.4), sa
clôture intégrale $C'_\alpha$ est une $C_\alpha$-algèbre finie, évidemment
contenue dans $A$. Soit $\mathfrak p'_\alpha$ l'idéal premier
$\mathfrak m\cap C'_\alpha$, et posons
$A''_\alpha=(C'_\alpha)_{\mathfrak p'_\alpha}$, de sorte que $A''_\alpha$
est un anneau local noethérien intègre et intégralement clos, et que l'on a
$A=\varinjlim A''_\alpha$ (5.13.3). Posons
$Y'_\alpha=\operatorname{Spec}(A''_\alpha)$,
$X'_\alpha=\operatorname{Spec}(B_0\otimes_{A_0}A''_\alpha)$ et soient
$f'_\alpha:X'_\alpha\to Y'_\alpha$, $x'_\alpha$ (resp. $y'_\alpha$) la
projection de $x$ (resp. $y$) dans $X'_\alpha$ (resp. $Y'_\alpha$). Puisque
$f$ est normal au point $x$, il existe $\alpha_0$ tel que, pour
$\alpha\geq\alpha_0$, $f'_\alpha$ soit normal au point $x'_\alpha$ ((6.7.8)
et (11.2.6))~;

\oldpage[IV-3]{141}

comme $X'_\alpha$ et $Y'_\alpha$ sont noethériens et que $A''_\alpha$ est
intégralement clos, on déduit de (6.5.4) que $X'_\alpha$ est normal au point
$x'_\alpha$. Mais les morphismes
$\operatorname{Spec}(A''_\beta)\to\operatorname{Spec}(A''_\alpha)$ pour
$\alpha\leq\beta$ sont dominants, donc, comme en vertu de (11.3.1) on peut
supposer que les $f'_\alpha$ sont plats pour $\alpha\geq\alpha_0$, il résulte
de (2.3.7) que toute composante irréductible de $X'_\beta$ domine une
composante irréductible de $X'_\alpha$. On conclut alors de (5.13.4), appliqué
aux préschémas
\[
\operatorname{Spec}(\mathcal O_{X_0,x_0}\otimes_{A_0}A''_\alpha),
\]
que $X$ est normal au point $x$.
\end{proof}

\begin{corollary}[11.3.14]
\label{IV.11.3.14-fr}
Soient $f:X\to Y$ un morphisme localement de présentation finie, $x$ un point
de $X$, $y=f(x)$. Si $Y$ est géométriquement unibranche au point $y$ (6.15.1)
et si le morphisme $f$ est normal au point $x$ (6.8.1), alors $X$ est
géométriquement unibranche au point $x$.
\end{corollary}

\begin{proof}
On peut évidemment se borner au cas où $Y=\operatorname{Spec}(A)$, $A$ étant
un anneau local intègre géométriquement unibranche, $y$ étant le point fermé de
$Y$. Soit $A'$ la clôture intégrale de $A$~; posons
$Y'=\operatorname{Spec}(A')$ et soit $y'$ le point fermé de $Y'$, de sorte que
le morphisme $g:Y'\to Y$ est radiciel au point $y$ (6.15.3), entier et
birationnel. Si l'on pose $X'=X\times_Y Y'$, et si $f':X'\to Y'$ et
$g':X'\to X$ sont les projections, $g'$ est entier et radiciel au point $x$
(6.15.3.1). Soit donc $x'$ l'unique point de $g'^{-1}(x)$, qui est au-dessus
de $y'$. Le morphisme $f'$ est de présentation finie et normal au point $x'$
(6.7.8), et par suite l'anneau local $\mathcal O_{X',x'}$ est intègre et
intégralement clos (11.3.13). D'autre part, on peut supposer $f$ plat (11.3.1),
donc $g'$ est un morphisme birationnel (6.15.4.1)~; par suite
$\operatorname{Spec}(\mathcal O_{X,x})$ est irréductible, et comme
$\mathcal O_{X,x}$ est réduit (11.3.13) il est intègre et géométriquement
unibranche en vertu de (6.15.5).
\end{proof}

\begin{proposition}[11.3.15]
\label{IV.11.3.15-fr}
Soient $A$ un anneau, $B$ une $A$-algèbre de présentation finie, $M$ un
$B$-module de présentation finie, qui est un $A$-module plat. Alors il existe
une suite finie $(f_i)_{1\leq i\leq n}$ d'éléments de $A$, tels que l'idéal
engendré par les $f_i$ soit égal à $A$, et que, pour $0\leq i<n$,
\[
M_{f_{i+1}}\Big/\Big(\sum_{j=1}^{i}f_jM_{f_{i+1}}\Big)
\]
soit un
$A_{f_{i+1}}\big/(\sum_{j=1}^{i}f_jA_{f_{i+1}})$-module libre.
\end{proposition}

On peut encore dire que les $D(f_i)$ forment un recouvrement ouvert de
$X=\operatorname{Spec}(A)$, et que si l'on pose $Z_1=D(f_1)$, puis, par
récurrence,
\[
Z_{i+1}=D(f_{i+1})\cap V(f_1A+\cdots+f_iA),
\]
les $Z_i$ forment une \emph{partition} de $X$ en ensembles localement fermés,
$A_{f_{i+1}}/(\sum_{j=1}^{i}f_jA_{f_{i+1}})$ étant l'anneau d'un sous-schéma
affine de $X$ ayant $Z_{i+1}$ pour espace sous-jacent.

\begin{proof}
Pour démontrer la proposition, on peut d'abord, en vertu de (11.2.7), supposer
qu'il existe un sous-anneau noethérien $A_0$ de $A$, une $A_0$-algèbre de type
fini $B_0$ et un $B_0$-module de type fini $M_0$, plat sur $A_0$ et tels que
$B=B_0\otimes_{A_0}A$, $M=M_0\otimes_{A_0}A$~; il est clair qu'il suffira de
prouver la proposition pour $A_0$, $B_0$ et $M_0$, car si les éléments
$f_i\in A_0$ vérifient dans ce cas les conditions de l'énoncé, ils les
vérifieront aussi pour $A$, $B$, $M$, puisque
\[
M_{f_{i+1}}\Big/\Big(\sum_{j=1}^{i}f_jM_{f_{i+1}}\Big)
=\left((M_0)_{f_{i+1}}\Big/\Big(\sum_{j=1}^{i}f_j(M_0)_{f_{i+1}}\Big)\right)
\otimes_{A_0}A
\]
(Bourbaki, \emph{Alg. comm.}, chap.~II, \S~2, n$^\circ$~7, prop.~18). On
peut donc se borner désormais au cas où $A$ est noethérien.

Notons maintenant que si $C$ est un anneau noethérien, $\mathfrak N$ son
nilradical et $P$ un $C$-module plat, alors il résulte de
$(\mathbf 0_{\mathrm{III}},10.1.2)$ que pour que $P$ soit un $C$-module libre,

\oldpage[IV-3]{142}

il faut et il suffit que $P\otimes_C(C/\mathfrak N)$ soit un
$(C/\mathfrak N)$-module libre. Notons d'autre part que si $C$ est un anneau
noethérien réduit, il existe $g\in C$ tel que $C_g$ soit un anneau intègre.
Utilisons maintenant le lemme (6.9.2)~: on peut définir par récurrence une suite
$(f_i)_{i\geq1}$ d'éléments de $A$ de la façon suivante~:
\begin{enumerate}
\item[$1^\circ$] $f_1$ est tel que $A_1=(A_{\mathrm{red}})_{f_1}$ soit
intègre et $M\otimes_A A_1$ un $A_1$-module libre~;
\item[$2^\circ$] si l'idéal $\mathfrak J_i$ engendré par $f_1,\ldots,f_i$
est $A$, $f_{i+1}=f_i$~;
\item[$3^\circ$] si au contraire $\mathfrak J_i\neq A$, $f_{i+1}$ est un
élément n'appartenant pas à $\mathfrak J_i$ tel que
$A_{i+1}=((A/\mathfrak J_i)_{\mathrm{red}})_{f_{i+1}}$ soit intègre et
$M\otimes_A A_{i+1}$ un $A_{i+1}$-module libre.
\end{enumerate}
Comme $A$ est noethérien, la suite croissante $(\mathfrak J_i)$ est
stationnaire, donc il existe $n$ tel que $f_1,\ldots,f_n$ engendrent l'idéal
$A$, et les $f_i$ répondent à la question puisque
\[
((A/\mathfrak J_i)_{\mathrm{red}})_{f_{i+1}}
=((A/\mathfrak J_i)_{f_{i+1}})_{\mathrm{red}}.
\]
\end{proof}

\begin{proposition}[11.3.16]
\label{IV.11.3.16-fr}
Soient $f:X\to Y$ un morphisme fidèlement plat de présentation finie,
$g:Y\to Z$ un morphisme tel que $g\circ f:X\to Z$ soit un morphisme de type
fini (resp. de présentation finie). Alors $g$ est un morphisme de type fini
(resp. de présentation finie).
\end{proposition}

\begin{proof}
Comme $f$ est surjectif et $g\circ f$ quasi-compact, $g$ est quasi-compact
(1.1.3). Montrons d'abord que si $g\circ f$ est de présentation finie, $g$ est
quasi-séparé. En effet, considérons un ouvert affine $W\subset Z$, et soit
$(V_i)$ un recouvrement affine fini de $g^{-1}(W)$~; il s'agit de voir que les
$V_i\cap V_j$ sont quasi-compacts (1.2.6 et 1.2.7). Pour chaque $i$, soit
$(U_{ih})$ un recouvrement ouvert affine fini de $f^{-1}(V_i)$~; comme $f$ est
de présentation finie, les $U_{ih}\cap U_{jk}$ sont tous quasi-compacts~; or,
puisque $f$ est surjectif, $V_i\cap V_j$ est réunion des images
$f(U_{ih}\cap U_{jk})$ pour $h$, $k$ variables, donc est quasi-compact puisque
$f$ est continue.

La question est donc locale sur $Y$ et on peut supposer que
$Y=\operatorname{Spec}(B)$ et $Z=\operatorname{Spec}(A)$ sont affines, $X$
étant réunion finie d'ouverts affines $X_i=\operatorname{Spec}(C_i)$, où les
$C_i$ sont des $B$-algèbres de type fini (resp. de présentation finie). Si
$X'$ est le préschéma somme des $X_i$, $p:X'\to X$ le morphisme qui coïncide
avec l'injection canonique dans chaque $X_i$, $p$ est un morphisme fidèlement
plat de présentation finie (1.6.5), donc $g\circ f\circ p$ est un morphisme de
type fini (resp. de présentation finie) et $f\circ p$ un morphisme fidèlement
plat de présentation finie. On peut par suite supposer aussi que
$X=\operatorname{Spec}(C)$ est affine, et on est donc ramené à prouver le
corollaire suivant.
\end{proof}

\begin{corollary}[11.3.17]
\label{IV.11.3.17-fr}
Soient $A$ un anneau, $B$ une $A$-algèbre, $C$ une $B$-algèbre de présentation
finie et qui soit un $B$-module fidèlement plat. Alors, si $C$ est une
$A$-algèbre de type fini (resp. de présentation finie), $B$ est une
$A$-algèbre de type fini (resp. de présentation finie).
\end{corollary}

\begin{proof}
Supposons d'abord que $g\circ f$ soit de type fini. Soit $(B_\lambda)$ la
famille filtrante croissante des sous-$A$-algèbres de type fini de $B$~; en
vertu de (8.8.2), il existe un indice $\alpha$ tel que
$C=C_\alpha\otimes_{B_\alpha}B$, où $C_\alpha$ est une
$B_\alpha$-algèbre de présentation finie~; en outre (8.10.5 et 11.2.6) on peut
supposer que $C_\alpha$ est un $B_\alpha$-module fidèlement plat. Pour
$\lambda\geq\alpha$, on a donc
$C=C_\lambda\otimes_{B_\lambda}B$, où $C_\lambda$ est un
$B_\lambda$-module fidèlement plat~; comme l'application $B_\lambda\to B$ est
injective, on en déduit qu'il en est de même de $C_\lambda\to C$. En outre,
comme $C=\varinjlim C_\lambda$, tout élément de $C$ appartient à un
$C_\lambda$, et par suite il existe $\lambda$ tel que l'application
$C_\lambda\to C$

\oldpage[IV-3]{143}

soit bijective, puisque $C$ est une $B$-algèbre de type fini. Mais alors
l'application $B_\lambda\to B$ est bijective par fidèle platitude, donc
$B=B_\lambda$ est une $A$-algèbre de type fini.

Supposons maintenant que $g\circ f$ soit de présentation finie, $C$ étant donc
une $A$-algèbre de présentation finie. D'après la première partie du
raisonnement, on a
\[
B=A[T_1,\ldots,T_n]/\widetilde{\mathfrak J}
\]
pour un idéal $\widetilde{\mathfrak J}$~; soit
$(\widetilde{\mathfrak J}_\lambda)$ la famille filtrante des idéaux de type
fini de $A[T_1,\ldots,T_n]$ contenus dans $\widetilde{\mathfrak J}$, de sorte
que $\widetilde{\mathfrak J}=\varinjlim\widetilde{\mathfrak J}_\lambda$ et
$B=\varinjlim B_\lambda$, avec
$B_\lambda=A[T_1,\ldots,T_n]/\widetilde{\mathfrak J}_\lambda$. Appliquant
comme ci-dessus (8.8.2), (8.10.5) et (11.2.6), on a
$C=C_\alpha\otimes_{B_\alpha}B$, où $C_\alpha$ est une
$B_\alpha$-algèbre de présentation finie et un $B_\alpha$-module fidèlement
plat~; on posera encore $C_\lambda=C_\alpha\otimes_{B_\alpha}B_\lambda$ pour
$\lambda\geq\alpha$, de sorte que
\[
C_\lambda=C_\alpha/\big(C_\alpha\otimes_{B_\alpha}
(\widetilde{\mathfrak J}_\lambda/\widetilde{\mathfrak J}_\alpha)\big)
\]
par platitude, et de même
\[
C=C_\alpha/\big(C_\alpha\otimes_{B_\alpha}
(\widetilde{\mathfrak J}/\widetilde{\mathfrak J}_\alpha)\big).
\]
Comme par hypothèse $C$ est une $A$-algèbre de présentation finie ainsi que
$C_\alpha$, $C_\alpha\otimes_{B_\alpha}
(\widetilde{\mathfrak J}/\widetilde{\mathfrak J}_\alpha)$ est un idéal de
type fini de $C_\alpha$ (1.4.4), et les
$C_\alpha\otimes_{B_\alpha}
(\widetilde{\mathfrak J}_\lambda/\widetilde{\mathfrak J}_\alpha)$
s'identifient à des idéaux de type fini de $C_\alpha$ dont
$C_\alpha\otimes_{B_\alpha}
(\widetilde{\mathfrak J}/\widetilde{\mathfrak J}_\alpha)$ est la réunion. Il
existe par suite un $\lambda\geq\alpha$ tel que
\[
C_\alpha\otimes_{B_\alpha}
(\widetilde{\mathfrak J}/\widetilde{\mathfrak J}_\alpha)
=C_\alpha\otimes_{B_\alpha}
(\widetilde{\mathfrak J}_\lambda/\widetilde{\mathfrak J}_\alpha),
\]
d'où $\widetilde{\mathfrak J}=\widetilde{\mathfrak J}_\lambda$ par fidèle
platitude, ce qui prouve que $B$ est une $A$-algèbre de présentation finie.
\end{proof}

\subsection{Descente de la platitude par des morphismes quelconques~: cas d'un préschéma de base artinien}
\label{subsection:IV.11.4-fr}

\begin{theorem}[11.4.1]
\label{IV.11.4.1-fr}
Soient $A$ un anneau, $\mathfrak J$ un idéal nilpotent de $A$,
$u_\lambda:A\to B_\lambda$ $(\lambda\in L)$ une famille d'homomorphismes
d'anneaux telle que l'intersection des noyaux des $u_\lambda$ soit réduite à
$0$. Soit $M$ un $A$-module tel que pour tout $\lambda\in L$,
$M\otimes_A B_\lambda$ soit un $B_\lambda$-module libre et que
$M\otimes_A(A/\mathfrak J)$ soit un $(A/\mathfrak J)$-module libre. Alors $M$
est un $A$-module libre. Si de plus l'ensemble d'indices $L$ est fini, on peut
remplacer partout «~libre~» par «~plat~» dans l'énoncé précédent.
\end{theorem}

\begin{proof}
Dans les deux cas il suffira de prouver que $M$ est un $A$-module \emph{plat}~:
en effet, lorsque $M\otimes_A(A/\mathfrak J)$ est un
$(A/\mathfrak J)$-module libre, il en résultera que $M$ est un $A$-module libre
en vertu de $(\mathbf 0_{\mathrm{III}},10.1.2)$.

Nous utiliserons le lemme suivant, qui généralise
$(\mathbf 0_{\mathrm{III}},10.2.1)$~:

\begin{lemma}[11.4.1.1]
\label{IV.11.4.1.1-fr}
Soit $A$ un anneau muni d'une filtration finie
$(\mathfrak J_n)_{0\leq n\leq N+1}$, avec $\mathfrak J_0=A$,
$\mathfrak J_{N+1}=0$. Soit $M$ un $A$-module muni de la filtration
$(\mathfrak J_nM)_{0\leq n\leq N+1}$, et désignons par
$\operatorname{gr}(A)$ et $\operatorname{gr}(M)$ l'anneau et le module gradués
correspondants. Supposons que $M\otimes_A(A/\mathfrak J_1)$ soit un
$(A/\mathfrak J_1)$-module plat et que l'homomorphisme canonique
\begin{equation}
\operatorname{gr}_0(M)\otimes_{\operatorname{gr}_0(A)}\operatorname{gr}(A)
\longrightarrow\operatorname{gr}(M)
\label{IV.11.4.1.2-fr}
\tag{11.4.1.2}
\end{equation}
soit injectif. Alors $M$ est un $A$-module plat.
\end{lemma}

L'homomorphisme canonique (11.4.1.2) se définit de la même façon que celui de
$(\mathbf 0_{\mathrm{III}},10.1.1)$ comme étant en degré $n$ l'homomorphisme
\[
(M/\mathfrak J_1M)\otimes(\mathfrak J_n/\mathfrak J_{n+1})
=(M\otimes\mathfrak J_n)/
(\operatorname{Im}(M\otimes\mathfrak J_{n+1})
+\operatorname{Im}(\mathfrak J_1M\otimes\mathfrak J_n))
\longrightarrow\mathfrak J_nM/\mathfrak J_{n+1}M
\]
provenant de l'homomorphisme canonique $M\otimes\mathfrak J_n\to\mathfrak J_nM$
par passage aux quotients. Le lemme se démontre par récurrence sur $N$,
puisqu'il n'y a rien à démontrer pour $N=0$.

\oldpage[IV-3]{144}

Les conditions sur $M$ entraînent, en vertu de l'hypothèse de récurrence, que
$M\otimes_A(A/\mathfrak J_N)$ est un $(A/\mathfrak J_N)$-module plat. Notons
maintenant que l'on a $\mathfrak J_N^2\subset\mathfrak J_{N+1}=0$ si $N\geq1$,
et
\[
(M/\mathfrak J_1M)\otimes(\mathfrak J_N/\mathfrak J_{N+1})
=(M/\mathfrak J_NM)\otimes(\mathfrak J_N/\mathfrak J_{N+1})
=(M/\mathfrak J_NM)\otimes\mathfrak J_N\,;
\]
donc l'homomorphisme canonique
\[
(M/\mathfrak J_NM)\otimes\mathfrak J_N
\longrightarrow\mathfrak J_NM=\mathfrak J_NM/\mathfrak J_N^2M
\]
est injectif. Appliquant $(\mathbf 0_{\mathrm{III}},10.2.1)$ à la filtration
$\mathfrak J_N$-préadique, on en conclut bien que $M$ est un $A$-module plat.

Pour appliquer ce lemme à (11.4.1), nous désignerons par $\mathfrak J_n$
l'idéal de $A$ intersection des images réciproques
$u_\lambda^{-1}(\mathfrak J^nB_\lambda)$~; il est immédiat que l'on a
$\mathfrak J_0=A$, $\mathfrak J_m\mathfrak J_n\subset\mathfrak J_{m+n}$ pour
$m\geq0$, $n\geq0$~; en outre, si $\mathfrak J^{N+1}=0$, on a aussi
$\mathfrak J_{N+1}=0$ puisque l'intersection des noyaux des $u_\lambda$ est
réduite à $0$.

Munissons $A$ de la filtration $(\mathfrak J_n)$, $M$ de la filtration
$(\mathfrak J_nM)$, et, pour chaque $\lambda$, munissons $B_\lambda$ et
$N_\lambda=M\otimes_A B_\lambda$ des filtrations $\mathfrak J$-préadiques~;
considérons pour chaque $\lambda$ le diagramme commutatif
\[
\xymatrix{
\operatorname{gr}^0(M)\otimes_{\operatorname{gr}^0(A)}\operatorname{gr}(A)
 \ar[r]^{\varphi} \ar[d]_{f_\lambda} & \operatorname{gr}(M) \ar[d]^{g_\lambda} \\
\operatorname{gr}_{\mathfrak J}^0(N_\lambda)
 \otimes_{\operatorname{gr}_{\mathfrak J}^0(B_\lambda)}
 \operatorname{gr}_{\mathfrak J}(B_\lambda)
 \ar[r]_{\varphi_\lambda} & \operatorname{gr}_{\mathfrak J}(N_\lambda)
}
\]
où les flèches horizontales sont les homomorphismes canoniques (11.4.1.2) et
les flèches verticales sont déduites des homomorphismes canoniques
$A\to B_\lambda$ et $M\to N_\lambda$. L'hypothèse que $N_\lambda$ est un
$B_\lambda$-module plat entraîne que $\varphi_\lambda$ est injective
$(\mathbf 0_{\mathrm{III}},10.2.1)$, donc
$\operatorname{Ker}(\varphi)\subset\operatorname{Ker}(f_\lambda)$. Posant
$A_0=A/\mathfrak J_1$, $M_0=M\otimes_A A_0$, notons que
\[
\operatorname{gr}_{\mathfrak J}^0(N_\lambda)
=N_\lambda/\mathfrak JN_\lambda
=(M/\mathfrak JM)\otimes_A B_\lambda
=(M/\mathfrak J_1M)\otimes_A(B_\lambda/\mathfrak JB_\lambda),
\]
car ce dernier produit tensoriel est égal à
\[
(M\otimes_A B_\lambda)\Big/
\big(\operatorname{Im}(M\otimes_A\mathfrak JB_\lambda)
 +\operatorname{Im}(\mathfrak J_1M\otimes_A B_\lambda)\big),
\]
et l'on a
$\operatorname{Im}(\mathfrak J_1M\otimes_A B_\lambda)
=\operatorname{Im}(M\otimes_A\mathfrak J_1B_\lambda)
\subset\operatorname{Im}(M\otimes_A\mathfrak JB_\lambda)$ puisque
$\mathfrak J_1B_\lambda\subset\mathfrak JB_\lambda$ par définition~; enfin,
la relation $\mathfrak J_1B_\lambda\subset\mathfrak JB_\lambda$ montre que l'on
a aussi
\[
(M/\mathfrak J_1M)\otimes_A(B_\lambda/\mathfrak JB_\lambda)
=(M/\mathfrak J_1M)\otimes_{A/\mathfrak J_1}
 (B_\lambda/\mathfrak JB_\lambda),
\]
si bien que finalement on a
\[
\operatorname{gr}_{\mathfrak J}^0(N_\lambda)
=M_0\otimes_{A_0}(B_\lambda/\mathfrak JB_\lambda)
\]
et par suite
\[
\operatorname{gr}_{\mathfrak J}^0(N_\lambda)
\otimes_{\operatorname{gr}_{\mathfrak J}^0(B_\lambda)}
\operatorname{gr}_{\mathfrak J}(B_\lambda)
=M_0\otimes_{A_0}\operatorname{gr}_{\mathfrak J}(B_\lambda).
\]
L'homomorphisme $f_\lambda$ peut donc s'écrire
\[
1\otimes\operatorname{gr}(u_\lambda):
M_0\otimes_{A_0}\operatorname{gr}(A)
\longrightarrow M_0\otimes_{A_0}\operatorname{gr}_{\mathfrak J}(B_\lambda),
\]

\oldpage[IV-3]{145}

et comme $M_0$ est par hypothèse un $A_0$-module \emph{plat}, le noyau de
$f_\lambda$ est égal à $M_0\otimes_{A_0}R_\lambda$, où $R_\lambda$ est le
noyau de
$\operatorname{gr}(u_\lambda):\operatorname{gr}(A)
\to\operatorname{gr}_{\mathfrak J}(B_\lambda)$. Tout revient donc à prouver
que
\[
\bigcap_{\lambda\in L}(M_0\otimes_{A_0}R_\lambda)=0.
\]
Or, par définition des $\mathfrak J_n$, l'intersection des noyaux des
homomorphismes
\[
\mathfrak J_nA/\mathfrak J_{n+1}A
\longrightarrow\mathfrak J^nB_\lambda/\mathfrak J^{n+1}B_\lambda,
\]
lorsque $\lambda$ parcourt $L$, est réduite à $0$, autrement dit
$\bigcap_{\lambda\in L}R_\lambda=0$. Cela étant, supposons d'abord que $M_0$
soit un $A_0$-module libre~; prenant une base de $M_0$, on voit aussitôt que
l'on a
\[
M_0\otimes_{A_0}\Big(\bigcap_{\lambda\in L}R_\lambda\Big)
=\bigcap_{\lambda\in L}(M_0\otimes_{A_0}R_\lambda),
\]
d'où la proposition dans ce cas. Lorsque $L$ est fini, la formule précédente
est encore vraie sous la seule hypothèse que $M_0$ est un $A_0$-module plat
$(\mathbf 0_{\mathrm I},6.1.3)$, ce qui termine la démonstration.
\end{proof}

\begin{remark}[11.4.2]
\label{IV.11.4.2-fr}
La conclusion de (11.4.1) peut être inexacte si $L$ est infini et si l'on
suppose seulement que $M\otimes_A(A/\mathfrak J)$ est un
$(A/\mathfrak J)$-module plat. Par exemple, soient $V$ un anneau de valuation
discrète, $K$ son corps des fractions, et soit $A=V[T]/(T^2)$ ($T$
indéterminée)~; prenons pour $\mathfrak J$ l'image de $(T)$ dans $A$, de sorte
que $A/\mathfrak J=V$, et prenons $M=K$, qui est un
$(A/\mathfrak J)$-module, donc égal à $M\otimes_A(A/\mathfrak J)$~; en outre
$M$ est un $(A/\mathfrak J)$-module plat, mais non un $A$-module plat, car
puisque $\mathfrak J$ est nilpotent, il résulterait de
$(\mathbf 0_{\mathrm{III}},10.1.2)$ que $M$ serait un $A$-module libre, ce qui
est absurde puisque $\mathfrak JK=0$. Considérons d'autre part l'idéal maximal
$\mathfrak m$ de l'anneau local noethérien $A$~; on a $\mathfrak mK=K$, donc
$M\otimes_A(A/\mathfrak m^n)=0$ quel que soit l'entier $n$~; les
$(A/\mathfrak m^n)$-modules $M\otimes_A(A/\mathfrak m^n)$ sont donc plats pour
tout $n$, et l'intersection des $\mathfrak m^n$ est réduite à $0$.
\end{remark}

\begin{corollary}[11.4.3]
\label{IV.11.4.3-fr}
Soient $A$ un anneau semi-local dont le radical $\mathfrak J$ est nilpotent
(par exemple un anneau artinien), $u_\lambda:A\to B_\lambda$ $(\lambda\in L)$
une famille d'homomorphismes d'anneaux telle que l'intersection des noyaux des
$u_\lambda$ soit réduite à $0$. Pour qu'un $A$-module $M$ soit plat, il faut et
il suffit que pour tout $\lambda\in L$, $M\otimes_A B_\lambda$ soit un
$B_\lambda$-module plat.
\end{corollary}

\begin{proof}
Comme $A/\mathfrak J$ est composé direct d'un nombre fini de corps (Bourbaki,
\emph{Alg. comm.}, chap.~II, \S~3, n$^\circ$~5, prop.~16) et que $\mathfrak J$
est nilpotent, $A$ est composé direct d'un nombre fini d'anneaux locaux $A_i$
dont le radical est nilpotent (\emph{loc. cit.}, \S~4, n$^\circ$~3, cor.~1 de
la prop.~15) et $M$ est par suite somme directe de $A_i$-modules $M_i$, chaque
$M_i$ étant annulé par les $A_j$ d'indice $j\neq i$~; pour que $M$ soit un
$A$-module plat, il faut et il suffit que chacun des $M_i$ soit un $A_i$-module
plat~; d'ailleurs l'intersection des noyaux des homomorphismes
\[
A_i\longrightarrow A\xrightarrow{u_\lambda}B_\lambda
\]
est réduite à $0$, et $M_i\otimes_{A_i}B_\lambda$ est facteur direct de
$M\otimes_A B_\lambda$. On peut donc se borner au cas où $A$ est en outre
local. Alors $A/\mathfrak J$ est un corps, donc
$M\otimes_A(A/\mathfrak J)$ est un $(A/\mathfrak J)$-module libre, et il
suffit de voir que pour tout $\lambda$, $M\otimes_A B_\lambda$ est un
$B_\lambda$-module libre, en vertu de (11.4.1). Mais si l'on pose
$\mathfrak J_\lambda=\mathfrak JB_\lambda$,
\[
(M\otimes_A B_\lambda)\otimes_{B_\lambda}(B_\lambda/\mathfrak J_\lambda)
=(M\otimes_A(A/\mathfrak J))\otimes_{A/\mathfrak J}
 (B_\lambda/\mathfrak J_\lambda)
\]
est un $(B_\lambda/\mathfrak J_\lambda)$-module libre, et
$\mathfrak J_\lambda$ est nilpotent. La conclusion résulte donc de l'hypothèse
que $M\otimes_A B_\lambda$ est un $B_\lambda$-module plat et de
$(\mathbf 0_{\mathrm{III}},10.1.2)$.
\end{proof}

\begin{corollary}[11.4.4]
\label{IV.11.4.4-fr}
Soient $A$ un anneau, $M$ un $A$-module~; on suppose qu'il existe un idéal
nilpotent $\mathfrak J$ de $A$ tel que $M\otimes_A(A/\mathfrak J)$ soit un
$(A/\mathfrak J)$-module libre. Alors l'ensemble des idéaux $\mathfrak K$ de
$A$ tels que $M\otimes_A(A/\mathfrak K)$ soit un $(A/\mathfrak K)$-module
libre admet un plus petit élément $\mathfrak K_0$ (qui est aussi le plus petit
des idéaux $\mathfrak K$ tels que $M\otimes_A(A/\mathfrak K)$ soit un
$(A/\mathfrak K)$-module plat).
\end{corollary}

\oldpage[IV-3]{146}

\begin{proof}
Pour qu'un homomorphisme $u:A\to A'$ soit tel que $M\otimes_A A'$ soit un
$A'$-module libre (resp. un $A'$-module plat), il faut et il suffit que $u$ se
factorise en $A\to A/\mathfrak K_0\to A'$ (ou encore que
$\mathfrak K_0A'=0$).

Le fait que l'intersection $\mathfrak K_0$ des idéaux $\mathfrak K$ pour
lesquels $M\otimes_A(A/\mathfrak K)$ est un $(A/\mathfrak K)$-module libre
soit le plus petit de ces idéaux résulte de (11.4.1) appliqué à l'anneau
$A/\mathfrak K_0$, à son idéal nilpotent $\mathfrak J/\mathfrak K_0$ et aux
homomorphismes $A/\mathfrak K_0\to A/\mathfrak K$, dont les noyaux ont $0$
pour intersection. Si $A'$ est une $(A/\mathfrak K_0)$-algèbre, on a
\[
M\otimes_A A'=
(M\otimes_A(A/\mathfrak K_0))\otimes_{A/\mathfrak K_0}A',
\]
donc $M\otimes_A A'$ est un $A'$-module libre. Réciproquement, si $A'$ est
une $A$-algèbre telle que $M\otimes_A A'$ soit un $A'$-module libre et si
$\mathfrak K$ est le noyau de l'homomorphisme $A\to A'$, il résulte de
(11.4.1) appliqué à l'anneau $A/\mathfrak K$, au
$(A/\mathfrak K)$-module $M\otimes_A(A/\mathfrak K)$, à l'idéal nilpotent
$(\mathfrak J+\mathfrak K)/\mathfrak K$ et à l'homomorphisme injectif
$A/\mathfrak K\to A'$, que $M\otimes_A(A/\mathfrak K)$ est un
$(A/\mathfrak K)$-module libre, donc que $\mathfrak K\supset\mathfrak K_0$.
Le fait que l'on puisse remplacer «~libre~» par «~plat~» dans ce qui précède
(en conservant naturellement l'hypothèse que $M\otimes_A(A/\mathfrak J)$ est
un $(A/\mathfrak J)$-module libre) résulte de
$(\mathbf 0_{\mathrm{III}},10.1.2)$, comme on l'a vu au début de la
démonstration de (11.4.1).
\end{proof}

\begin{proposition}[11.4.5]
\label{IV.11.4.5-fr}
Soient $Y$ un préschéma irréductible, $f:X\to Y$ un morphisme de présentation
finie, $\mathcal F$ un $\mathcal O_X$-Module de présentation finie. Il existe
alors un ouvert non vide $U$ dans $Y$ et un sous-schéma fermé $Z$ de $U$, de
présentation finie sur $U$, ayant la propriété suivante~: pour tout changement
de base $U'\to U$, si l'on pose $X'=X\times_YU'$, $f'=f_{(U')}$ et
$\mathcal F'=\mathcal F\otimes_{\mathcal O_Y}\mathcal O_{U'}$, alors, pour que
$\mathcal F'$ soit $f'$-plat, il faut et il suffit que le morphisme $U'\to U$
se factorise en $U'\to Z\to U$. Un tel schéma $Z$ a même espace sous-jacent que
$U$. Supposons de plus $Y$ affine, et soit $(W_i)$ un recouvrement fini de $X$
par des ouverts affines~; alors, on peut supposer $U$ choisi de sorte que, si
$U'=\operatorname{Spec}(A')$ est un schéma affine et $U'\to U$ un morphisme qui
se factorise en $U'\to Z\to U$, chacun des
$\Gamma(W_i\times_YU',\mathcal F')$ est un $A'$-module libre.
\end{proposition}

\begin{proof}
On peut évidemment se borner au cas où $Y=\operatorname{Spec}(A)$ est affine.
Utilisant (8.9.1), il y a un sous-anneau noethérien $A_1$ de $A$, un morphisme
de type fini $f_1:X_1\to Y_1=\operatorname{Spec}(A_1)$ et un
$\mathcal O_{X_1}$-Module cohérent $\mathcal F_1$ tels que $f=(f_1)_{(Y)}$ et
$\mathcal F=\mathcal F_1\otimes_{\mathcal O_{Y_1}}\mathcal O_Y$~; on peut en
outre supposer que les $W_i$ sont images réciproques d'ouverts affines de
$X_1$. Notons en outre que $Y_1$ est irréductible, le morphisme $Y\to Y_1$
étant dominant $(\mathbf I,1.2.7)$. Supposons la proposition démontrée pour
$Y_1$, $X_1$ et $\mathcal F_1$, et soient $U_1$ l'ouvert de $Y_1$ et $Z_1$ le
sous-schéma fermé de $U_1$ ayant les propriétés voulues, $U=U_1\times_{Y_1}Y$
et $Z=Z_1\times_{Y_1}Y$ leurs images réciproques. Alors, si $U'\to U$ est un
changement de base tel que
$\mathcal F'=\mathcal F\otimes_{\mathcal O_Y}\mathcal O_{U'}
=\mathcal F_1\otimes_{\mathcal O_{Y_1}}\mathcal O_{U'}$ soit $f'$-plat, le
morphisme $U'\to U_1$ se factorise en $U'\to Z_1\to U_1$~; mais comme
$U'\to U_1$ se factorise aussi en $U'\to U\to U_1$, la définition du produit
de préschémas montre que $U'\to U$ se factorise en $U'\to Z\to U$.

On peut donc se borner au cas où $A$ est \emph{noethérien}~; soit
$\mathfrak N$ son nilradical, qui est nilpotent, et posons
$A_0=A_{\mathrm{red}}=A/\mathfrak N$,
$Y_0=\operatorname{Spec}(A_0)=Y_{\mathrm{red}}$,
$X_0=X\times_YY_0$, $\mathcal F_0=\mathcal F\otimes_{\mathcal O_Y}
\mathcal O_{Y_0}$, $W_{i0}=W_i\times_YY_0$~; si
$M_i=\Gamma(W_i,\mathcal F)$, $M_{i0}=\Gamma(W_{i0},\mathcal F_0)$ est donc
égal à $M_i\otimes_A A_0$. Comme $A_0$ est \emph{intègre}, on peut, en vertu du
théorème de platitude générique (6.9.2), en remplaçant au besoin $Y$ par un
ouvert non vide de $Y$, supposer que les $M_{i0}$ soient des $A_0$-modules
libres. En vertu de (11.4.4), il y a donc pour chaque $i$ un idéal
$\mathfrak J_i\subset\mathfrak N$ tel

\oldpage[IV-3]{147}

que les $A$-algèbres $A'$ pour lesquelles les $M_i\otimes_A A'$ sont des
$A'$-modules libres (ou plats) sont exactement celles pour lesquelles
$\mathfrak J_iA'=0$. Soit $\mathfrak J=\sum_i\mathfrak J_i$, qui est un idéal
contenu dans $\mathfrak N$~; dire que $\mathcal F\otimes_A A'$ est $A'$-plat
équivaut à dire que les $M_i\otimes_A A'$ sont tous des $A'$-modules plats,
donc que $\mathfrak JA'=0$. Il en résulte que si l'on prend
$Z=V(\mathfrak J)$, on répond à la question, car pour qu'un morphisme
$U'\to U$ ait la propriété de l'énoncé, il faut et il suffit évidemment que
pour tout ouvert affine $V'\subset U'$, le morphisme $V'\to U$ ait la même
propriété.
\end{proof}

\begin{corollary}[11.4.6]
\label{IV.11.4.6-fr}
Soient $A$ un anneau tel que $\operatorname{Spec}(A)$ soit irréductible,
$\mathfrak p$ son unique idéal premier minimal, $B$ une $A$-algèbre de
présentation finie, $M$ un $B$-module de présentation finie. Supposons que
$M_\mathfrak p$ soit un $A_\mathfrak p$-module plat~; alors il existe
$t\in A-\mathfrak p$ tel que $M_t$ soit un $A_t$-module libre.
\end{corollary}

\begin{proof}
Appliquant (11.4.5) à $Y=\operatorname{Spec}(A)$, $X=\operatorname{Spec}(B)$,
$\mathcal F=\widetilde M$, on peut (en remplaçant au besoin $A$ par $A_h$, où
$h\in A-\mathfrak p$) supposer qu'il existe un idéal de type fini
$\mathfrak J$ dans $A$ tel que les $A$-algèbres $A'$ pour lesquelles
$M\otimes_A A'$ soit un $A'$-module libre (ou plat) sont exactement celles
telles que $\mathfrak JA'=0$. En particulier, l'hypothèse que
$M_\mathfrak p$ soit un $A_\mathfrak p$-module plat entraîne
$\mathfrak JA_\mathfrak p=0$, ou encore $\mathfrak J_\mathfrak p=0$. Mais
comme $\mathfrak J$ est de type fini, il existe $t\in A-\mathfrak p$ tel que
$\mathfrak J_t=0$, ou encore $\mathfrak JA_t=0$, et par suite $M_t$ est un
$A_t$-module libre.
\end{proof}

\begin{proposition}[11.4.7]
\label{IV.11.4.7-fr}
Soient $A$ un anneau noethérien, $\mathfrak J$ un idéal de $A$, $M$ un
$A$-module idéalement séparé $(\mathbf 0_{\mathrm{III}},10.2.1)$ pour la
topologie $\mathfrak J$-préadique. Soit $(\mathfrak p_i)_{1\leq i\leq r}$ une
famille finie d'idéaux premiers de $A$ contenant $\mathfrak J$~; pour tout
entier $n\geq0$, soit $\mathfrak p_i^{(n)}$ la puissance symbolique $n$-ième de
$\mathfrak p_i$ (noyau de l'homomorphisme
$A\to A_{\mathfrak p_i}/\mathfrak p_i^nA_{\mathfrak p_i}$)~; posons
$\mathfrak J_n=\bigcap_{i=1}^{r}\mathfrak p_i^{(n)}$, de sorte que
$\mathfrak J_n\supset\mathfrak J^n$, et supposons que la topologie définie par
la filtration $(\mathfrak J_n)$ soit identique à la topologie
$\mathfrak J$-préadique (autrement dit que, pour tout $n$, il existe $m$ tel
que $\mathfrak J^n\supset\mathfrak J_m$). Pour que $M$ soit un $A$-module
plat, il faut et il suffit que $M\otimes_A(A/\mathfrak J)$ soit un
$(A/\mathfrak J)$-module plat et que, pour tout $i$,
$M\otimes_AA_{\mathfrak p_i}=M_{\mathfrak p_i}$ soit un
$A_{\mathfrak p_i}$-module plat.
\end{proposition}

\begin{proof}
Comme $M$ est idéalement séparé, il suffit, en vertu de
$(\mathbf 0_{\mathrm{III}},10.2.1)$, de montrer que, pour tout $n\geq0$,
$M\otimes_A(A/\mathfrak J^n)$ est un $(A/\mathfrak J^n)$-module plat~;
puisque tout $\mathfrak J^n$ contient un $\mathfrak J_m$, il revient au même
de prouver que pour tout $n$, $M\otimes_A(A/\mathfrak J_n)$ est un
$(A/\mathfrak J_n)$-module plat. Or, comme
$\mathfrak J_1\supset\mathfrak J$, $M\otimes_A(A/\mathfrak J_1)$ est un
$(A/\mathfrak J_1)$-module plat~; dans l'anneau $A/\mathfrak J_n$, l'idéal
$\mathfrak J_1/\mathfrak J_n$ est nilpotent et enfin l'intersection des noyaux
des homomorphismes
\[
A/\mathfrak J_n\longrightarrow
A_{\mathfrak p_i}/\mathfrak p_i^nA_{\mathfrak p_i}
\]
est nulle dans $A/\mathfrak J_n$, par définition de $\mathfrak J_n$. Il suffit
donc, par (11.4.1), de vérifier que
\[
M\otimes_A(A_{\mathfrak p_i}/\mathfrak p_i^nA_{\mathfrak p_i})
\]
est un $(A_{\mathfrak p_i}/\mathfrak p_i^nA_{\mathfrak p_i})$-module plat, ce
qui résulte de l'hypothèse que $M\otimes_AA_{\mathfrak p_i}$ est un
$A_{\mathfrak p_i}$-module plat.
\end{proof}

\begin{remark}[11.4.8]
\label{IV.11.4.8-fr}
L'hypothèse faite dans (11.4.7) sur la topologie définie par les
$\mathfrak J_n$ est vérifiée si, pour tout $n$ assez grand,
$\operatorname{Ass}(A/\mathfrak J^n)$ est contenu dans l'ensemble des
$\mathfrak p_i$. En effet, $\mathfrak J^n$ est alors intersection d'idéaux
primaires pour les $\mathfrak p_i$, dont chacun contient une puissance
symbolique de $\mathfrak p_i$, d'où la conclusion. En particulier~:
\end{remark}

\begin{corollary}[11.4.9]
\label{IV.11.4.9-fr}
Soient $A$ un anneau noethérien, $\mathfrak J$ un idéal nilpotent de $A$, $M$
un $A$-module. Pour que $M$ soit un $A$-module plat (resp. libre), il faut et
il suffit que $M\otimes_A(A/\mathfrak J)$ soit un $(A/\mathfrak J)$-module
plat (resp. libre) et que pour tout idéal premier
$\mathfrak p\in\operatorname{Ass}(A)$, $M_\mathfrak p$ soit un
$A_\mathfrak p$-module plat.
\end{corollary}
% IV3_B22_APPEND_ANCHOR

\oldpage[IV-3]{148}

L'assertion relative au cas où $M\otimes_A(A/\mathfrak J)$ est libre se déduit
encore de l'assertion relative au cas où $M\otimes_A(A/\mathfrak J)$ est plat
par $(\mathbf 0_{\mathrm{III}},10.1.2)$.

\begin{corollary}[11.4.10]
\label{IV.11.4.10-fr}
Soient $A$ un anneau noethérien, $\mathfrak J$ un idéal de $A$, $M$ un
$A$-module. On suppose que $M$ est idéalement séparé pour la topologie
$\mathfrak J$-préadique et que $\operatorname{gr}_{\mathfrak J}(A)$ est un
$(A/\mathfrak J)$-module plat. Pour que $M$ soit un $A$-module plat, il faut et
il suffit que $M\otimes_A(A/\mathfrak J)$ soit un $(A/\mathfrak J)$-module plat
et que, pour tout $\mathfrak p\in\operatorname{Ass}(A/\mathfrak J)$,
$M_{\mathfrak p}$ soit un $A_{\mathfrak p}$-module plat.
\end{corollary}

\begin{proof}
Compte tenu de (11.4.8), tout revient à montrer que
$\operatorname{Ass}(A/\mathfrak J^n)$ est contenu dans
$\operatorname{Ass}(A/\mathfrak J)$ pour tout $n$. Or, si $a\in A$ n'appartient
à aucun des $\mathfrak p\in\operatorname{Ass}(A/\mathfrak J)$, l'homothétie de
rapport $a$ est injective dans $A/\mathfrak J$~; comme chacun des
$\mathfrak J^k/\mathfrak J^{k+1}$ est un $(A/\mathfrak J)$-module plat, $a$ est
aussi un élément $(\mathfrak J^k/\mathfrak J^{k+1})$-régulier, donc $a$ est
$(A/\mathfrak J^n)$-régulier pour tout $n$ (Bourbaki, \emph{Alg. comm.},
chap.~III, \S~2, n$^\circ$~8, cor.~1 du th.~1), et par suite n'appartient à
aucun idéal premier associé à $A/\mathfrak J^n$, d'où le corollaire (Bourbaki,
\emph{Alg. comm.}, chap.~II, \S~1, n$^\circ$~1, prop.~2).
\end{proof}

\begin{proposition}[11.4.11]
\label{IV.11.4.11-fr}
Soient $A$ un anneau local artinien, d'idéal maximal $\mathfrak m$,
$Y=\operatorname{Spec}(A)$, $y$ l'unique point de $Y$, $f:X\to Y$ un morphisme
localement de type fini, $\mathcal F$ un $\mathcal O_X$-Module cohérent. Soient
$(A'_\alpha)$ une famille d'anneaux locaux, et pour chaque $\alpha$,
$u_\alpha:A\to A'_\alpha$ un homomorphisme d'anneaux (nécessairement local).
Posons $Y'_\alpha=\operatorname{Spec}(A'_\alpha)$,
$X'_\alpha=X\times_Y Y'_\alpha$,
$\mathcal F'_\alpha=\mathcal F\otimes_A A'_\alpha$. Soit $x$ un point de $X$,
et supposons vérifiées les conditions suivantes~:
\begin{enumerate}[label=(\roman*)]
  \item L'intersection des noyaux des $u_\alpha$ est réduite à $0$.
  \item L'extension $k(x)$ du corps résiduel $k$ de $A$ est primaire (4.3.1)
  (condition automatiquement vérifiée si $k$ est séparablement clos).
  \item Pour tout $\alpha$, il existe un point $x'_\alpha\in X'_\alpha$ dont
  les projections respectives dans $X$ et $Y'_\alpha$ sont $x$ et le point
  fermé $y'_\alpha$ de $Y'_\alpha$, et tel que $\mathcal F'_\alpha$ soit
  $Y'_\alpha$-plat au point $x'_\alpha$.
\end{enumerate}
Dans ces conditions, $\mathcal F$ est $f$-plat au point $x$.
\end{proposition}

\begin{proof}
La question étant locale sur $X$, on peut évidemment se borner au cas où $f$
est un morphisme de type fini et supposer $\mathcal F_x\neq0$. Nous procéderons
en plusieurs étapes.

\medskip
\noindent\textbf{I)} \emph{Réduction au cas où $A'_\alpha$ est un anneau local
d'un $Y$-préschéma de type fini et où le corps résiduel $k'_\alpha$ de
$A'_\alpha$ est une extension finie de $k$.}

\medskip
Comme la réduction se fait séparément sur chaque $A'_\alpha$, on peut supprimer
dans cette partie l'indice $\alpha$. Soit $\mathfrak m'$ l'idéal maximal de
$A'$. Considérons $A'$ comme limite inductive de ses sous-$A$-algèbres de type
fini $B_\lambda$, et posons $\mathfrak n_\lambda=\mathfrak m'\cap B_\lambda$~;
$A'$ est aussi limite inductive de ses sous-anneaux locaux
$B'_\lambda=(B_\lambda)_{\mathfrak n_\lambda}$ (5.13.3), et on a évidemment
$\mathfrak m'\cap B'_\lambda=\mathfrak n'_\lambda$, idéal maximal de
$B'_\lambda$. Il existe donc (11.2.6) un indice $\lambda$ tel que, en posant
$Z'_\lambda=\operatorname{Spec}(B'_\lambda)$,
$\mathcal F\otimes_A B'_\lambda$ soit $Z'_\lambda$-plat au point $x'_\lambda$,
projection de $x'$, la projection de $x'_\lambda$ sur $Z'_\lambda$ étant le
point fermé $z'_\lambda$ de $Z'_\lambda$.

On peut donc supposer qu'il existe un schéma affine $Z'$ de type fini sur $Y$
et un point $z'$ de $Z'$ tel que $A'=\mathcal O_{Z',z'}$, et que, si l'on pose
$T'=X\times_Y Z'$, il existe un point $t'\in T'$ dont les projections dans
$Z'$ et $X$ sont $z'$ et $x$, et tel que $\mathcal F\otimes_Y Z'$ soit
$Z'$-plat au point $t'$. Soient $Z'_1$ (resp. $X_1$) un sous-préschéma fermé de
$Z'$ (resp. $X$) ayant pour espace sous-jacent l'adhérence de $\{z'\}$ (resp.
$\{x\}$), et posons $T'_1=X_1\times_Y Z'_1$. L'ensemble $U$ des points de
$T'$ où $\mathcal F\otimes_Y Z'$ est $Z'$-plat est ouvert dans $T'$ (11.1.1) et
contient $t'$, donc

\oldpage[IV-3]{149}

$V=U\cap T'_1$ est ouvert non vide dans $T'_1$. L'anneau $A$, étant artinien,
est un anneau de Jacobson, donc (10.4.6) $Z'_1$ et $T'_1$ sont des préschémas de
Jacobson~; il existe par suite dans $V$ un point $t'_0$ fermé dans $V$, et son
image $z'_0$ dans $Z'_1$ est un point fermé de $Z'_1$ (10.4.7). Soient $f_1$ la
restriction de $f$ à $X_1$ et $f'_1:T'_1\to Z'_1$ la projection canonique~;
$V\cap f_1'^{-1}(z'_0)$ est un ouvert non vide dans
$f_1^{-1}(y)\otimes_{k(y)}k(z'_0)$, et comme ce dernier préschéma est plat sur
$f_1^{-1}(y)$, un point maximal $t'_1$ de $V\cap f_1'^{-1}(z'_0)$ est
nécessairement au-dessus de l'unique point maximal $x$ de $X_1$ (2.3.4).
Enfin, $k(z'_0)$ est une extension finie de $k$ (\textbf I, 6.4.2) et
l'homomorphisme $A=\mathcal O_{Y,y}\to A'=\mathcal O_{Z',z'}$ se factorise en
$\mathcal O_{Y,y}\to\mathcal O_{Z',z'_0}\to\mathcal O_{Z',z'}$, donc son noyau
est contenu dans celui de $\mathcal O_{Y,y}\to\mathcal O_{Z',z'_0}$. Ceci
achève la réduction annoncée.

\medskip
\noindent\textbf{II)} \emph{Réduction au cas où les $A'_\alpha$ sont en nombre
fini et sont des $A$-algèbres finies.} --- Soit $\mathfrak m'_\alpha$ l'idéal
maximal de $A'_\alpha$~; comme $A'_\alpha$ est un anneau local noethérien,
l'intersection des $\mathfrak m'^n_\alpha$ est $0$ (\textbf I, 7.3.5)~;
l'intersection des $u_\alpha^{-1}(\mathfrak m'^n_\alpha)$, pour tous les indices
$n$ et $\alpha$, est donc égale à l'intersection des noyaux des $u_\alpha$, donc
est réduite à $0$ par l'hypothèse (i). Puisque $A$ est artinien, il y a un nombre
fini de ces idéaux dont l'intersection est déjà $0$~; notons-les
$u_i^{-1}(\mathfrak m_i'^{\,n_i})$ $(1\leq i\leq r)$. Comme les $A'_i$ sont
noethériens, les $\mathfrak m_i'^{\,h}/\mathfrak m_i'^{\,h+1}$ sont des
$(A'_i/\mathfrak m'_i)$-modules de longueur finie, et comme
$A'_i/\mathfrak m'_i$ est un $k$-espace vectoriel de rang fini, on voit que
$A''_i=A'_i/\mathfrak m_i'^{\,n_i}$ est une $A$-algèbre finie et un anneau local
artinien. La réduction annoncée résulte donc de (2.1.4).

\medskip
\noindent\textbf{III)} \emph{Fin de la démonstration.} --- Supposons désormais
que les $A'_i$ $(1\leq i\leq r)$ soient en nombre fini et soient des
$A$-algèbres finies. Pour tout $i$, le corps résiduel $k_i$ de $A'_i$ est une
extension finie de $k$~; utilisant l'hypothèse (ii), on en conclut que l'image
réciproque de $x$ dans $X'_i$ est réduite au seul point $x'_i$ (4.3.2). Soient
$Y'$ le préschéma somme des $Y'_i$, $X'=X\times_Y Y'$, somme des $X'_i$,
$\mathcal F'=\mathcal F\otimes_Y Y'$. L'hypothèse entraîne que
$\mathcal F'$ est $Y'$-plat aux points de l'image réciproque de $x$ par la
projection $p:X'\to X$. Comme $p$ est de type fini, il existe par suite un
ouvert $U'\supset p^{-1}(x)$ tel que $\mathcal F'$ soit $Y'$-plat aux points de
$U'$ (11.1.1). En outre, le morphisme $Y'\to Y$ est fini puisque les $A'_i$
sont des $A$-algèbres finies~; donc $p$ est un morphisme fini (\textbf{II},
6.1.5), par suite fermé (\textbf{II}, 6.1.10), et il existe donc un voisinage
affine $U$ de $x$ dans $X$ tel que $p^{-1}(U)\subset U'$. Soient $B$ et $A'$
les anneaux des schémas $U$ et $Y'$ ($A'$ étant le composé direct des $A'_i$)~;
remplaçant $X$ par $U$, on a donc $\mathcal F=\widetilde M$, où $M$ est un
$B$-module et, par hypothèse, $M'=M\otimes_A A'$ est un $A'$-module plat
(2.1.2)~; comme l'homomorphisme $A\to A'$ est injectif par construction, on
peut appliquer (11.4.3), qui prouve que $M$ est un $A$-module plat, C.Q.F.D.
\end{proof}

\begin{proposition}[11.4.12]
\label{IV.11.4.12-fr}
Soient $A$ un anneau local artinien de corps résiduel $k$,
$Y=\operatorname{Spec}(A)$, $f:X\to Y$ un morphisme localement de type fini,
$\mathcal F$ un $\mathcal O_X$-Module cohérent. Soient $(A'_\alpha)$ une famille
d'anneaux locaux, et pour chaque $\alpha$, $u_\alpha:A\to A'_\alpha$ un
homomorphisme d'anneaux. Posons $Y'_\alpha=\operatorname{Spec}(A'_\alpha)$,
$X'_\alpha=X\times_Y Y'_\alpha$, $f'_\alpha=f_{(Y'_\alpha)}$,
$\mathcal F'_\alpha=\mathcal F\otimes_A A'_\alpha$. Soit $x$ un point de $X$,
et supposons vérifiées les conditions suivantes~:
\begin{enumerate}[label=(\roman*)]
  \item L'intersection des noyaux des $u_\alpha$ est réduite à $0$.
  \item Pour tout $\alpha$, $\mathcal F'_\alpha$ est $f'_\alpha$-plat en tous
  les points $x'_\alpha\in X'_\alpha$ dont les projections respectives dans
  $X$ et $Y'_\alpha$ sont $x$ et le point fermé $y'_\alpha$ de $Y'_\alpha$.
\end{enumerate}
Alors $\mathcal F$ est $f$-plat au point $x$.
\end{proposition}

\begin{proof}
% IV3_B22_P149_ANCHOR

\oldpage[IV-3]{150}

Par hypothèse, l'intersection des noyaux des $u_\alpha$ est réduite à $0$~;
comme $A$ est artinien, il y a déjà un nombre fini de ces noyaux dont
l'intersection est $0$, donc on peut se borner au cas où la famille
$(A'_\alpha)$ est finie. Soit $k'$ une clôture algébrique de $k$~; on sait
$(\mathbf 0_{\mathrm{III}},10.3.1)$ qu'il existe un homomorphisme local
$A\to B$ faisant de $B$ un $A$-module plat, tel que $B$ soit un anneau local
artinien, entier sur $A$ et que $B\otimes_A k$ soit isomorphe à $k'$.

Par platitude, les noyaux des homomorphismes
$B\to B'_\alpha=B\otimes_A A'_\alpha$ se déduisent de ceux des $u_\alpha$ par
tensorisation avec $B$, et comme ils sont en nombre fini, leur intersection est
réduite à $0$ $(\mathbf 0_{\mathrm I},6.1.3)$. Considérons les anneaux
$B'_{\alpha\beta}$, localisés de $B'_\alpha$ en ses idéaux maximaux~; on sait
(Bourbaki, \emph{Alg. comm.}, chap.~II, \S~3, n$^\circ$~3, cor.~2 du th.~1)
que l'intersection des noyaux des homomorphismes
$B'_\alpha\to B'_{\alpha\beta}$ (pour un $\alpha$ donné) est réduite à $0$~;
on en conclut que l'intersection des noyaux des homomorphismes composés
$v_{\alpha\beta}:B\to B'_\alpha\to B'_{\alpha\beta}$ ($\alpha$ et $\beta$
variables) est réduite à $0$. D'autre part, comme $B$ est entier sur $A$,
$B'_\alpha$ est entier sur $A'_\alpha$, donc ses idéaux maximaux sont au-dessus
de l'idéal maximal de $A'_\alpha$. Si l'on pose
$Z'_{\alpha\beta}=\operatorname{Spec}(B'_{\alpha\beta})$,
$T'_{\alpha\beta}=X'_\alpha\times_{Y'_\alpha}Z'_{\alpha\beta}$,
$f'_{\alpha\beta}=f_{(Z'_{\alpha\beta})}$,
$\mathcal G=\mathcal F\otimes_A B$ et
$\mathcal G'_{\alpha\beta}=\mathcal G\otimes_B B'_{\alpha\beta}$, on voit donc
que les hypothèses (i) et (ii) sont encore vérifiées lorsqu'on remplace
respectivement $A$, $A'_\alpha$, $u_\alpha$, $\mathcal F$ et $x$ par $B$,
$B'_{\alpha\beta}$, $v_{\alpha\beta}$, $\mathcal G$ et un point $t$ de
$T=X\otimes_A B$ au-dessus de $x$. Comme le corps résiduel de $B$ est
séparablement clos, on déduit de (11.4.11) que $\mathcal G$ est plat sur $B$ au
point $t$. Mais puisque $B$ est un $A$-module fidèlement plat, on conclut de
(2.1.4) que $\mathcal F$ est plat sur $A$ au point $x$, ce qui prouve (11.4.12).
\end{proof}

\begin{corollary}[11.4.13]
\label{IV.11.4.13-fr}
Soient $A$ un anneau artinien local, $Y=\operatorname{Spec}(A)$,
$f:X\to Y$ un morphisme localement de type fini, $\mathcal F$ un
$\mathcal O_X$-Module cohérent. Soit $(Y'_\alpha)$ une famille de
$Y$-préschémas, et pour tout $\alpha$, posons
$X'_\alpha=X\times_Y Y'_\alpha$, $f'_\alpha=f_{(Y'_\alpha)}$,
$\mathcal F'_\alpha=\mathcal F\otimes_Y Y'_\alpha$. Soit $x$ un point de $X$ et
supposons vérifiées les hypothèses suivantes~:
\begin{enumerate}[label=(\roman*)]
  \item L'intersection des noyaux des homomorphismes
  $A\to\Gamma(Y'_\alpha,\mathcal O_{Y'_\alpha})$ correspondant aux morphismes
  structuraux est réduite à $0$.
  \item Pour tout $\alpha$, $\mathcal F'_\alpha$ est $f'_\alpha$-plat en tous
  les points $x'_\alpha\in X'_\alpha$ dont la projection dans $X$ est $x$.
\end{enumerate}
Alors $\mathcal F$ est $f$-plat au point $x$.
\end{corollary}

\begin{proof}
En effet, pour tout $y'_\alpha\in Y'_\alpha$, considérons le schéma local
$Y''_{\alpha,y'_\alpha}=\operatorname{Spec}(\mathcal O_{Y'_\alpha,y'_\alpha})$~;
en vertu de (2.1.4),
$\mathcal F'_\alpha\otimes_{Y'_\alpha}Y''_{\alpha,y'_\alpha}$ est
$Y''_{\alpha,y'_\alpha}$-plat aux points dont les projections sur $X$ et
$Y''_{\alpha,y'_\alpha}$ sont $x$ et le point fermé de
$Y''_{\alpha,y'_\alpha}$. En outre le noyau de l'homomorphisme
$A\to\Gamma(Y'_\alpha,\mathcal O_{Y'_\alpha})$ est l'intersection des noyaux
des homomorphismes
$A\to\Gamma(Y''_{\alpha,y'_\alpha},\mathcal O_{Y''_{\alpha,y'_\alpha}})$, car
on se ramène aussitôt au cas où $Y'_\alpha$ est affine, et il suffit alors
d'appliquer Bourbaki, \emph{Alg. comm.}, chap.~II, \S~3, n$^\circ$~3, cor.~2 du
th.~1. En remplaçant la famille $(Y'_\alpha)$ par la famille des
$Y''_{\alpha,y'_\alpha}$, on est donc ramené à (11.4.12).
\end{proof}

\subsection{Descente de la platitude par des morphismes quelconques : cas général}
\label{subsection:IV.11.5-fr}

\begin{theorem}[11.5.1]
\label{IV.11.5.1-fr}
Soient $Y$ un préschéma localement noethérien, $f:X\to Y$ un morphisme
localement de type fini, $\mathcal F$ un $\mathcal O_X$-Module cohérent, $x$ un
point de $X$, $y=f(x)$. Soit $(Y'_\alpha)$ une famille de $Y$-préschémas locaux
$Y'_\alpha=\operatorname{Spec}(A'_\alpha)$ telle que les images des points
fermés $y'_\alpha$ des $Y'_\alpha$ soient toutes égales à $y$. Pour tout
$\alpha$, soient $\mathfrak m'_\alpha$ l'idéal maximal de $A'_\alpha$, et
$u_\alpha:\mathcal O_y\to A'_\alpha$ l'homomorphisme canonique
(\textbf I, 2.4.4)~; on suppose que les intersections finies des idéaux
$u_\alpha^{-1}(\mathfrak m_\alpha'^{\,n_\alpha})$ forment un système fondamental
de voisinages de $0$ dans $\mathcal O_y$. Posons
$X'_\alpha=X\times_Y Y'_\alpha$, $f'_\alpha=f_{(Y'_\alpha)}$,
$\mathcal F'_\alpha=\mathcal F\otimes_Y Y'_\alpha$, et supposons vérifiée l'une
des hypothèses suivantes~:
\begin{enumerate}[label=(\roman*)]
  \item Pour tout $\alpha$, $\mathcal F'_\alpha$ est $f'_\alpha$-plat en tous
  les points dont la projection dans $X$ est égale à $x$ et la projection dans
  $Y'_\alpha$ égale à $y'_\alpha$.
  \item Pour tout $\alpha$, il existe un $x'_\alpha\in X'_\alpha$ dont la
  projection sur $X$ est $x$ et la projection dans $Y'_\alpha$ égale à
  $y'_\alpha$, tel que $\mathcal F'_\alpha$ soit $f'_\alpha$-plat au point
  $x'_\alpha$, et $k(x)$ est une extension primaire de $k(y)$.
\end{enumerate}
Dans ces conditions, $\mathcal F$ est $f$-plat au point $x$.
\end{theorem}

\begin{proof}
Soit $\mathfrak m_y$ l'idéal maximal de $\mathcal O_y$~; comme $\mathcal O_y$
et $\mathcal O_x$ sont noethériens et $\mathcal O_y\to\mathcal O_x$ un
homomorphisme local, il suffit, en vertu de
$(\mathbf 0_{\mathrm{III}},10.2.2)$, de prouver que pour tout $n>0$,
$\mathcal F_x/\mathfrak m_y^n\mathcal F_x$ est un
$(\mathcal O_y/\mathfrak m_y^n)$-module plat. Désignons par
$(\mathfrak J_\lambda)$ la famille des intersections finies des
$u_\alpha^{-1}(\mathfrak m_\alpha'^{\,n_\alpha})$~; par hypothèse, il existe
$\lambda$ tel que $\mathfrak J_\lambda\subset\mathfrak m_y^n$, et comme
\[
\mathcal F_x\otimes_{\mathcal O_y}(\mathcal O_y/\mathfrak m_y^n)
=\bigl(\mathcal F_x\otimes_{\mathcal O_y}(\mathcal O_y/\mathfrak J_\lambda)\bigr)
 \otimes_{\mathcal O_y/\mathfrak J_\lambda}(\mathcal O_y/\mathfrak m_y^n),
\]
il suffira de prouver que $\mathcal F_x/\mathfrak J_\lambda\mathcal F_x$ est un
$(\mathcal O_y/\mathfrak J_\lambda)$-module plat. Or, pour chaque $\alpha$ tel
que $\mathfrak J_\lambda\subset u_\alpha^{-1}(\mathfrak m_\alpha'^{\,n_\alpha})$,
on a, par passage aux quotients, un homomorphisme
$u'_\alpha:\mathcal O_y/\mathfrak J_\lambda\to
A'_\alpha/\mathfrak m_\alpha'^{\,n_\alpha}$, et l'intersection des noyaux des
$u'_\alpha$ est réduite à $0$. Compte tenu de (\textbf I, 3.6.1), on voit qu'on
est ramené à (11.4.11) dans le cas (ii), et à (11.4.12) dans le cas (i).
\end{proof}

\begin{corollary}[11.5.2]
\label{IV.11.5.2-fr}
Soient $Y$ un préschéma localement noethérien, $f:X\to Y$ un morphisme
localement de type fini, $\mathcal F$ un $\mathcal O_X$-Module cohérent, $x$ un
point de $X$, $y=f(x)$~; posons $A=\mathcal O_y$. Soit $u:A\to A'$ un
homomorphisme de $A$ dans un anneau de Zariski $A'$, tel que l'image réciproque
$u^{-1}(\mathfrak r')$ du radical $\mathfrak r'$ de $A'$ soit l'idéal maximal
$\mathfrak m$ de $A$~; supposons en outre que l'homomorphisme
$\widehat u:\widehat A\to\widehat{A'}$ soit injectif. Posons
$Y'=\operatorname{Spec}(A')$, $X'=X\times_Y Y'$, $f'=f_{(Y')}$,
$\mathcal F'=\mathcal F\otimes_Y Y'$. Pour que $\mathcal F$ soit $f$-plat au
point $x$, il faut et il suffit que $\mathcal F'$ soit $f'$-plat en tout point
dont la projection dans $X$ est égale à $x$ et dont la projection dans $Y'$ est
égale à un point fermé de $Y'$.

Si en outre $A'$ est une $A$-algèbre finie, on peut dans ce qui précède
remplacer l'hypothèse que $\widehat u$ est injectif par l'hypothèse que $u$ est
injectif.
\end{corollary}

\begin{proof}
\oldpage[IV-3]{151}
% IV3_B22_P151_ANCHOR

Comme $A$ (resp. $A'$) s'identifie à un sous-anneau de $\widehat A$ (resp.
$\widehat{A'}$) (Bourbaki, \emph{Alg. comm.}, chap.~III, \S~3, n$^\circ$~3,
prop.~6), on voit d'abord que $u$ lui-même est injectif et que $\widehat u$ est
son prolongement par continuité à $\widehat A$.

Soit $(\mathfrak m'_\alpha)$ la famille des idéaux maximaux de $A'$~; comme on a
\[
\mathfrak m_\alpha'^{\,n}\widehat{A'}
=\widehat{\mathfrak m}_\alpha'^{\,n},
\qquad\text{et}\qquad
\widehat{\mathfrak m}_\alpha'^{\,n}\cap A'
=\mathfrak m_\alpha'^{\,n},
\]
et que les $\mathfrak m_\alpha'^{\,n}$ sont ouverts dans $A'$, on a
$u^{-1}(\mathfrak m_\alpha'^{\,n})
=A\cap\widehat u^{-1}(\widehat{\mathfrak m}_\alpha'^{\,n})$, et il suffira de
montrer que dans $\widehat A$, les intersections finies des
$\widehat u^{-1}(\widehat{\mathfrak m}_\alpha'^{\,n_\alpha})$ forment un système
fondamental de voisinages de $0$, ce qui permettra d'appliquer (11.5.1) aux
homomorphismes composés $v_\alpha\circ u:A\to A'_{\mathfrak m'_\alpha}$,
où $v_\alpha:A'\to A'_{\mathfrak m'_\alpha}$ est l'homomorphisme canonique.
Comme $\widehat A$ est complet, il suffira de montrer que l'intersection des
$\widehat u^{-1}(\widehat{\mathfrak m}_\alpha'^{\,n_\alpha})$ est réduite à
$0$ (Bourbaki, \emph{Alg. comm.}, chap.~III, \S~2, n$^\circ$~7, prop.~8, où on
peut dans la démonstration remplacer la suite décroissante par un ensemble
filtrant quelconque). Or, pour tout $\alpha$ fixé, l'intersection des idéaux
$\widehat{\mathfrak m}_\alpha'^{\,n}\widehat{A'}_{\widehat{\mathfrak m}'_\alpha}$
pour $n>0$ est réduite à $0$ dans l'anneau local noethérien
$\widehat{A'}_{\widehat{\mathfrak m}'_\alpha}$. D'autre part les
$\widehat{\mathfrak m}'_\alpha$ sont les idéaux maximaux de $\widehat{A'}$, donc
l'homomorphisme canonique
$\widehat{A'}\to\prod_\alpha\widehat{A'}_{\widehat{\mathfrak m}'_\alpha}$ est
injectif (Bourbaki, \emph{Alg. comm.}, chap.~II, \S~3, n$^\circ$~3, cor.~2 du
th.~1), et comme par hypothèse $\widehat u:\widehat A\to\widehat{A'}$ est aussi
injectif, cela achève la démonstration dans le cas général. La dernière
assertion résulte de ce que $\widehat A$ est un $A$-module fidèlement plat
($\mathbf 0_{\mathrm I},7.3.5$) et $\widehat{A'}=A'\otimes_A\widehat A$ puisque
$A'$ est par hypothèse un $A$-module de type fini (Bourbaki, \emph{Alg. comm.},
chap.~III, \S~3, n$^\circ$~4, th.~3 et chap.~IV, \S~2, n$^\circ$~5, cor.~3 de la
prop.~9).
\end{proof}

\oldpage[IV-3]{152}

\begin{proposition}[11.5.3]
\label{IV.11.5.3-fr}
Soient $f:X\to Y$ un morphisme localement de présentation finie,
$\mathcal F$ un $\mathcal O_X$-Module quasi-cohérent de présentation finie, $x$
un point de $X$, $y=f(x)$, $g=(\psi,\theta):Y'\to Y$ un morphisme propre de
présentation finie. Supposons que~:
\begin{enumerate}[label=(\roman*)]
  \item l'homomorphisme $\theta_y:\mathcal O_y\to
  (g_*\mathcal O_{Y'})_y$ soit injectif~;
  \item pour tout $x'\in X'=X\times_Y Y'$ dont la projection dans $X$ est $x$,
  $\mathcal F'=\mathcal F\otimes_Y Y'$ soit $Y'$-plat au point $x'$.
\end{enumerate}
Alors $\mathcal F$ est $f$-plat au point $x$.
\end{proposition}

\begin{proof}
La question étant locale sur $X$, on peut supposer $f$ de présentation finie.
Soit $p:X'\to X$ la projection canonique. Comme
$f'=f_{(Y')}:X'\to Y'$ est de présentation finie (1.6.2) et que $\mathcal F'$
est un $\mathcal O_{X'}$-Module de présentation finie
$(\mathbf 0_{\mathrm I},5.2.5)$, il résulte de (11.3.1) que l'ensemble $U'$ des
points de $X'$ où $\mathcal F'$ est $f'$-plat est ouvert. Comme $U'$ contient
$p^{-1}(x)$ par hypothèse, et que $p$ est propre, donc fermé, $U'$ contient un
ensemble de la forme $p^{-1}(U)$, où $U$ est un voisinage de $x$. Remplaçant $X$
par $U$, on peut donc supposer déjà que $\mathcal F'$ est $f'$-plat.

D'autre part, en tenant compte de (\textbf I, 3.6.5), (\textbf{II}, 5.4.2) et
(2.1.4), on peut remplacer $Y$ par $\operatorname{Spec}(\mathcal O_y)$,
autrement dit supposer que $Y=\operatorname{Spec}(A)$, où $A$ est un anneau
local. Sous ces conditions, on va prouver que $\mathcal F$ est $f$-plat. En
vertu de (5.13.3), $A$ est limite inductive de sous-anneaux locaux noethériens
$A_\lambda$ tels que l'injection canonique $A_\lambda\to A$ soit un
homomorphisme local. En vertu de (8.9.1), on peut supposer que
$X=X_\lambda\otimes_{A_\lambda}A$, $f=f_\lambda\otimes1_A$,
$\mathcal F=\mathcal F_\lambda\otimes_{A_\lambda}A$, pour un $\lambda$
convenable, $f_\lambda:X_\lambda\to Y_\lambda=\operatorname{Spec}(A_\lambda)$
étant un morphisme de type fini et $\mathcal F_\lambda$ un
$\mathcal O_{X_\lambda}$-Module cohérent. De même, on peut supposer que
$Y'=Y'_\lambda\otimes_{A_\lambda}A$, $g=g_\lambda\otimes1_A$, où
$g_\lambda:Y'_\lambda\to Y_\lambda$ est un morphisme de présentation finie~;
en outre, (8.10.5, (xii)) permet de supposer $g_\lambda$ propre. Comme par
hypothèse l'homomorphisme $A\to\Gamma(Y',\mathcal O_{Y'})$ est injectif et que
$A_\lambda$ est un sous-anneau de $A$, l'homomorphisme
$A_\lambda\to\Gamma(Y',\mathcal O_{Y'})$ est aussi injectif. Enfin, en vertu de
(11.2.6), on peut supposer $\lambda$ pris tel que
$\mathcal F'_\lambda=\mathcal F_\lambda\otimes_{Y_\lambda}Y'_\lambda$ soit
$f'_\lambda$-plat, puisque $\mathcal F'=\mathcal F'_\lambda\otimes_{Y'_\lambda}Y'$.
Ces remarques prouvent que l'on peut désormais supposer l'anneau $A$
noethérien, les autres hypothèses de (11.5.3) étant vérifiées.

Posons alors $B=\widehat A$, $Z=\operatorname{Spec}(B)$~; comme $B$ est un
$A$-module fidèlement plat $(\mathbf 0_{\mathrm I},7.3.5)$, il revient au même
de dire que $\mathcal F$ est $f$-plat ou que $\mathcal F\otimes_Y Z$ est
$Z$-plat (2.1.4). De même, si l'on pose $Z'=Y'\times_Y Z$, le morphisme
$Z'\to Y'$ est fidèlement plat (2.2.13), donc il revient au même de dire que
$\mathcal F'$ est $f'$-plat ou que $\mathcal F'\otimes_{Y'}Z'$ est $Z'$-plat~;
enfin $h=g_{(Z)}:Z'\to Z$ est propre (\textbf{II}, 5.4.2) et de type fini
(1.5.2), $\widehat A$ est noethérien, et si $z$ est son point fermé,
l'homomorphisme $\mathcal O_z\to(h_*\mathcal O_{Z'})_z$ est injectif, car il
résulte de (2.3.1) que $h_*\mathcal O_{Z'}=g_*\mathcal O_{Y'}\otimes_Y Z$,
et notre assertion résulte de l'hypothèse (i) et de la définition des modules
plats $(\mathbf 0_{\mathrm I},6.1.1)$.

\oldpage[IV-3]{153}

On peut donc désormais supposer l'anneau local noethérien $A$ complet~; la
démonstration sera achevée si l'on prouve que l'intersection des noyaux des
homomorphismes $A=\mathcal O_y\to\mathcal O_{y'}$ (où $y'$ parcourt
$g^{-1}(y)$) est réduite à $0$. En effet, les $\mathcal O_{y'}$ sont des anneaux
locaux noethériens, donc pour chaque $y'$ l'intersection des
$\mathfrak m_{y'}^n$ $(n\geq0)$ est réduite à $0$~; si $\mathfrak a_{n,y'}$ est
l'image réciproque dans $A$ de $\mathfrak m_{y'}^n$, les intersections finies
des $\mathfrak a_{n,y'}$ sont des voisinages de $0$ dans $A$ et l'intersection de
tous les $\mathfrak a_{n,y'}$ est réduite à $0$~; elles forment donc un système
fondamental de voisinages de $0$ dans $A$ (Bourbaki, \emph{Alg. comm.}, chap.~III,
\S~2, n$^\circ$~7, prop.~8, où l'on peut remplacer la suite décroissante par un
ensemble filtrant quelconque), et l'on pourra par suite appliquer (11.5.1).
Or, soit $s\in A$ un élément appartenant au noyau de chacun des homomorphismes
$A\to\mathcal O_{y'}$~; l'image $s'$ de $s$ dans
$\Gamma(Y',\mathcal O_{Y'})$ est donc une section de $\mathcal O_{Y'}$ au-dessus
de $Y'$ telle que $s'_{y'}=0$ pour tout $y'\in g^{-1}(y)$~; il y a par suite un
voisinage $V'$ de $g^{-1}(y)$ dans $Y'$ tel que $s'|V'=0$. Mais comme $g$ est
fermé, $V'$ contient un ensemble de la forme $g^{-1}(V)$, où $V$ est un
voisinage ouvert de $y$ dans $Y$~; or $Y$ est un schéma local, donc le seul
voisinage du point fermé $y$ de $Y$ est $Y$ tout entier, autrement dit $V'=Y'$,
$s'=0$, et comme $A\to\Gamma(Y',\mathcal O_{Y'})$ est injectif par hypothèse,
$s=0$. C.Q.F.D.
\end{proof}

Le cas particulier suivant de (11.5.3) nous sera utile au chap.~V.

\begin{corollary}[11.5.4]
\label{IV.11.5.4-fr}
Soit $f=(\psi,\theta):X\to Y$ un morphisme propre de présentation finie, et
soit $p:X\times_YX\to X$ la première projection. On suppose que
$\theta:\mathcal O_Y\to f_*(\mathcal O_X)$ est injectif. Alors, pour que $f$
soit plat, il faut et il suffit que $p$ le soit.
\end{corollary}

\begin{proposition}[11.5.5]
\label{IV.11.5.5-fr}
Soient $A$ un anneau, $Y=\operatorname{Spec}(A)$, $f:X\to Y$ un morphisme
localement de présentation finie, $\mathcal F$ un $\mathcal O_X$-Module
quasi-cohérent de présentation finie, $x$ un point de $X$. Soit
$A\to A'$ un homomorphisme injectif faisant de $A'$ une algèbre entière sur $A$.
Posons $Y'=\operatorname{Spec}(A')$, $X'=X\times_Y Y'$, $f'=f\times1$,
$\mathcal F'=\mathcal F\otimes_Y Y'$. Alors, si $\mathcal F'$ est $f'$-plat en
tout point de $X'$ dont la projection dans $X$ est égale à $x$, $\mathcal F$ est
$f$-plat au point $x$.
\end{proposition}

\begin{proof}
Supposons d'abord que $A'$ soit une $A$-algèbre finie et de présentation finie~;
alors le morphisme $Y'\to Y$ est propre (\textbf{II}, 6.1.11) et de présentation
finie, donc les hypothèses de (11.5.3) sont vérifiées, d'où la conclusion. Dans
le cas général, la proposition résultera de ce cas particulier, du fait que $A'$
est limite inductive de ses sous-$A$-algèbres finies $A'_\lambda$ et des deux
lemmes suivants~:
\end{proof}

\begin{lemma}[11.5.5.1]
\label{IV.11.5.5.1-fr}
Toute $A$-algèbre finie $A'$ est limite inductive de $A$-algèbres $A'_\lambda$
qui sont finies et de présentation finie.
\end{lemma}

\begin{proof}
On raisonne comme dans (1.9.3.1). On a en effet $A'=B/\mathfrak J$, où $B$ est
une $A$-algèbre finie qui est un $A$-module libre, et $\mathfrak J$ un idéal de
$B$ (1.4.7.1). Or, $\mathfrak J$ est limite inductive des idéaux
$\mathfrak J_\lambda\subset\mathfrak J$ de $B$ qui sont de type fini (et a
fortiori des $A$-modules de type fini), donc, par l'exactitude du foncteur
$\varinjlim$, $A'$ est limite inductive des $A$-algèbres
$B/\mathfrak J_\lambda$~; or $B/\mathfrak J_\lambda$ est par définition un
$A$-module de présentation finie, donc aussi (1.4.7) une $A$-algèbre de
présentation finie.
\end{proof}
% IV3_B22_P153_END

\oldpage[IV-3]{154}

Pour appliquer ce lemme à la situation de (11.5.5) on notera en outre que si
l'homomorphisme $A\to A'$ est injectif, il en est ainsi \emph{a fortiori} de
$A\to A'_{\lambda}$, pour tout $\lambda$.

\begin{lemma}[11.5.5.2]
\label{IV.11.5.5.2-fr}
Soient $A$ un anneau, $A'$ une $A$-algèbre, $(A'_{\lambda})$ un système inductif
de $A$-algèbres tel que $A'=\varinjlim A'_{\lambda}$~; on pose
$Y=\operatorname{Spec}(A)$, $Y'=\operatorname{Spec}(A')$,
$Y'_{\lambda}=\operatorname{Spec}(A'_{\lambda})$. Soient $f:X\to Y$ un
morphisme de présentation finie, $\mathcal F$ un $\mathcal O_X$-Module
quasi-cohérent de présentation finie~; on pose $X'=X\times_Y Y'$,
$f'=f_{(Y')}$, $\mathcal F'=\mathcal F\otimes_Y Y'$,
$X'_{\lambda}=X\times_Y Y'_{\lambda}$, $f'_{\lambda}=f_{(Y'_{\lambda})}$,
$\mathcal F'_{\lambda}=\mathcal F\otimes_Y Y'_{\lambda}$. Soit $x$ un point de
$X$, tel que $\mathcal F'$ soit $f'$-plat en tous les points $x'\in X'$
au-dessus de $x$~; alors il existe $\lambda$ tel que $\mathcal F'_{\lambda}$
soit $f'_{\lambda}$-plat en tous les points de $X'_{\lambda}$ au-dessus de $x$.
\end{lemma}

\begin{proof}
Soit $U'$ l'ensemble des $x'\in X'$ tels que $\mathcal F'$ soit $f'$-plat au
point $x'$~; on sait (11.3.1) que $U'$ est ouvert dans $X'$ puisque $f'$ est de
présentation finie (1.6.2)~; de même l'ensemble $U'_{\lambda}$ des points de
$X'_{\lambda}$ où $\mathcal F'_{\lambda}$ est $f'_{\lambda}$-plat est ouvert
dans $X'_{\lambda}$, et on sait en outre (11.2.6) que $U'$ est réunion des
$v_{\lambda}^{-1}(U'_{\lambda})$, où
$v_{\lambda}:X'\to X'_{\lambda}$ est la projection canonique. Considérons le
schéma $T=\operatorname{Spec}(k(x))$~; posons $T'=T\times_Y Y'$,
$T'_{\lambda}=T\times_Y Y'_{\lambda}$, et soient
$w_{\lambda}:T'\to T'_{\lambda}$, $g':T'\to X'$,
$g'_{\lambda}:T'_{\lambda}\to X'_{\lambda}$ les projections canoniques. Posons
$V'_{\lambda}=g_{\lambda}'^{-1}(U'_{\lambda})$,
$V'=g'^{-1}(U')=\bigcup_{\lambda}w_{\lambda}^{-1}(V'_{\lambda})$. Par
hypothèse on a (compte tenu de (\textbf I, 3.6.1)) $V'=T'$~; comme $T'$ est
quasi-compact, il existe $\lambda$ tel que
$w_{\lambda}^{-1}(V'_{\lambda})=T'$. On déduit alors de (8.3.3) appliqué aux
parties fermées quasi-compactes $T'_{\lambda}-V'_{\lambda}$ de
$T'_{\lambda}$, qu'il existe $\mu\geq\lambda$ tel que
$T'_{\mu}=V'_{\mu}$~; cela signifie que $\mathcal F'_{\mu}$ est
$f'_{\mu}$-plat en tous les points de $X'_{\mu}$ dont la projection dans $X$
est $x$. C.Q.F.D.
\end{proof}

\subsection{Descente de la platitude par des morphismes quelconques : cas d'un
préschéma de base unibranche}
\label{subsection:IV.11.6-fr}

\begin{theorem}[11.6.1]
\label{IV.11.6.1-fr}
Soient $A$ un anneau local intègre géométriquement unibranche
$(\mathbf 0,23.2.1)$, $Y=\operatorname{Spec}(A)$, $f:X\to Y$ un morphisme
localement de présentation finie, $\mathcal F$ un $\mathcal O_X$-Module
quasi-cohérent de présentation finie. Soient $A'$ un anneau local,
$u:A\to A'$ un homomorphisme local injectif~; on pose
$Y'=\operatorname{Spec}(A')$, $X'=X\times_Y Y'$, $f'=f_{(Y')}$,
$\mathcal F'=\mathcal F\otimes_Y Y'$. Soient $x$ un point de $X$ dont la
projection $f(x)=y$ est le point fermé de $Y$, $x'$ un point de $X'$ dont les
projections dans $X$ et $Y'$ sont respectivement $x$ et le point fermé $y'$ de
$Y'$. Alors, si $\mathcal F'$ est $f'$-plat au point $x'$, $\mathcal F$ est
$f$-plat au point $x$.
\end{theorem}

\begin{proof}
On peut se borner au cas où $f$ est de présentation finie, la question étant
locale sur $X$. Nous procéderons en plusieurs étapes.

\medskip
\noindent\textbf{I)} \emph{Réduction au cas où $A$ et $A'$ sont des anneaux
locaux intégralement clos.}

Comme $u$ est injectif et $A$ intègre, il existe un idéal premier
$\mathfrak p'$ de $A'$ tel que $u^{-1}(\mathfrak p')=0$
$(\mathbf 0_{\mathrm I},1.5.8)$~; l'homomorphisme composé
$A\xrightarrow{u}A'\to A''=A'/\mathfrak p'$ est donc injectif et local, et si
$Y''=\operatorname{Spec}(A'')$, $X''=X'\otimes_{A'}A''=X\times_Y Y''$,
$\mathcal F''=\mathcal F'\otimes_{A'}A''$, $\mathcal F''$ est $Y''$-plat aux
points de $X''$ au-dessus de $x'$ (2.1.4)~; remplaçant au besoin $A'$ par $A''$
et tenant compte de (\textbf I, 3.4.7), on peut donc supposer d'abord que $A'$
est intègre. Si $K'$ est le corps des fractions de $A'$, il y a alors un anneau
de valuation $B'$ dans $K'$ qui domine $A'$~; l'homomorphisme composé
$A\xrightarrow{u}A'\to B'$ étant injectif et local, le même raisonnement que
% IV3_B23_P154_END

\oldpage[IV-3]{155}

précédemment permet de remplacer $A'$ par $B'$~; on peut donc supposer l'anneau
local $A'$ intégralement clos, $A$ étant un sous-anneau local de $A'$ dominé par
$A'$. Soit $A_1$ la clôture intégrale de $A$~; il est clair que
$A\subset A_1\subset A'$, et par hypothèse $A_1$ est un anneau local~; si
$\mathfrak m$, $\mathfrak m_1$, $\mathfrak m'$ sont les idéaux maximaux de
$A$, $A_1$, $A'$, on a $\mathfrak m\subset\mathfrak m_1\subset\mathfrak m'$~;
en effet, $\mathfrak m_1$ est le \emph{seul} idéal premier de $A_1$ au-dessus de
$\mathfrak m$, puisque $A_1$ est un anneau local (Bourbaki,
\emph{Alg. comm.}, chap.~V, \S~2, n$^\circ$~1, prop.~1)~; comme
$\mathfrak m'\cap A=\mathfrak m$, on a
$\mathfrak m'\cap A_1\cap A=\mathfrak m$, donc
$\mathfrak m'\cap A_1=\mathfrak m_1$. Posons
$Y_1=\operatorname{Spec}(A_1)$, $X_1=X\times_Y Y_1$, $f_1=f_{(Y_1)}$,
$\mathcal F_1=\mathcal F\otimes_Y Y_1$, et soit $x_1$ la projection de $x'$ dans
$X_1$~; désignons d'autre part par $y_1$ l'unique point fermé de $Y_1$, de sorte
que $y_1=f_1(x_1)$. Par hypothèse le morphisme
$\operatorname{Spec}(k(y_1))\to\operatorname{Spec}(k(y))$ est radiciel, d'où
l'on conclut, par la transitivité des fibres (\textbf I, 3.6.4) et
(\textbf I, 3.5.7), que le morphisme $f_1^{-1}(y_1)\to f^{-1}(y)$ est radiciel,
et en particulier que $x_1$ est le \emph{seul} point de $X_1$ dont les
projections dans $X$ et $Y_1$ soient $x$ et $y_1$ respectivement~; en outre, on
a vu que $y_1$ est le seul point de $Y_1$ dont la projection dans $Y$ soit $y$,
donc $x_1$ est le seul point de $X_1$ dont la projection dans $X$ soit $x$. Si
l'on prouve que $\mathcal F_1$ est $f_1$-plat au point $x_1$, on pourra donc
appliquer (11.5.5), d'où résultera la conclusion. On est donc ramené au cas où
$A$ lui aussi est \emph{intégralement clos}.

\medskip
\noindent\textbf{II)} \emph{Réduction au cas où $A$ et $A'$ sont des anneaux
locaux de $\mathbf Z$-algèbres de type fini intégralement clos.}

On peut considérer $A'$ comme limite inductive filtrante de ses
sous-$\mathbf Z$-algèbres (intègres) de type fini $B'_{\lambda}$~; en outre,
comme $A'$ est intégralement clos et que la clôture intégrale d'une
$\mathbf Z$-algèbre de type fini est aussi une $\mathbf Z$-algèbre de type fini
(7.8.3), on voit que $A'$ est limite inductive de ses sous-$\mathbf Z$-algèbres
de type fini $B'_{\lambda}$ \emph{intégralement closes}~; si
$\mathfrak m'_{\lambda}=\mathfrak m'\cap B'_{\lambda}$, $A'$ est aussi limite
inductive des sous-anneaux locaux
$(B'_{\lambda})_{\mathfrak m'_{\lambda}}=A'_{\lambda}$ dominés par $A'$
(5.13.3). Pour tout $\lambda$, $B'_{\lambda}\cap A$ est aussi limite inductive
de ses sous-$\mathbf Z$-algèbres de type fini $B_{\alpha\lambda}$, donc
$A=\varinjlim_{\alpha,\lambda}B_{\alpha\lambda}$, et comme plus haut on peut
remplacer $B_{\alpha\lambda}$ dans cette formule par sa clôture intégrale
(contenue par hypothèse dans $B'_{\lambda}\cap A$), puis par l'anneau local
$A_{\alpha\lambda}=(B_{\alpha\lambda})_{\mathfrak m_{\alpha\lambda}}$, où
$\mathfrak m_{\alpha\lambda}=\mathfrak m\cap B_{\alpha\lambda}
=\mathfrak m'_{\lambda}\cap B_{\alpha\lambda}$, de sorte que
$A_{\alpha\lambda}$ est dominé par $A'_{\lambda}$ et par $A$. Posons
$Y_{\alpha\lambda}=\operatorname{Spec}(A_{\alpha\lambda})$~; il résulte de
(8.9.1) qu'il existe un couple $(\alpha,\lambda)$ assez grand, un morphisme
$f_{\alpha\lambda}:X_{\alpha\lambda}\to Y_{\alpha\lambda}$ de type fini et un
$\mathcal O_{X_{\alpha\lambda}}$-Module cohérent $\mathcal F_{\alpha\lambda}$
tels que
\[
X=X_{\alpha\lambda}\times_{Y_{\alpha\lambda}}Y,
\qquad f=f_{\alpha\lambda}\times1_Y,
\qquad \mathcal F=\mathcal F_{\alpha\lambda}\otimes_{Y_{\alpha\lambda}}Y~;
\]
si $x_{\alpha\lambda}$ est la projection de $x$ dans $X_{\alpha\lambda}$, il
suffira de montrer que $\mathcal F_{\alpha\lambda}$ est
$f_{\alpha\lambda}$-plat au point $x_{\alpha\lambda}$. Comme $A'$ est limite
inductive des $A'_{\mu}$ pour $\mu\geq\lambda$, $X'$ est limite projective des
$X_{\alpha\lambda}\times_{Y_{\alpha\lambda}}Y'_{\mu}=X'_{\mu}$, et on a aussi
$\mathcal F'=\mathcal F'_{\alpha\mu}\otimes_{Y'_{\mu}}Y'$, où
$\mathcal F'_{\alpha\mu}
=\mathcal F_{\alpha\lambda}\otimes_{Y_{\alpha\lambda}}Y'_{\mu}$.
Appliquant (11.2.6), on voit qu'on peut prendre $\mu$ assez grand pour que
$\mathcal F'_{\alpha\mu}$ soit $Y'_{\mu}$-plat au point $x'_{\mu}$, projection
de $x'$ dans $X'_{\mu}$, et en outre, par construction des $A'_{\mu}$, la
projection de $x'_{\mu}$ dans $Y'_{\mu}$ est le point fermé de $Y'_{\mu}$.

\medskip
\noindent\textbf{III)} \emph{Réduction au cas où le corps résiduel $k'$ de
$A'$ est une extension finie et radicielle du corps résiduel $k$ de $A$.}

On peut en premier lieu répéter le raisonnement de la partie I) de la
démonstration de (11.4.11), tenant compte du fait que $\mathbf Z$ est un anneau
de
% IV3_B23_P155_END

\oldpage[IV-3]{156}

Jacobson~; on se réduit ainsi au cas où $k'$ est une extension \emph{finie} de
$k$, ce que l'on va supposer dans ce qui suit. Soient $k''$ la plus grande
extension séparable de $k$ contenue dans $k'$, $k_1$ une extension galoisienne
finie de $k$ contenant $k''$, de sorte que $k''\otimes_k k_1$ est composée
directe de corps isomorphes à $k_1$~; comme $k'$ est une extension radicielle de
$k''$, $k'\otimes_k k_1$ est donc composée directe d'extensions radicielles de
$k_1$. Il existe un anneau local $A_1$ qui est une $A$-algèbre et un $A$-module
\emph{libre de type fini}, tel que $A_1/\mathfrak m A_1$ soit $k$-isomorphe à
$k_1$ $(\mathbf 0_{\mathrm{III}},10.3.1.2)$~; de façon précise, on peut supposer
que $k_1=k[T]/(r)$, où $r$ est un polynôme unitaire irréductible et séparable de
$k[T]$ de degré $n$~; si $R$ est un polynôme unitaire de $A[T]$ dont l'image
canonique est $r$ (et qui est donc de degré $n$), on peut prendre
$A_1=A[T]/(R)$. Or, si $K$ est le corps des fractions de $A$, il est clair que
$R$ est un polynôme irréductible et séparable de $K[T]$~; on en déduit donc
d'abord que $A_1$ est un anneau intègre dont le corps des fractions
$K_1=K[T]/(R)$ est une extension séparable de $K$. En outre, si $t$ est l'image
canonique de $T$ dans $A_1$, les $t^j$ $(0\leq j<n)$ forment une base du
$A$-module $A_1$, et leurs images dans $k_1$ une base sur $k$~; on déduit de là
que
\[
d=\det\bigl(\operatorname{Tr}_{A_1/A}(t^{i+j})\bigr)
\]
est un élément de $A$ dont la classe dans $k$ est $\ne0$, et qui est par suite
inversible. Le même raisonnement que dans (6.12.4.1, I)) prouve alors que le
morphisme $\operatorname{Spec}(A_1)\to\operatorname{Spec}(A)$ est plat et a ses
fibres régulières~; on conclut par suite de (6.5.4, (ii)) que $A_1$ est
\emph{intégralement clos}. Considérons alors l'anneau
$A'_1=A'\otimes_A A_1$~; c'est un $A'$-module libre de type fini, donc un anneau
semi-local (Bourbaki, \emph{Alg. comm.}, chap.~IV, \S~2, n$^\circ$~5, cor.~3 de
la prop.~9)~; en outre, les idéaux maximaux de cette $A'$-algèbre finie sont
tous au-dessus de l'idéal maximal $\mathfrak m'$ de $A'$, et \emph{a fortiori}
contiennent $\mathfrak m A'_1$. Mais
\[
A'_1/\mathfrak m A'_1=(A'/\mathfrak m A')\otimes_k k_1,
\]
et comme $k_1$ est une extension séparable et finie de $k$, le radical de
$A'_1/\mathfrak m A'_1$ est égal à
$(\mathfrak m'/\mathfrak m A')\otimes_k k_1$ (Bourbaki, \emph{Alg.},
chap.~VIII, \S~7, n$^\circ$~2, cor.~2 de la prop.~3)~; si $\mathfrak n_i$
$(1\leq i\leq r)$ sont les idéaux maximaux de $A'_1$, les corps
$A'_1/\mathfrak n_i$ sont donc les corps composant l'algèbre
$k'\otimes_k k_1$, autrement dit ce sont des extensions finies
\emph{radicielles} de $k_1$. D'ailleurs, comme $A\to A'$ est un homomorphisme
injectif, il en est de même de $A_1\to A'_1$, $A_1$ étant un $A$-module plat~;
l'homomorphisme canonique
\[
A'_1\longrightarrow\prod_{i=1}^{r}(A'_1)_{\mathfrak n_i}
\]
étant lui aussi injectif (Bourbaki, \emph{Alg. comm.}, chap.~II, \S~3,
n$^\circ$~3, th.~1), il en est de même du composé de $A_1$ dans ce produit.
Mais $A_1$ est intègre, et les noyaux des homomorphismes
$A_1\to(A'_1)_{\mathfrak n_i}$ sont en nombre fini~; comme leur intersection
est nulle, c'est que l'un d'eux déjà est nul. En d'autres termes, il y a un
$B_1=(A'_1)_{\mathfrak n_i}$ tel que l'homomorphisme $A_1\to B_1$ soit
\emph{injectif} et \emph{local}. Posons $Y'_1=\operatorname{Spec}(B_1)$,
$X'_1=X'\times_{Y'}Y'_1$~; $\mathcal F'_1=\mathcal F'\otimes_{Y'}Y'_1$ est
$Y'_1$-plat en tous les points de $X'_1$ au-dessus de $x'$~; d'ailleurs l'idéal
maximal de $B_1$ est le seul au-dessus de $\mathfrak m'$, par suite tous ces
points ont pour projection dans $Y'_1$ le point fermé $y'_1$. Soit $x'_1$ un de
ces points. Posons d'autre part $Y_1=\operatorname{Spec}(A_1)$,
$X_1=X\times_Y Y_1$, $\mathcal F_1=\mathcal F\otimes_Y Y_1$~; si $x_1$ est la
projection de $x'_1$ dans $X_1$, la projection de $x_1$ dans $X$ est $x$ et sa
projection dans $Y_1$ est le point fermé $y_1$. Si l'on prouve que
$(\mathcal F_1)_{x_1}$ est un $\mathcal O_{y_1}$-module plat, il en résultera
que $\mathcal F_x$ est un $\mathcal O_y$-module plat~; en effet
$\mathcal O_{y_1}$ est un $\mathcal O_y$-module plat, donc
$(\mathcal F_1)_{x_1}$ est un $\mathcal O_y$-module plat
$(\mathbf 0_{\mathrm I},6.2.1)$.
% IV3_B23_P156_END

\oldpage[IV-3]{157}

Mais $(\mathcal F_1)_{x_1}=\mathcal F_x\otimes_{\mathcal O_x}\mathcal O_{x_1}$,
et $\mathcal O_{x_1}$ est un $\mathcal O_x$-module \emph{fidèlement plat}~;
donc (2.2.11, (iii)) $\mathcal F_x$ est un $\mathcal O_y$-module plat. Comme
$X'_1=X_1\times_{Y_1}Y'_1$, $\mathcal F'_1=\mathcal F_1\otimes_{Y_1}Y'_1$,
on est bien ramené à la situation de l'énoncé (11.6.1), avec $A$ et $A'$
remplacés par $A_1$ et $B_1$ respectivement.

\medskip
\noindent\textbf{IV)} \emph{Fin de la démonstration.} --- On est finalement
réduit à prouver (11.6.1) sous les hypothèses supplémentaires suivantes~:
\begin{enumerate}[label=(\roman*)]
  \item $A$ et $A'$ sont des anneaux locaux de $\mathbf Z$-algèbres de type fini
  (donc des anneaux \emph{excellents} (7.8.3))~;
  \item $A$ est intégralement clos~;
  \item le corps résiduel $k'$ de $A'$ est une extension finie et radicielle du
  corps résiduel $k$ de $A$.
\end{enumerate}
On sait alors ((7.8.3) et (2.3.8)) que sous les conditions (i) et (ii), si
$\mathfrak m$ et $\mathfrak m'$ sont les idéaux maximaux de $A$ et $A'$
respectivement, la topologie $\mathfrak m$-adique sur $A$ est induite par la
topologie $\mathfrak m'$-adique de $A'$. Le complété
$\widehat u:\widehat A\to\widehat{A'}$ est donc un homomorphisme injectif.
D'autre part, comme le morphisme
$\operatorname{Spec}(k')\to\operatorname{Spec}(k)$ est radiciel, il en est de
même du morphisme $f'^{-1}(y')\to f^{-1}(y)$ (\textbf I, 3.5.7), et il n'y a
donc qu'un seul point $x'$ de $X'$ dont les projections dans $X$ et $Y'$ soient
$x$ et $y'$ respectivement. On peut donc appliquer le résultat de (11.5.2).
C.Q.F.D.
\end{proof}

\begin{corollary}[11.6.2]
\label{IV.11.6.2-fr}
Soient $A$ un anneau local intègre unibranche, $Y=\operatorname{Spec}(A)$,
$f:X\to Y$ un morphisme localement de présentation finie, $\mathcal F$ un
$\mathcal O_X$-Module quasi-cohérent de présentation finie. Soient $A'$ un
anneau local, $A\to A'$ un homomorphisme local injectif~; on pose
$Y'=\operatorname{Spec}(A')$, $X'=X\times_Y Y'$, $f'=f_{(Y')}$,
$\mathcal F'=\mathcal F\otimes_Y Y'$. Soit $x$ un point de $X$ dont la
projection $f(x)=y$ est le point fermé de $Y$~; on suppose que $\mathcal F'$ est
$f'$-plat en tous les points $x'$ de $X'$ dont les projections dans $X$ et $Y'$
sont respectivement $x$ et le point fermé $y'$ de $Y'$. Alors $\mathcal F$ est
$f$-plat au point $x$.
\end{corollary}

\begin{proof}
On peut en effet reprendre la partie I) de la démonstration de (11.6.1), qui
prouve (avec les mêmes notations) que si $\mathcal F_1$ est $f_1$-plat en tous
les points $x_1$ de $X_1$ dont les projections respectives dans $X$ et $Y_1$
sont $x$ et $y_1$, alors $\mathcal F$ est $f$-plat au point $x$~; on est ainsi
ramené au cas où $A$ est intégralement clos, donc géométriquement unibranche, et
la conclusion résulte alors de (11.6.1).
\end{proof}

\subsection{Contre-exemples}
\label{subsection:IV.11.7-fr}

\begin{env}[11.7.1]
\label{IV.11.7.1-fr}
Considérons d'abord le cas où $A$ est un anneau local \emph{artinien}, et où les
hypothèses de (11.4.11) sont remplies \emph{sauf} la condition (ii) concernant
le corps résiduel $k$ de $A$. Nous allons voir que la conclusion de (11.4.11)
peut alors être en défaut. Soit $k$ un corps admettant une extension
\emph{galoisienne} $k'$ de degré $[k':k]>1$, et désignons par $\Gamma$ le groupe
de Galois de $k'$. Soit $A$ une $k$-algèbre ayant une base de 3 éléments
$1,a,b$ avec la table de multiplication $a^2=ab=ba=b^2=0$, de sorte que $A$ est
un anneau local artinien dont l'idéal maximal $\mathfrak m=ka+kb$ est de carré
nul. Soit $A'=A\otimes_k k'$, qui est une $k'$-algèbre de base $1,a,b$, anneau
local artinien d'idéal maximal
$\mathfrak m'=k'a+k'b=\mathfrak m A'$, de carré nul~; $A$ s'identifie
canoniquement à un sous-anneau de $A'$. Soit $\mathfrak J$ le sous-$k'$-espace
vectoriel de $\mathfrak m'$ engendré par $a+\gamma b$, où $\gamma\in k'$
n'appartient pas à $k$~; il est clair que $\mathfrak J$ est un idéal de $A'$.
Posons $B=A'/\mathfrak J$~; c'est un anneau artinien qui est un $A$-module
\emph{non plat}~; sinon (Bourbaki, \emph{Alg. comm.}, chap.~II, \S~3,
n$^\circ$~2, cor.~2 de la prop.~5), $B$ serait un $A$-module \emph{libre}~;
comme $A'$ est lui aussi un $A$-module libre, et que l'homomorphisme canonique
$A'/\mathfrak m A'\to B/\mathfrak m B$ est bijectif, l'homomorphisme canonique
$A'\to B=A'/\mathfrak J$ serait lui
% IV3_B23_P157_END

\oldpage[IV-3]{158}

aussi bijectif (\emph{loc. cit.}, n$^\circ$~2, cor. de la prop.~6) ce qui est
absurde. Autrement dit, si l'on pose $Y=\operatorname{Spec}(A)$,
$X=\operatorname{Spec}(B)$, $\mathcal O_X$ n'est pas $Y$-plat en l'unique point
$x$ de $X$. Mais soit $A_1=B$, et $Y_1=\operatorname{Spec}(A_1)$, et considérons
l'homomorphisme canonique $A\to A_1$, qui est local et injectif puisque
$\mathfrak J\cap A=0$~; si
$B_1=B\otimes_A A_1=B\otimes_A B$, nous allons voir qu'il y a un point de
$X_1=X\times_Y Y_1=\operatorname{Spec}(B_1)$ où $\mathcal O_{X_1}$ est
$Y_1$-plat. Pour cela, remarquons que l'on a
\[
B\otimes_A B
=\frac{A'\otimes_A A'}
{\operatorname{Im}(\mathfrak J\otimes_A A')
+\operatorname{Im}(A'\otimes_A\mathfrak J)}.
\]
Or la structure de $C'=A'\otimes_A A'$ s'obtient aisément~; on considère la
$A$-algèbre produit $C=\prod_{\sigma\in\Gamma}A'_{\sigma}$, où tous les
$A'_{\sigma}$ sont égaux à $A'$, et l'application canonique
$\varphi:A'\otimes_A A'\to C$ telle que
$\varphi(x\otimes y)=(\sigma(x)y)_{\sigma\in\Gamma}$ (le groupe $\Gamma$
opérant canoniquement sur $A'$)~; en passant aux quotients, on en déduit un
homomorphisme $C'/\mathfrak m C'\to C/\mathfrak m C$ qui n'est autre que
l'homomorphisme canonique
$k'\otimes_k k'\to\prod_{\sigma}k'_{\sigma}$ (avec $k'_{\sigma}=k'$ pour tout
$\sigma\in\Gamma$)~; on sait que ce dernier est bijectif (Bourbaki,
\emph{Alg.}, chap.~VIII, \S~8, prop.~4), donc il en est de même de $\varphi$,
puisque $C'$ et $C$ sont des $A$-modules libres (Bourbaki,
\emph{Alg. comm.}, chap.~II, \S~3, n$^\circ$~2, cor. de la prop.~6). De ce qui
précède, il résulte que $B\otimes_A B$ est une $A$-algèbre semi-locale, composée
directe des $A$-algèbres locales $A'/(\sigma(\mathfrak J)+\mathfrak J)$. Celle
de ces algèbres correspondant à l'identité de $\Gamma$ est isomorphe à
$A'/\mathfrak J$, donc est un $A_1$-module plat, mais comme $\gamma\notin k$, il
y a au moins un $\sigma\in\Gamma$ tel que
$\sigma(\mathfrak J)\ne\mathfrak J$~; alors
$\sigma(\mathfrak J)+\mathfrak J=\mathfrak m'$, et $A'/\mathfrak m'$ n'est pas
un $(A'/\mathfrak J)$-module plat, puisque c'est le quotient de
$A'/\mathfrak J$ par un idéal non nul.
\end{env}

\begin{env}[11.7.2]
\label{IV.11.7.2-fr}
Nous allons maintenant montrer que le résultat de (11.5.4) perd sa validité
lorsqu'on ne suppose plus que $f$ soit un morphisme \emph{propre} (et
\emph{a fortiori} (11.5.3) cesse d'être exact lorsqu'on ne suppose plus $g$
propre). Soient $k$ un corps, $A_0$ l'anneau de polynômes $k[S,T]$, $A$
l'anneau quotient $A_0/A_0ST$~; $Y=\operatorname{Spec}(A)$ est donc la courbe
réductible formée des deux «~axes de coordonnées~» dans le plan affine
$V_k^2=\operatorname{Spec}(A_0)$. L'anneau $A$ admet deux idéaux premiers
minimaux $\mathfrak p_1=A_0S/A_0ST$, $\mathfrak p_2=A_0T/A_0ST$, et comme $A$
est réduit, il se plonge canoniquement dans $B=B_1\oplus B_2$, où
$B_1=A/\mathfrak p_1$, $B_2=A/\mathfrak p_2$~; d'ailleurs $B_1$ s'identifie à
$k[T]$ et $B_2$ à $k[S]$, donc ce sont des anneaux intègres intégralement clos
et par suite $Z=\operatorname{Spec}(B)$ n'est autre que le normalisé du
préschéma $Y$ par rapport à $R(Y)$ (\textbf{II}, 6.3.8), somme des deux schémas
$Z_1=\operatorname{Spec}(B_1)$, $Z_2=\operatorname{Spec}(B_2)$. Désignons par
$y$ le «~point double~» de $Y$, correspondant à l'idéal maximal
$\mathfrak p_1+\mathfrak p_2=\mathfrak m$ de $A$, par $z_1$ et $z_2$ les
points de $Z$ qui se projettent en $y$, correspondant respectivement aux idéaux
maximaux $\mathfrak m_1=(T)$ et $\mathfrak m_2=(S)$ de $B_1$ et $B_2$. Nous
désignerons par $X$ le sous-préschéma de $Z$ induit sur le complémentaire de
$z_2$ dans $Z$~; on a donc
$X=\operatorname{Spec}(B_1\oplus(B_2)_S)$~; il est immédiat que
l'homomorphisme $A\to A'=B_1\oplus(B_2)_S$ est injectif, mais le morphisme
correspondant $f:X\to Y$ n'est pas fermé (car $f(Z_2-\{z_2\})$ n'est pas fermé
dans $Y$, bien que $Z_2-\{z_2\}$ soit fermé dans $X$)~; \emph{a fortiori} il
n'est pas propre. Nous allons voir maintenant que $f$ n'est pas plat au point
$z_1$~; il suffira de montrer $(\mathbf 0_{\mathrm I},6.6.2)$ que
$\mathcal O_{z_1}$ n'est pas un $\mathcal O_y$-module fidèlement plat, et pour
cela il suffit de voir que l'homomorphisme canonique
$\mathcal O_y\to\mathcal O_{z_1}$ n'est pas injectif~; mais cela est immédiat
car $\mathcal O_{z_1}$ est un anneau intègre, tandis que $\mathcal O_y$ admet
des diviseurs de zéro. Cependant, la première projection
$p:X\times_Y X\to X$ est un isomorphisme~: en effet, on a
\[
B_1\otimes_A B_1=(A/\mathfrak p_1)\otimes_A(A/\mathfrak p_1)=A/\mathfrak p_1,
\]
$(B_2)_S\otimes_A(B_2)_S=(B_2\otimes_A B_2)_S=(B_2)_S$ pour la même raison,
et enfin $B_1\otimes_A(B_2)_S=0$, car l'image canonique de $S$ dans $B_1$ est
nulle.
\end{env}

\begin{env}[11.7.3]
\label{IV.11.7.3-fr}
L'exemple précédent peut se généraliser~: on considère sur un corps $k$ une
courbe algébrique réduite $Y$ admettant un seul «~point double ordinaire~» $y$
(notion qui sera définie plus tard de façon générale), et sa normalisée $Z$, de
sorte que le morphisme $g:Z\to Y$ est fini, que la restriction de $g$ à
$Z-g^{-1}(y)$ est un isomorphisme sur $Y-\{y\}$, et que $g^{-1}(y)$ se réduit à
deux points «~simples~» $z_1,z_2$~; en outre le préschéma $g^{-1}(y)$ est somme
des deux préschémas $\operatorname{Spec}(k(z_1))$,
$\operatorname{Spec}(k(z_2))$, canoniquement isomorphes à
$\operatorname{Spec}(k(y))$. Soit $X$ le sous-préschéma de $Z$ induit sur
l'ouvert $Z-\{z_2\}$~; le morphisme $f:X\to Y$, restriction de $g$ à $X$, n'est
pas propre, sans quoi (\textbf{II}, 5.4.3) il en serait de même de l'injection
canonique $j:X\to Z$, qui n'est pas fermée. Le morphisme $f$ est radiciel, car
pour tout $y'\in Y$, la fibre $f^{-1}(y')$ ne comprend qu'un seul point $x'$,
$k(y')\to k(x')$ étant un isomorphisme~; on en conclut d'abord que le morphisme
diagonal $\Delta_f:X\to X\times_Y X$ est bijectif (1.7.7.1) et d'autre part,
comme $f$ est non ramifié (17.4.2, d$'$), $\Delta_f$ est une immersion ouverte
(17.4.2, b$'$)~; par suite $\Delta_f$ est un isomorphisme, et la première
projection $p:X\times_Y X\to X$ l'isomorphisme réciproque. Cependant $f$ n'est
pas plat au point $z_1$~; sinon $\mathcal O_{z_1}$ serait un
$\mathcal O_y$-module fidèlement plat $(\mathbf 0_{\mathrm I},6.6.2)$, et comme
$\mathcal O_y$ contient deux idéaux premiers minimaux $\mathfrak p_1$,
$\mathfrak p_2$ distincts (correspondant aux deux «~branches~» de $Y$ au point
$y$) il existerait dans $\mathcal O_{z_1}$ deux idéaux premiers dont les images
réciproques par $u:\mathcal O_y\to\mathcal O_{z_1}$ seraient
$\mathfrak p_1$ et $\mathfrak p_2$ $(\mathbf 0_{\mathrm I},6.5.1)$~; mais cela
est absurde, car $\mathcal O_{z_1}$ n'a que deux idéaux premiers distincts, $0$
et l'idéal maximal $\mathfrak m_1$, et $u^{-1}(\mathfrak m_1)$ est l'idéal
maximal $\mathfrak m$ de $\mathcal O_y$.
\end{env}

\begin{env}[11.7.4]
\label{IV.11.7.4-fr}
On notera que dans l'exemple précédent l'homomorphisme $u$ est \emph{injectif}
lorsque $Y$ est \emph{irréductible} (on peut par exemple prendre
$Y=\operatorname{Spec}(k[S,T]/(S(S^2+T^2)-(S^2-T^2)))$, «~cubique à point
double~»)~; on peut en effet alors (en remplaçant $Y$ par un voisinage affine de
$y$) supposer que $Y=\operatorname{Spec}(A)$, où $A$ est intègre, d'où
$Z=\operatorname{Spec}(B)$, où $B$ est la clôture intégrale de $A$~; comme $B$,
$\mathcal O_y$ et $\mathcal O_{z_1}$ s'identifient alors à des sous-anneaux du
corps
% IV3_B23_P158_END

\oldpage[IV-3]{159}

des fractions de $A$, $u$ est évidemment injectif. On notera par contre que
l'homomorphisme $\widehat u:\widehat{\mathcal O}_y\to
\widehat{\mathcal O}_{z_1}$ n'est pas injectif, car
$\widehat{\mathcal O}_{z_1}$ est un anneau intègre local ($z_1$ étant un point
simple), tandis que $\widehat{\mathcal O}_y$ a deux idéaux premiers minimaux
distincts (correspondant aux deux «~branches~» de $Y$) et admet donc des
diviseurs de zéro. Cela donne un exemple montrant que dans l'énoncé (11.5.2), on
ne peut pas remplacer l'hypothèse que $\widehat u$ est injectif par l'hypothèse
que $u$ lui-même est injectif, même lorsque $A'$ est un anneau local. Il suffit
en effet de prendre (avec les notations précédentes) $A=\mathcal O_y$,
$A'=\mathcal O_{z_1}$, $Y=\operatorname{Spec}(A)$,
$X=\operatorname{Spec}(A')$~; le raisonnement de (11.7.3) prouve encore que la
première projection $p:X\times_Y X\to X$ est un isomorphisme, bien que $f$ ne
soit pas plat au point $z_1$.
\end{env}

\begin{env}[11.7.5]
\label{IV.11.7.5-fr}
Les exemples de (11.7.2) et (11.7.3) expliquent la restriction aux anneaux
locaux \emph{unibranches} dans (11.6.1) et (11.6.2). Nous allons voir maintenant
que dans (11.6.1), on ne peut affaiblir l'hypothèse sur $A$ en supposant
seulement $A$ unibranche. Considérons en effet l'anneau local intègre complet
$A=\mathbf R[[U,V]]/(U^2+V^2)$ qui est unibranche mais non géométriquement
unibranche (6.5.11). On sait (\emph{loc. cit.}) que si $u,v$ sont les images de
$U$ et $V$ dans $A$, la clôture intégrale de $A$ est $\overline A=A[t]$ avec
$t=v/u$, tel que $t^2=-1$, si bien que $\overline A$ est isomorphe à
$\mathbf C[[U]]$. Posons $Y=\operatorname{Spec}(A)$,
$X=\operatorname{Spec}(\overline A)$ (normalisé de $Y$ (\textbf{II}, 6.3.8)) et
soient $y$ et $x$ les points fermés de $Y$ et $X$ respectivement~; nous allons
montrer que pour une $A$-algèbre locale convenable $A'$, si l'on pose
$Y'=\operatorname{Spec}(A')$, $X'=X\times_Y Y'$, et si $y'$ désigne le point
fermé de $Y'$, $\mathcal O_{X'}$ est $Y'$-plat en un point de $X'$ dont les
projections dans $X$ et $Y'$ sont respectivement $x$ et $y'$, mais n'est pas
$Y'$-plat en tous les points ayant ces projections~; il en résultera évidemment
(2.1.4) que $\mathcal O_X$ n'est pas $Y$-plat au point $x$ (ce qui est d'ailleurs
trivial \emph{a priori}, $\overline A$ n'étant pas un $A$-module libre).

Soit $B=A\otimes_{\mathbf R}\mathbf C$, isomorphe à
$\mathbf C[[U,V]]/(U+iV)(U-iV)$~; $B$ a deux idéaux premiers minimaux
$\mathfrak p'$, $\mathfrak p''$ engendrés respectivement par $u+iv$ et $u-iv$,
et $\mathfrak n=\mathfrak p'+\mathfrak p''$ est l'idéal maximal de l'anneau
local complet $B$. Soit $\overline B=\overline A\otimes_{\mathbf R}\mathbf C$~;
$\overline B$ est composée directe de deux algèbres isomorphes à
$\mathbf C[[U]]$, engendrées par les idempotents
$e'=(1\otimes1+t\otimes i)/2$ et $e''=(1\otimes1-t\otimes i)/2$~; comme
l'homomorphisme $A\to\overline A$ est injectif, il en est de même de
$B\to\overline B$, et les images de $u+iv$ et de $u-iv$ par cette injection sont
respectivement $ue'$ et $ue''$~; on en conclut aussitôt que $\overline B$
s'identifie canoniquement à $(B/\mathfrak p')\oplus(B/\mathfrak p'')$. Cela
étant, prenons pour $A'$ l'anneau local $B/\mathfrak p'$~; alors
$\overline A\otimes_A A'$ s'identifie à $\overline B\otimes_B A'$. Mais on a
$(B/\mathfrak p')\otimes_B(B/\mathfrak p')=B/\mathfrak p'$ et
$(B/\mathfrak p')\otimes_B(B/\mathfrak p'')=B/\mathfrak n$, donc
$\overline A\otimes_A A'$ est isomorphe à $A'\oplus(B/\mathfrak n)$. Cela
établit notre assertion, car $B/\mathfrak n=A'/(\mathfrak n/\mathfrak p')$
n'est pas un $A'$-module plat (sans quoi ce serait un $A'$-module libre
(Bourbaki, \emph{Alg. comm.}, chap.~II, \S~3, n$^\circ$~2, cor.~2 de la
prop.~5), ce qui est absurde).
\end{env}

\subsection{Un critère valuatif de platitude}
\label{subsection:IV.11.8-fr}

\begin{theorem}[11.8.1]
\label{IV.11.8.1-fr}
Soient $f:X\to Y$ un morphisme localement de présentation finie, $\mathcal F$ un
$\mathcal O_X$-Module quasi-cohérent de présentation finie, $x$ un point de
$X$, $y=f(x)$. On suppose l'anneau local $\mathcal O_y$ intègre (resp. réduit et
noethérien). Pour que $\mathcal F$ soit $f$-plat au point $x$, il faut et il
suffit que, pour tout anneau de valuation (resp. tout anneau de valuation
discrète) $A'$ et tout homomorphisme local $\mathcal O_y\to A'$, la condition
suivante soit satisfaite~: en posant $Y'=\operatorname{Spec}(A')$,
$X'=X\times_Y Y'$, $f'=f_{(Y')}$, le $\mathcal O_{X'}$-Module
$\mathcal F'=\mathcal F\otimes_Y Y'$ est $f'$-plat en tous les points $x'$ de
$X'$ dont les projections respectives dans $X$ et $Y'$ sont $x$ et le point
fermé $y'$ de $Y'$.
\end{theorem}

\begin{proof}
La condition étant évidemment nécessaire (2.1.4), reste à prouver qu'elle est
suffisante. On peut évidemment ((\textbf I, 3.6.5) et (\textbf I, 2.4.4)) se
borner au cas où $Y=\operatorname{Spec}(A)$ est le spectre d'un anneau local $A$
et $y$ le point fermé de $Y$.

\medskip
\noindent\textbf{(i)} \emph{Cas où $A$ est intègre.} --- Soient $K$ le corps
des fractions de $A$, $A_1$ la clôture intégrale de $A$~; si l'on pose
$Y_1=\operatorname{Spec}(A_1)$, $X_1=X\times_Y Y_1$, $f_1=f_{(Y_1)}$, il suffit,
en vertu de (11.5.5), de montrer que
$\mathcal F_1=\mathcal F\otimes_Y Y_1$ est $f_1$-plat en \emph{tous} les points
$x_1$ de $X_1$ dont $x$ est la projection dans $X$. Or, si $f_1(x_1)=y_1$, on a
$\mathfrak i_{y_1}=\mathfrak m_1$, où $\mathfrak m_1$ est un idéal premier de
$A_1$ dont $\mathfrak m=\mathfrak i_y$ est la trace sur $A$ ($\mathfrak m_1$
est d'ailleurs nécessairement un idéal maximal). Soit alors $A'$ un anneau de
valuation pour $K$ qui domine $\mathcal O_{y_1}=(A_1)_{\mathfrak m_1}$~;
l'homomorphisme $A\to\mathcal O_{y_1}$ étant local, il en est de même de
$A\to A'$. Il existe alors au moins un
% IV3_B23_P159_END

\oldpage[IV-3]{160}

point $x'\in X'$ dont les projections dans $X_1$ et $Y'$ sont respectivement
$x_1$ et $y'$ (\textbf I, 3.4.7)~; comme par hypothèse $\mathcal F'$ est
$f'$-plat au point $x'$, et que $\mathcal O_{y_1}$ est intégralement clos, on
peut appliquer (11.6.1), et on en déduit que $\mathcal F_1$ est $f_1$-plat au
point $x_1$, d'où le théorème dans ce cas.

\medskip
\noindent\textbf{(ii)} \emph{Cas où $A$ est réduit et noethérien.} --- Soient
$\mathfrak p_i$ $(1\leq i\leq m)$ les idéaux minimaux de $A$~; comme $A$ est
réduit, il s'identifie canoniquement à un sous-anneau du produit des
$A_i=A/\mathfrak p_i$, qui sont des anneaux locaux noethériens~; posant
$Y_i=\operatorname{Spec}(A_i)$, $X_i=X\times_Y Y_i$, $f_i=f_{(Y_i)}$, il
résulte de (11.5.2) qu'il suffit de démontrer que pour chaque $i$,
$\mathcal F_i=\mathcal F\otimes_Y Y_i$ est $f_i$-plat en tout point $x_i$ de
$X_i$ dont les projections dans $X$ et $Y_i$ sont $x$ et le point fermé $y_i$
de $Y_i$ respectivement. Or, comme $A_i$ est un anneau local noethérien intègre,
il existe dans son corps des fractions $K_i$ un sous-anneau $A''_i$ contenant
$A_i$, qui soit une $A_i$-algèbre \emph{finie} (donc un anneau semi-local
noethérien) et dont les anneaux locaux soient \emph{géométriquement
unibranches} ((6.15.5) et $(\mathbf 0,23.2.5)$). Comme les idéaux maximaux de
$A''_i$ sont alors nécessairement au-dessus de l'idéal maximal de $A_i$, on
déduit encore de (\textbf I, 3.4.7) et de (11.5.2) qu'il suffit, en posant
$Y''_i=\operatorname{Spec}(A''_i)$, $X''_i=X\times_Y Y''_i$,
$f''_i=f_{(Y''_i)}$, de prouver que
$\mathcal F''_i=\mathcal F\otimes_Y Y''_i$ est $f''_i$-plat en tout point
$x''_i$ de $X''_i$ dont les projections dans $X$ et $Y''_i$ sont $x$ et le
point fermé $y''_i$ de $Y''_i$ respectivement. Or, soit $A'$ un anneau de
valuation discrète pour $K_i$ dominant $A''_i$, et soit $x'$ un point de $X'$
dont les projections dans $X''_i$ et dans $Y'$ sont $x''_i$ et $y'$
respectivement (\textbf I, 3.4.7)~; comme $\mathcal O_{y''_i}$ est
géométriquement unibranche, on peut encore appliquer (11.6.1) et on en déduit
bien que $\mathcal F''_i$ est $f''_i$-plat au point $x''_i$.
\end{proof}

\begin{remarks}[11.8.2]
\label{IV.11.8.2-fr}
\begin{enumerate}[label=(\roman*)]
  \item Dans l'énoncé de (11.8.1), on peut se borner à supposer que la condition
  sur $\mathcal F'$ est vérifiée pour les anneaux de valuation $A'$ \emph{complets}
  et dont le corps résiduel est \emph{algébriquement clos}. On sait en effet que
  tout anneau de valuation $A'$ est dominé par un tel anneau $A''$
  (\textbf{II}, 7.1.2), et que si $A'$ est un anneau de valuation discrète, on
  peut supposer qu'il en est de même de $A''$
  $(\mathbf 0_{\mathrm{III}},10.3.1)$.
  \item La démonstration de (11.8.1) se simplifie lorsqu'on suppose non seulement
  que $A$ est intègre et noethérien, mais que son complété $\widehat A$ est lui
  aussi \emph{intègre}. Remplaçant $X$ par $X\otimes_A\widehat A$ et raisonnant
  comme dans la démonstration de (11.5.3), on peut en effet se ramener alors à
  prouver (11.8.1) lorsque $A=\mathcal O_y$ est intègre, noethérien et
  \emph{complet}. Or, on sait (\textbf{II}, 7.1.7) qu'un tel anneau $A$ est
  dominé par un anneau de valuation discrète complet~; la conclusion résulte
  donc directement de (11.5.2).
\end{enumerate}
\end{remarks}

\subsection{Familles séparantes et universellement séparantes d'homomorphismes
de faisceaux de modules}
\label{subsection:IV.11.9-fr}

\begin{env}[11.9.1]
\label{IV.11.9.1-fr}
Soient $X$ un préschéma, $(f_{\lambda})_{\lambda\in L}$ une famille de morphismes
$f_{\lambda}:Z_{\lambda}\to X$, $\mathcal F$ un $\mathcal O_X$-Module
quasi-cohérent~; pour tout $\lambda\in L$, supposons donnés un
$\mathcal O_{Z_{\lambda}}$-Module quasi-cohérent $\mathcal G_{\lambda}$ et un
homomorphisme
\begin{equation}
u_{\lambda}:\mathcal F\longrightarrow(f_{\lambda})_*(\mathcal G_{\lambda}).
\label{IV.11.9.1.1-fr}
\tag{11.9.1.1}
\end{equation}
\end{env}
% IV3_B23_P160_END

\oldpage[IV-3]{161}

On dit que la famille $(u_\lambda)$ (ou la famille correspondante des
$u^\#_\lambda:(f_\lambda)^*(\mathcal F)\to\mathcal G_\lambda$) est
\emph{séparante} si l'intersection des noyaux des $u_\lambda$ est
\emph{nulle}. En d'autres termes, cela signifie que pour \emph{tout} ouvert
$U$ de $X$, et toute section $t$ de $\mathcal F$ au-dessus de $U$, telle que,
pour \emph{tout} $\lambda$, la section $u_\lambda(t)$ (qui, par définition est
une section de $\mathcal G_\lambda$ au-dessus de $f_\lambda^{-1}(U)$) soit
nulle, alors $t$ est elle-même nulle.

\begin{env}[11.9.2]
\label{IV.11.9.2-fr}
Avec les notations de (11.9.1), soit $M$ un second ensemble d'indices~; pour
tout $\lambda\in L$, soit $(g_{\lambda\mu})_{\mu\in M}$ une famille de
morphismes $g_{\lambda\mu}:Z'_{\lambda\mu}\to Z_\lambda$~; pour tout couple
$(\lambda,\mu)$, supposons donnés un $\mathcal O_{Z'_{\lambda\mu}}$-Module
quasi-cohérent $\mathcal H_{\lambda\mu}$ et un homomorphisme
$v_{\lambda\mu}:\mathcal G_\lambda\to
(g_{\lambda\mu})_*(\mathcal H_{\lambda\mu})$~; posons
$h_{\lambda\mu}=f_\lambda\circ g_{\lambda\mu}:Z'_{\lambda\mu}\to X$ et
considérons l'homomorphisme composé
\[
w_{\lambda\mu}:\mathcal F\xrightarrow{u_\lambda}
(f_\lambda)_*(\mathcal G_\lambda)
\xrightarrow{(f_\lambda)_*(v_{\lambda\mu})}
(h_{\lambda\mu})_*(\mathcal H_{\lambda\mu}).
\]
Supposons que, pour tout $\lambda\in L$, la famille
$(v_{\lambda\mu})_{\mu\in M}$ soit séparante~: alors, il en est de même de la
famille des $(f_\lambda)_*(v_{\lambda\mu})$ ($\mu\in M$), comme on le voit
aussitôt. On en conclut que, pour que la famille $(u_\lambda)$ soit séparante,
il faut et il suffit que la famille
$(w_{\lambda\mu})_{(\lambda,\mu)\in L\times M}$ le soit.
\end{env}

\begin{env}[11.9.3]
\label{IV.11.9.3-fr}
Pour vérifier que la famille $(u_\lambda)$ est séparante (avec les notations de
(11.9.1)), on peut évidemment se ramener tout d'abord au cas où $X$ est
\emph{affine}, la propriété étant locale sur $X$. On peut en outre supposer que
$Z_\lambda=X$ \emph{pour tout} $\lambda\in L$. En effet, soit $M$ un ensemble
d'indices, somme d'une famille $(M_\lambda)_{\lambda\in L}$, et pour tout
$\lambda\in L$, soit $(Y_{\lambda\mu})_{\mu\in M_\lambda}$ un recouvrement
ouvert affine de $Z_\lambda$~; soit
$j_{\lambda\mu}:Y_{\lambda\mu}\to Z_\lambda$ l'injection canonique et posons
$\mathcal H_{\lambda\mu}=j_{\lambda\mu}^*(\mathcal G_\lambda)
=\mathcal G_\lambda|Y_{\lambda\mu}$. Si l'on considère l'homomorphisme
canonique
$v_{\lambda\mu}=\rho_{\mathcal G_\lambda}:\mathcal G_\lambda\to
(j_{\lambda\mu})_*(j_{\lambda\mu}^*(\mathcal G_\lambda))
=(j_{\lambda\mu})_*(\mathcal H_{\lambda\mu})$ relatif à $j_{\lambda\mu}$
$(\mathbf 0_{\mathrm I},4.4.3.2)$, il est immédiat que pour chaque
$\lambda\in L$, la famille $(v_{\lambda\mu})_{\mu\in M_\lambda}$ est
séparante. En vertu de (11.9.2), on est donc ramené à prouver que la famille des
homomorphismes composés $((f_\lambda)_*(v_{\lambda\mu}))\circ u_\lambda$ est
séparante, autrement dit on est ramené au cas où les $Z_\lambda$ sont affines.
Mais alors les $(f_\lambda)_*(\mathcal G_\lambda)$ sont des
$\mathcal O_X$-Modules quasi-cohérents (\textbf I, 1.6.2) et en vertu de la
définition, on peut remplacer les $Z_\lambda$ par $X$ et les
$\mathcal G_\lambda$ par les $(f_\lambda)_*(\mathcal G_\lambda)$, d'où notre
assertion.

On notera en outre que si $L$ est \emph{fini} et les $f_\lambda$
\emph{quasi-compacts}, on peut, dans la réduction précédente, supposer que les
$M_\lambda$ sont aussi \emph{finis}, donc on est dans ce cas ramené à vérifier
qu'une famille \emph{finie} d'homomorphismes de $\mathcal F$ dans des
$\mathcal O_X$-Modules quasi-cohérents est séparante.
\end{env}

\begin{env}[11.9.4]
\label{IV.11.9.4-fr}
Considérons donc le cas où $Z_\lambda=X$ pour tout $\lambda$, et où
$X=\operatorname{Spec}(A)$ est affine~; alors on a
$\mathcal F=\widetilde M$ et $\mathcal G_\lambda=\widetilde N_\lambda$, où
$M$ et $N_\lambda$ sont des $A$-modules, et $u_\lambda=\widetilde\varphi_\lambda$,
où les $\varphi_\lambda:M\to N_\lambda$ sont des $A$-homomorphismes. Dire que
la famille $(u_\lambda)$ est séparante signifie alors que, pour tout $s\in A$,
l'intersection des noyaux des
$(\varphi_\lambda)_s:M_s\to(N_\lambda)_s$ est réduite à 0. On dit alors aussi
que la famille $(\varphi_\lambda)$ est \emph{séparante}. On notera que si $L$
est \emph{fini}, il revient au même de dire que l'intersection des noyaux des
$\varphi_\lambda$ est 0, car on a alors
\[
\bigcap_{\lambda\in L}\operatorname{Ker}((\varphi_\lambda)_s)
=\left(\bigcap_{\lambda\in L}\operatorname{Ker}(\varphi_\lambda)\right)_s
\quad(\mathbf 0_{\mathrm I},1.3.2).
\]
Mais cette relation n'est plus exacte en général lorsque $L$ est \emph{infini},
et le fait que l'intersection des noyaux des $\varphi_\lambda$ soit 0
\emph{n'entraîne donc pas},
\oldpage[IV-3]{162}
en général, que la famille $(\varphi_\lambda)$ soit séparante. Par exemple,
supposons que $A$ soit un anneau de valuation discrète, d'idéal maximal
$\mathfrak m$, et considérons la famille des homomorphismes
$\varphi_k:A\to A/\mathfrak m^k$, dont l'intersection des noyaux
$\mathfrak m^k$ est réduite à 0~; cette famille n'est toutefois pas séparante,
car les fibres de tous les $\mathfrak m^k$ au point générique $x$ de
$X=\operatorname{Spec}(A)$ (qui est ouvert dans $X$) sont égales à
$\mathcal O_x=k(x)$, corps des fractions de $A$, et leur intersection n'est
donc pas réduite à 0.
\end{env}

\begin{env}[11.9.5]
\label{IV.11.9.5-fr}
Nous allons principalement nous occuper dans ce qui suit du problème du
\emph{changement de base} pour les familles séparantes. Les notations étant
celles de (11.9.1), considérons un morphisme $g:X'\to X$ et posons
$\mathcal F'=\mathcal F\otimes_{\mathcal O_X}\mathcal O_{X'}
=g^*(\mathcal F)$, et, pour tout $\lambda$,
$Z'_\lambda=Z_\lambda\times_X X'$, $f'_\lambda=f_\lambda\times 1$,
$\mathcal G'_\lambda=\mathcal G_\lambda\otimes_{\mathcal O_{Z_\lambda}}
\mathcal O_{Z'_\lambda}=(g'_\lambda)^*(\mathcal G_\lambda)$, où
$g'_\lambda:Z'_\lambda\to Z_\lambda$ est la projection canonique. Pour tout
$\lambda$, on désigne alors par
$u'_\lambda:\mathcal F'\to(f'_\lambda)_*(\mathcal G'_\lambda)$
l'homomorphisme obtenu comme suit~: soit
\[
\psi:(f_\lambda)_*(\mathcal G_\lambda)\longrightarrow
(f_\lambda)_*((g'_\lambda)_*((g'_\lambda)^*(\mathcal G_\lambda)))
=g_*((f'_\lambda)_*(\mathcal G'_\lambda))
\]
l'homomorphisme $(f_\lambda)_*(\rho_{\mathcal G_\lambda})$, où
$\rho_{\mathcal G_\lambda}$ est l'homomorphisme canonique
$(\mathbf 0_{\mathrm I},4.4.3.2)$ correspondant à $g'_\lambda$. Alors
$u'_\lambda$ est défini comme le composé
\[
g^*(\mathcal F)\xrightarrow{g^*(u_\lambda)}
g^*((f_\lambda)_*(\mathcal G_\lambda))
\xrightarrow{g^*(\psi)}g^*g_*((f'_\lambda)_*(\mathcal G'_\lambda))
\xrightarrow{\sigma}(f'_\lambda)_*(\mathcal G'_\lambda)
\]
où $\sigma=\sigma_{(f'_\lambda)_*(\mathcal G'_\lambda)}$ est
l'homomorphisme canonique $(\mathbf 0_{\mathrm I},4.4.3.3)$ correspondant à
$g$. On dira pour abréger que $u'_\lambda$ est \emph{déduit de $u_\lambda$
par le changement de base} $g$. Lorsque $f_\lambda:Z_\lambda\to X$ est un
morphisme \emph{affine}, on a
$(f'_\lambda)_*(\mathcal G'_\lambda)=g^*((f_\lambda)_*(\mathcal G_\lambda))$,
et dans ce cas on a donc simplement $u'_\lambda=g^*(u_\lambda)$.

On peut aussi interpréter $u'_\lambda$ de la façon suivante~: il suffit de
connaître la valeur de $u'_\lambda(t')$, lorsque $t'$ est une section de
$\mathcal F'$ au-dessus d'un ouvert $U'$ de $X'$, du type particulier
suivant~: $t'$ est la restriction à $U'$ de l'image canonique par
$\Gamma(U,\mathcal F)\to\Gamma(g^{-1}(U),\mathcal F')$ d'une section $t$ de
$\mathcal F$ au-dessus d'un ouvert $U$ de $X$ contenant $g(U')$ (ces sections
engendrant en effet le $\mathcal O_{X'}$-Module $\mathcal F'$
$(\mathbf 0_{\mathrm I},3.7.1)$). Considérons la section $u_\lambda(t)$ de
$\mathcal G_\lambda$ au-dessus de $f_\lambda^{-1}(U)$, et son image canonique
$t''$ par
$\Gamma(f_\lambda^{-1}(U),\mathcal G_\lambda)\to
\Gamma((g'_\lambda)^{-1}(f_\lambda^{-1}(U)),\mathcal G'_\lambda)$~; alors
$u'_\lambda(t')$ est la restriction de $t''$ à $(f'_\lambda)^{-1}(U')$.
Considérons en particulier le cas où $Z_\lambda$ est un sous-préschéma induit
sur un ouvert de $X$, $\mathcal G_\lambda=j_\lambda^*(\mathcal F)$, où
$j_\lambda:Z_\lambda\to X$ est l'injection canonique, et où $u_\lambda$ est
l'homomorphisme canonique
$\rho_{\mathcal F}:\mathcal F\to(j_\lambda)_*(j_\lambda^*(\mathcal F))
=(j_\lambda)_*(\mathcal G_\lambda)$ $(\mathbf 0_{\mathrm I},4.4.3.2)$
correspondant à $j_\lambda$. Alors $Z'_\lambda$ est induit sur un ouvert de
$X'$, et l'interprétation précédente montre que $u'_\lambda$ n'est autre que
l'homomorphisme canonique
$\rho_{\mathcal F'}:\mathcal F'\to
(j'_\lambda)_*((j'_\lambda)^*(\mathcal F'))
=(j'_\lambda)_*(\mathcal G'_\lambda)$ correspondant à l'injection canonique
$j'_\lambda:Z'_\lambda\to X'$.
\end{env}

\begin{env}[11.9.6]
\label{IV.11.9.6-fr}
Sous les conditions de (11.9.5), supposons que $X$ et $X'$ soient affines, et
que l'on veuille prouver que pour toute section $t'$ de $\mathcal F'$ au-dessus
de $X'$, dont les images par tous les $u'_\lambda$ sont nulles, alors $t'$ est
elle-même nulle. Alors, on peut aussi se borner au cas où $Z_\lambda=X$ pour
tout $\lambda\in L$. En effet, avec les notations de (11.9.3) et de (11.9.5),
si l'on pose $Y'_{\lambda\mu}=(g'_\lambda)^{-1}(Y_{\lambda\mu})$,
l'homomorphisme $v'_{\lambda\mu}$ déduit de $v_{\lambda\mu}$ par le changement
de base $g$ n'est autre que l'homomorphisme canonique
\oldpage[IV-3]{163}
$\rho_{\mathcal G'_\lambda}:\mathcal G'_\lambda\to
(j'_{\lambda\mu})_*((j'_{\lambda\mu})^*(\mathcal G'_\lambda))$
correspondant à l'injection canonique
$j'_{\lambda\mu}:Y'_{\lambda\mu}\to Z'_\lambda$, comme on l'a vu dans
(11.9.5). L'assertion résulte alors des raisonnements de (11.9.2) et (11.9.3),
$Y_{\lambda\mu}$ et $v_{\lambda\mu}$ étant remplacés par
$Y'_{\lambda\mu}$ et $v'_{\lambda\mu}$.

On a une réduction semblable lorsque l'on veut prouver que la famille
$(u'_\lambda)$ est \emph{séparante} ($X$ et $X'$ étant affines)~: cela résulte
encore de (11.9.2) et (11.9.3).
\end{env}

\begin{proposition}[11.9.7]
\label{IV.11.9.7-fr}
Avec les notations de (11.9.1) et (11.9.5), supposons
$X=\operatorname{Spec}(A)$ et $X'=\operatorname{Spec}(A')$ affines, et
supposons en outre que $A'$ soit un $A$-module projectif. Alors, si la famille
$(u_\lambda)$ est séparante, toute section $t'$ de $\mathcal F'$ au-dessus de
$X'$ dont les images par tous les $u'_\lambda$ sont nulles, est elle-même
nulle.
\end{proposition}

\begin{proof}
On a vu (11.9.6) qu'on peut se borner au cas où tous les $Z_\lambda$ sont égaux
à $X$. La proposition est alors conséquence du lemme suivant~:
\end{proof}

\begin{lemma}[11.9.7.1]
\label{IV.11.9.7.1-fr}
Soient $A$ un anneau, $(M_\lambda)_{\lambda\in L}$ une famille de $A$-modules,
$M$ un $A$-module et pour chaque $\lambda$,
$u_\lambda:M\to M_\lambda$ un homomorphisme. Supposons que l'intersection des
noyaux des $u_\lambda$ soit réduite à 0. Alors, pour tout $A$-module projectif
$N$, l'intersection des noyaux des homomorphismes
$u_\lambda\otimes 1:M\otimes_A N\to M_\lambda\otimes_A N$ est réduite à 0.
\end{lemma}

\begin{proof}
En effet, $N$ est facteur direct d'un $A$-module libre $L$, et il suffit
évidemment de prouver que l'intersection des noyaux des homomorphismes
$v_\lambda:M\otimes_A L\to M_\lambda\otimes_A L$ est réduite à 0, puisque
$u_\lambda\otimes 1:M\otimes_A N\to M_\lambda\otimes_A N$ est la restriction
de $v_\lambda$. Mais l'assertion résulte alors trivialement de l'hypothèse.
\end{proof}

\begin{remark}[11.9.8]
\label{IV.11.9.8-fr}
Nous ignorons si, sous les hypothèses de la proposition (11.9.7), la famille
$(u'_\lambda)$ est \emph{séparante}~: il faudrait en effet (11.9.4) prouver
qu'une section $t'$ de $\mathcal F'$ \emph{au-dessus d'un ouvert}
$D(h')\subset X'$ (où $h'\in A'$) telle que les $u'_\lambda(t')$ soient toutes
nulles, est elle-même nulle. Or, on ne peut appliquer la proposition (11.9.4)
à $D(h')=\operatorname{Spec}(A'_{h'})$, car du fait que $A'$ soit un
$A$-module projectif (même libre), il ne résulte pas que $A'_{h'}$ soit un
$A$-module projectif. Par exemple, on peut prendre pour $A$ un anneau de
valuation discrète, pour $A'$ un anneau de valuation discrète qui soit un
$A$-module libre de rang 2, et pour $A'_{h'}$ le corps des fractions de $A'$.
\end{remark}

On a toutefois le résultat suivant~:

\begin{corollary}[11.9.9]
\label{IV.11.9.9-fr}
Soient $X$ un préschéma artinien, $g:X'\to X$ un morphisme plat (on notera que
ces deux conditions sont satisfaites si $X$ est le spectre d'un \emph{corps}
et $g$ un morphisme \emph{quelconque}). Alors, avec les notations de (11.9.1)
et (11.9.5), si la famille $(u_\lambda)$ est séparante, il en est de même de la
famille $(u'_\lambda)$.
\end{corollary}

\begin{proof}
On peut évidemment se borner au cas où $X=\operatorname{Spec}(A)$ est le
spectre d'un anneau local artinien $A$ (\textbf I, 6.2.2)~; on note ensuite que
pour tout ouvert affine $U'=\operatorname{Spec}(A')$ de $X'$, $A'$ est un
$A$-module plat, donc \emph{projectif} $(\mathbf 0_{\mathrm{III}},10.1.3)$.
Il suffit donc d'appliquer (11.9.7) à tout ouvert affine de $X'$ pour obtenir
le corollaire.
\end{proof}

\begin{theorem}[11.9.10]
\label{IV.11.9.10-fr}
Soient $X$ un préschéma, $(u_\lambda)$ une famille d'homomorphismes
(11.9.1.1), $g:X'\to X$ un morphisme, $(u'_\lambda)$ la famille
d'homomorphismes déduite de $(u_\lambda)$ par le changement de base $g$
(11.9.5).
\begin{enumerate}[label=(\roman*)]
  \item Si $g$ est un morphisme fidèlement plat et si la famille
  $(u'_\lambda)$ est séparante, alors la famille $(u_\lambda)$ est séparante.
  \oldpage[IV-3]{164}
  \item Supposons que $g$ soit un morphisme plat, et en outre que l'une des
  deux conditions suivantes soit vérifiée~:
  \begin{enumerate}[label=\emph{\alph*)}]
    \item $L$ est fini et les $f_\lambda$ sont quasi-compacts.
    \item Le morphisme $g$ est localement de présentation finie.
  \end{enumerate}
  Alors, si la famille $(u_\lambda)$ est séparante, il en est de même de la
  famille $(u'_\lambda)$.
\end{enumerate}
\end{theorem}

\begin{proof}
\begin{enumerate}[label=(\roman*)]
  \item En vertu de (2.2.8), il suffit de montrer que si une section $t$ de
  $\mathcal F$ au-dessus d'un ouvert $U$ de $X$ appartient au noyau de chacun
  des $u_\lambda$, son image $t'$ par l'homomorphisme canonique
  \[
  \Gamma(\rho):\Gamma(U,\mathcal F)\longrightarrow
  \Gamma(g^{-1}(U),g^*(\mathcal F))
  =\Gamma(U,g_*(g^*(\mathcal F)))
  \]
  est nulle. Or les images de $t'$ par les $g^*(u_\lambda)$ sont les images
  des $u_\lambda(t)$ par l'homomorphisme
  \[
  \Gamma(U,(f_\lambda)_*(\mathcal G_\lambda))\longrightarrow
  \Gamma(g^{-1}(U),g^*((f_\lambda)_*(\mathcal G_\lambda))),
  \]
  donc sont nulles, et \emph{a fortiori} on a $u'_\lambda(t')=0$ pour tout
  $\lambda$, donc $t'=0$ par hypothèse, ce qui prouve (i).

  \item La question étant locale sur $X$ et $X'$, on peut se borner au cas où
  $X=\operatorname{Spec}(A)$ et $X'=\operatorname{Spec}(A')$ sont affines,
  $A'$ étant un $A$-module plat, et à prouver que, pour toute section $z'$ de
  $\mathcal F'$ au-dessus de $X'$ dont les images par tous les $u'_\lambda$
  sont nulles, alors $z'$ est elle-même nulle. On a vu en outre (11.9.6) que
  l'on peut alors supposer que $Z_\lambda=X$ pour tout $\lambda$.

  Distinguons maintenant les deux cas.

  \emph{a)} Si $L$ est fini et les $f_\lambda$ quasi-compacts, on a vu (11.9.3)
  qu'on peut encore se ramener au cas où $Z_\lambda=X$ pour tout
  $\lambda\in L$, et où en outre $L$ est fini. Il revient au même alors de dire
  que l'intersection des noyaux des
  $u_\lambda:\mathcal F\to\mathcal G_\lambda$ est nulle, ou que
  l'homomorphisme
  $u=(u_\lambda):\mathcal F\to\mathcal G=\bigoplus_\lambda\mathcal G_\lambda$
  est injectif. Comme $u'=g^*(u)$ est injectif puisque $g$ est plat, la
  proposition est démontrée dans ce cas.

  \emph{b)} Soient $\mathcal F=\widetilde M$,
  $\mathcal G_\lambda=\widetilde N_\lambda$, et posons
  $M'=M\otimes_A A'$, $N'_\lambda=N_\lambda\otimes_A A'$, de sorte que
  $\mathcal F'=\widetilde{M'}$, $\mathcal G'_\lambda=\widetilde{N'_\lambda}$~;
  par abus de langage, nous noterons encore $u_\lambda$ l'homomorphisme
  $M\to N_\lambda$, et $u'_\lambda$ l'homomorphisme
  $u_\lambda\otimes 1:M'\to N'_\lambda$. Donnons-nous un élément $z'\in M'$
  tel que $u'_\lambda(z')=0$ pour tout $\lambda$~; il s'agit de prouver que
  l'on a $z'=0$. Or, l'hypothèse que $g$ est plat et de présentation finie
  entraîne, d'après (11.3.15), qu'il existe une suite finie
  $(s_i)_{1\leq i\leq n}$ d'éléments de $A$, telle que, si l'on pose
  $\mathfrak J_k=s_1A+\cdots+s_kA$, et
  $A_i=A_{s_i}/\mathfrak J_{i-1}A_{s_i}$, l'anneau
  $A'_i=A'\otimes_A A_i$ soit un $A_i$-module \emph{libre} pour
  $1\leq i\leq n$, et $\mathfrak J_n=A$. La proposition sera établie si nous
  prouvons pour $1\leq i\leq n$ l'assertion suivante~:
  \begin{enumerate}[label=$(*_i)$]
    \item \emph{Il existe un entier $m_i>0$ tel que $s_j^{m_i}z'=0$ pour
    $j\leq i$.}
  \end{enumerate}
  En effet, posant alors $k=m_n$ et notant que les $s_i^k$ ($1\leq i\leq n$)
  engendrent aussi l'idéal unité de $A$, l'assertion $(*_n)$ montrera que $z'$,
  combinaison linéaire des $s_i^kz'$, est nul.

  Prouvons $(*_i)$ par récurrence sur $i$, l'assertion étant vide pour $i=0$.
  Supposons donc $i>0$ et soit $m$ un multiple commun des $m_j$ pour
  $j\leq i-1$. Remarquons que (pour $1\leq h\leq n$) si
  $\mathfrak J'_h$ est l'idéal engendré par les $s_j^m$ ($0\leq j\leq h$),
  $\mathfrak J_h/\mathfrak J'_h$ est nilpotent~; remplacer les $s_j$ par
  $s_j^m$ pour $1\leq j\leq n$ revient donc à remplacer, pour
  $1\leq h\leq n$, $A_h$ par
  $A_{s_h}/\mathfrak J'_{h-1}A_{s_h}=B_h$, de sorte que
  $A_h=B_h/(\mathfrak J_{h-1}/\mathfrak J'_{h-1})B_h$~; comme $A'$ est un
  $A$-module plat, il résulte de $(\mathbf 0_{\mathrm{III}},10.1.2)$ que
  $B'_h=A'\otimes_A B_h$ est encore un $B_h$-module libre. On peut donc
  \oldpage[IV-3]{165}
  remplacer tous les $s_j$ ($1\leq j\leq n$) par $s_j^m$ sans changer les
  propriétés des $\mathfrak J_h$ et des $A'_h$, et supposer par la suite que
  $m=1$. Alors $z'$, étant annulé par $\mathfrak J_{i-1}$, s'identifie à un
  élément de
  $\operatorname{Hom}_{A'}(A'/\mathfrak J_{i-1}A',M')$, et comme
  $\mathfrak J_{i-1}A'$ est un idéal de type fini de $A'$, ce module
  d'homomorphismes s'identifie lui-même à
  $\operatorname{Hom}_A(A/\mathfrak J_{i-1},M)\otimes_A A'$
  $(\mathbf 0_{\mathrm I},6.2.2)$. Soit
  \[
  v_\lambda=\operatorname{Hom}(1,u_\lambda):
  \operatorname{Hom}_A(A/\mathfrak J_{i-1},M)\longrightarrow
  \operatorname{Hom}_A(A/\mathfrak J_{i-1},N_\lambda)
  \]
  l'homomorphisme déduit de $u_\lambda$~; la famille $(v_\lambda)$ est elle
  aussi \emph{séparante}. En effet, pour tout $t\in A$, on a
  \[
  (\operatorname{Hom}_A(A/\mathfrak J_{i-1},M))_t
  =\operatorname{Hom}_{A_t}(A_t/(\mathfrak J_{i-1})_t,M_t)
  \]
  et de même en remplaçant $M$ par $N_\lambda$, puisque l'idéal
  $\mathfrak J_{i-1}$ est de type fini $(\mathbf 0_{\mathrm I},1.3.5)$~;
  comme par hypothèse l'intersection des noyaux des $(u_\lambda)_t$ est nulle,
  il en est de même de l'intersection des noyaux des
  $(v_\lambda)_t=\operatorname{Hom}(1,(u_\lambda)_t)$, d'où notre assertion
  (11.9.4). Remplaçant $A$ par $A/\mathfrak J_{i-1}$, $M$ par
  $\operatorname{Hom}_A(A/\mathfrak J_{i-1},M)$, $N_\lambda$ par
  $\operatorname{Hom}_A(A/\mathfrak J_{i-1},N_\lambda)$ (qui sont des
  $(A/\mathfrak J_{i-1})$-modules), $u_\lambda$ par $v_\lambda$ et enfin $A'$
  par $A'/\mathfrak J_{i-1}A'$, on voit qu'on peut se ramener au cas où, dans
  la situation initiale, l'élément $s=s_i\in A$ est tel que $A'_s$ soit un
  $A_s$-module libre. Or la famille des
  $(u_\lambda)_s:M_s\to(N_\lambda)_s$ est séparante par hypothèse~; comme on a
  $(u'_\lambda)_s(z'/1)=0$, il résulte de (11.9.7) que l'on a $z'/1=0$ dans
  $M'_s$~; mais cela signifie qu'il existe un entier $r$ tel que $s^rz'=0$
  dans $M'$, ce qui achève de prouver $(*_i)$ par récurrence.
\end{enumerate}
\end{proof}

\begin{remark}[11.9.11]
\label{IV.11.9.11-fr}
Bornons-nous pour simplifier au cas où $Z_\lambda=X$ pour tout $\lambda$. Il
faut noter alors que si la famille des homomorphismes
$u_\lambda:\mathcal F\to\mathcal G_\lambda$ est séparante, \emph{il ne
s'ensuit pas nécessairement} que, pour tout $x\in X$, l'intersection des
noyaux des homomorphismes
$(u_\lambda)_x:\mathcal F_x\to(\mathcal G_\lambda)_x$ soit réduite à 0. Par
exemple, soient $X$ un préschéma de Jacobson localement noethérien, de
dimension $\geq1$, et $\mathcal F$ un $\mathcal O_X$-Module cohérent~; pour
tout point fermé $x$ de $X$, et tout entier $n\geq0$,
$\mathfrak m_x^n\mathcal F$ est un $\mathcal O_X$-Module cohérent de support
contenu dans $\{x\}$. La famille des homomorphismes canoniques
$\mathcal F\to\mathcal F/\mathfrak m_x^n\mathcal F$ (où $x$ parcourt
l'ensemble $X_0$ des points fermés de $X$ et $n$ l'ensemble des entiers
$\geq0$) est \emph{séparante}~: en effet, si $t$ est une section de
$\mathcal F$ au-dessus d'un ouvert $U$ de $X$ dont les images dans les
$\Gamma(U,\mathcal F/\mathfrak m_x^n\mathcal F)$ sont toutes nulles, il en
résulte aussitôt que pour tout point fermé $x\in U$, on a $t_x=0$~; comme
l'ensemble des points fermés contenus dans $U$ est très dense dans $U$, cela
entraîne bien $t=0$ (10.2.1). Cependant, si l'on prend
$\mathcal F=\mathcal O_X$, et si $y\in X$ est un point \emph{non fermé} de
$X$, on a $(\mathcal O_X/\mathfrak m_x^n\mathcal O_X)_y=0$ pour tout point
fermé $x$ de $X$, mais $\mathcal O_{X,y}\neq0$, et l'intersection des noyaux
des homomorphismes
$\mathcal O_{X,y}\to(\mathcal O_X/\mathfrak m_x^n\mathcal O_X)_y$ est égale
à $\mathcal O_{X,y}$.
\end{remark}

\begin{lemma}[11.9.12]
\label{IV.11.9.12-fr}
Les notations étant celles de (11.9.1) et (11.9.5), supposons la famille
$(u_\lambda)$ séparante~; supposons en outre que $X$ soit un $S$-préschéma, où
$S=\operatorname{Spec}(A)$ est un schéma affine, et que
$X'=X\times_S S'$, où $S'=\operatorname{Spec}(A')$, $A'$ étant une
$A$-algèbre~; supposons enfin que les $\mathcal G_\lambda$ soient $S$-plats.
Soit $(A'_\alpha)_{\alpha\in I}$ la famille filtrante des sous-$A$-algèbres de
type fini de $A'$~; pour tout $\alpha\in I$, posons
$S''_\alpha=\operatorname{Spec}(A'_\alpha)$,
$X''_\alpha=X\times_S S''_\alpha$,
$Z''_{\alpha\lambda}=Z_\lambda\times_X X''_\alpha$, et soient
$f''_{\alpha\lambda}:Z''_{\alpha\lambda}\to X''_\alpha$,
$\mathcal F''_\alpha$, $\mathcal G''_{\alpha\lambda}$ et
$u''_{\alpha\lambda}$ les morphismes, Modules et homomorphismes de Modules
déduits de $f_\lambda$, $\mathcal F$, $\mathcal G_\lambda$ et $u_\lambda$ au
moyen du changement de base $X''_\alpha\to X$. Alors, si, pour tout
$\alpha\in I$, la famille $(u''_{\alpha\lambda})_{\lambda\in L}$ est
séparante, il en est de même de $(u'_\lambda)_{\lambda\in L}$.
\end{lemma}

\begin{proof}
Il s'agit de prouver que si une section $t'$ de $\mathcal F'$ au-dessus d'un
ouvert affine $U'$ de $X'$ est telle que $u'_\lambda(t')=0$ pour tout
$\lambda\in L$, on a $t'=0$. Si $h_\alpha:X'\to X''_\alpha$ est la projection
\oldpage[IV-3]{166}
canonique, il résulte de (8.2.11) qu'il existe un indice $\alpha$ et un ouvert
quasi-compact $U''_\alpha\subset X''_\alpha$ tels que
$U'=h_\alpha^{-1}(U''_\alpha)$~; en outre, en vertu de (8.5.2, (i)), on peut
supposer qu'il y a une section $t''_\alpha$ de $\mathcal F''_\alpha$ au-dessus
de $U''_\alpha$ telle que $t'$ soit l'image canonique de $t''_\alpha$. Quitte à
remplacer $S$, $X$, $Z_\lambda$, $f_\lambda$, $\mathcal F$,
$\mathcal G_\lambda$ et $u_\lambda$ par $S''_\alpha$, $U''_\alpha$,
$(f''_{\alpha\lambda})^{-1}(U''_\alpha)$ et les restrictions correspondantes
de $\mathcal F''_\alpha$, $\mathcal G''_{\alpha\lambda}$ et
$u''_{\alpha\lambda}$, on peut donc supposer que $U'=X'$, que $t'$ est
l'image canonique d'une section $t$ de $\mathcal F$ au-dessus de $X$ et que
l'homomorphisme $A\to A'$ est injectif, ou encore, si $p:S'\to S$ est le
morphisme correspondant, que l'homomorphisme
$\mathcal O_S\to p_*(\mathcal O_{S'})$ intervenant dans la définition de $p$
est injectif. Il en résulte aussitôt en vertu de la platitude de
$\mathcal G_\lambda$ sur $S$ (et en se ramenant par exemple au cas où
$Z_\lambda$ est affine sur $S$ (\textbf I, 1.6.3 et 1.6.5)) que
l'homomorphisme canonique
$\rho:\mathcal G_\lambda\to(g'_\lambda)_*(\mathcal G'_\lambda)$ est lui aussi
injectif. Mais l'homomorphisme composé
\[
\Gamma(X,\mathcal F)\xrightarrow{\Gamma(u_\lambda)}
\Gamma(Z_\lambda,\mathcal G_\lambda)\xrightarrow{\Gamma(\rho)}
\Gamma(Z'_\lambda,\mathcal G'_\lambda)
\]
est égal par définition à l'homomorphisme composé
\[
\Gamma(X,\mathcal F)\xrightarrow{\Gamma(\rho)}
\Gamma(X',\mathcal F')\xrightarrow{\Gamma(u'_\lambda)}
\Gamma(Z'_\lambda,\mathcal G'_\lambda)
\]
donc l'image de $t$ par ces homomorphismes composés est
$u'_\lambda(t')=0$~; en vertu de l'injectivité de l'homomorphisme
$\Gamma(Z_\lambda,\mathcal G_\lambda)\xrightarrow{\Gamma(\rho)}
\Gamma(Z'_\lambda,\mathcal G'_\lambda)$ on en conclut que $u_\lambda(t)=0$
pour tout $\lambda\in L$, d'où $t=0$ par hypothèse, et finalement $t'=0$.
\end{proof}

\begin{proposition}[11.9.13]
\label{IV.11.9.13-fr}
Les notations étant celles de (11.9.1) et (11.9.5), supposons que $X$ soit un
préschéma sur un corps $k$, et que, en posant $S=\operatorname{Spec}(k)$, on
ait $X'=X\times_S S'$, où $S'$ est un $k$-préschéma quelconque. Alors, si la
famille $(u_\lambda)_{\lambda\in L}$ est séparante, il en est de même de
$(u'_\lambda)$.
\end{proposition}

\begin{proof}
On peut se borner au cas où $S'=\operatorname{Spec}(A')$ est affine. Si $A'$
est une $k$-algèbre de type fini, le morphisme $g:X'\to X$ est plat et de
présentation finie, et on est donc dans les conditions d'application de
(11.9.10, (ii), b)). Dans le cas général, on considère $A'$ comme limite
inductive de ses $k$-sous-algèbres $A'_\alpha$ de type fini, et on applique à
chaque $A'_\alpha$ le résultat de (11.9.10, (ii), b))~; on conclut alors à
l'aide du lemme (11.9.12), puisque les $\mathcal G_\lambda$ sont $S$-plats.
\end{proof}

\begin{env}[11.9.14]
\label{IV.11.9.14-fr}
Gardons toujours les notations de (11.9.1) et (11.9.5) et supposons que $X$
soit un $S$-préschéma. Si pour tout changement de base
$g:X\times_S S'\to X$, où $S'\to S$ est un morphisme quelconque, la famille
$(u'_\lambda)$ correspondante est séparante, nous dirons que la famille
$(u_\lambda)$ est \emph{universellement séparante relativement à} $S$.
Lorsque la famille $(u_\lambda)$ est réduite à un seul élément $u$, nous
dirons encore que $u$ est \emph{universellement injectif}, relativement à
$S$. Il est clair alors que pour tout morphisme $h:S'\to S$, la famille
$(u'_\lambda)$ correspondante est universellement séparante relativement à
$S'$~; inversement, si $h$ est \emph{fidèlement plat} et si $(u'_\lambda)$ est
universellement séparante relativement à $S'$, alors $(u_\lambda)$ est
universellement séparante relativement à $S$, comme il résulte aussitôt de
(11.9.10, (i)) et du fait que pour tout morphisme $S''\to S$, le morphisme
correspondant $S''\times_S S'\to S''$ est fidèlement plat.
\end{env}

\oldpage[IV-3]{167}

\begin{proposition}[11.9.15]
\label{IV.11.9.15-fr}
Les notations étant celles de (11.9.1), supposons que $X$ soit un
$S$-préschéma, les $\mathcal G_\lambda$ étant $S$-plats. Soit $S_0$ un
sous-préschéma fermé de $S$ défini par un Idéal quasi-cohérent nilpotent
$\mathcal J$ de $\mathcal O_S$, tel que les
$(\mathcal O_S/\mathcal J)$-Modules $\mathcal J^k/\mathcal J^{k+1}$ soient
tous localement libres. Soit $(u_{\lambda0})$ la famille d'homomorphismes
obtenue à partir du changement de base
$X\leftarrow X_0=X\times_S S_0$. Alors, si la famille $(u_{\lambda0})$ est
séparante (resp. universellement séparante relativement à $S_0$), la famille
$(u_\lambda)$ est séparante (resp. universellement séparante relativement à
$S$).
\end{proposition}

\begin{proof}
Notons que si $S'\to S$ est un changement de base quelconque et
$S'_0=S'\times_S S_0$, $S'_0$ est un sous-préschéma fermé de $S'$ défini par
un idéal quasi-cohérent nilpotent $\mathcal J'$ de $\mathcal O_{S'}$, tel que
$\mathcal J'^k/\mathcal J'^{k+1}$ soit un
$(\mathcal O_{S'}/\mathcal J')$-Module localement libre pour tout $k$
(2.1.8, (i))~; comme en outre les
$\mathcal G'_\lambda=\mathcal G_\lambda\otimes_S S'$ sont $S'$-plats, on voit
que l'assertion relative aux familles universellement séparantes est
conséquence de l'assertion relative aux familles séparantes. Pour prouver
cette dernière, on peut (11.9.3) se ramener au cas où
$S=\operatorname{Spec}(A)$, $X=\operatorname{Spec}(B)$ sont affines,
$Z_\lambda=X$ pour tout $\lambda$, $\mathcal F=\widetilde M$,
$\mathcal G_\lambda=\widetilde N_\lambda$, où $M$ et les $N_\lambda$ sont des
$B$-modules, les $N_\lambda$ étant des $A$-modules plats. En outre, la question
étant locale sur $X$ et $S$, il suffit de voir que si $t\in M$ est tel que
$u_\lambda(t)=0$ pour tout $\lambda$, alors $t=0$. On a
$\mathcal J=\widetilde{\mathfrak J}$, où $\mathfrak J$ est un idéal nilpotent
de $A$, tel que les $\mathfrak J^k/\mathfrak J^{k+1}$ soient des
$(A/\mathfrak J)$-modules libres, et par hypothèse les
$u_{\lambda0}:M/\mathfrak JM\to N_\lambda/\mathfrak JN_\lambda$ forment une
famille séparante. Supposons que $\mathfrak J^{n+1}=0$ ($n$ entier $\geq0$) et
raisonnons par récurrence sur $n$, l'assertion étant triviale pour $n=0$. Si
$\bar t$ est la classe de $t$ dans $M/\mathfrak J^nM$, la classe de
$u_\lambda(t)$ dans $N_\lambda/\mathfrak J^nN_\lambda$ est nulle pour tout
$\lambda$, donc, par l'hypothèse de récurrence, $\bar t=0$, autrement dit on a
$t\in\mathfrak J^nM$. Or
$\mathfrak J^n=\mathfrak J^n/\mathfrak J^{n+1}$ est un
$(A/\mathfrak J)$-module libre~; si $(e_\alpha)$ est une base de ce module, on
peut donc écrire $t=\sum_\alpha e_\alpha t^0_\alpha$, avec
$t^0_\alpha\in M/\mathfrak JM$, nul sauf pour un nombre fini d'indices.
D'autre part, puisque $N_\lambda$ est un $A$-module plat,
$\mathfrak J^nN_\lambda$ s'identifie à
$\mathfrak J^n\otimes_{(A/\mathfrak J)}(N_\lambda/\mathfrak JN_\lambda)$ et
l'on peut par suite écrire
\[
u_\lambda(t)=\sum_\alpha e_\alpha u_{\lambda0}(t^0_\alpha)
=\sum_\alpha e_\alpha\otimes u_{\lambda0}(t^0_\alpha).
\]
Comme par hypothèse $u_\lambda(t)=0$, on en déduit que
$u_{\lambda0}(t^0_\alpha)=0$ pour tout $\alpha$ et tout $\lambda$~; d'où
$t^0_\alpha=0$ pour tout $\alpha$ puisque la famille $(u_{\lambda0})$ est
séparante, et par suite $t=0$. C.Q.F.D.
\end{proof}

\begin{theorem}[11.9.16]
\label{IV.11.9.16-fr}
Les notations étant celles de (11.9.1), supposons que $X$ soit un
$S$-préschéma localement noethérien, $\mathcal F$ un $\mathcal O_X$-Module
cohérent et que les $\mathcal G_\lambda$ soient $S$-plats. Pour tout $s\in S$,
soit $((u_\lambda)_s)_{\lambda\in L}$ la famille obtenue à partir de
$(u_\lambda)$ par le changement de base
$X_s=X\times_S\operatorname{Spec}(k(s))\to X$. Alors, pour que la famille
$(u_\lambda)$ soit universellement séparante relativement à $S$, il faut et il
suffit que pour tout $s\in S$, la famille $((u_\lambda)_s)$ soit séparante.
\end{theorem}

\begin{proof}
La nécessité de la condition découle trivialement des définitions.
Inversement, supposons la condition de l'énoncé vérifiée, et prouvons d'abord
que la famille $(u_\lambda)$ est séparante. On peut (11.9.3) se réduire au cas
où $S=\operatorname{Spec}(A)$, $X=\operatorname{Spec}(B)$ sont affines,
$Z_\lambda=X$ pour tout $\lambda$, $\mathcal F=\widetilde M$,
$\mathcal G_\lambda=\widetilde N_\lambda$, où $B$ est noethérien, $M$ est un
$B$-module de type fini et les $N_\lambda$ des $A$-modules plats et se borner à
prouver que, si $t\in M$ est tel que $u_\lambda(t)=0$ pour tout $\lambda$,
alors $t=0$. Pour montrer que $t=0$, il suffit de prouver que pour tout idéal
maximal $\mathfrak p$ de $B$, l'image $t_\mathfrak p$ de $t$ dans
$M_\mathfrak p$ est nulle (Bourbaki, \emph{Alg. comm.}, chap. II, \S~3,
n\textsuperscript{o}~3, cor. 1 du th. 1). On peut donc se borner à montrer que
l'intersection
\oldpage[IV-3]{168}
des noyaux des
$(u_\lambda)_\mathfrak p:M_\mathfrak p\to(N_\lambda)_\mathfrak p$ déduits
des $(u_\lambda)$ par le changement de base
$\operatorname{Spec}(B_\mathfrak p)\to\operatorname{Spec}(B)$ est réduite à
0. Autrement dit, on est ramené au cas où $B$ est un anneau local noethérien,
et en considérant l'idéal premier de $A$ image réciproque de l'idéal maximal de
$B$, on peut aussi supposer que $A$ est un anneau local, d'idéal maximal
$\mathfrak m$. Alors, comme $\mathfrak mB$ est contenu dans l'idéal maximal de
$B$, et que $M$ est un $B$-module de type fini, l'intersection des
$\mathfrak m^nM$ est réduite à 0 $(\mathbf 0_{\mathrm I},7.3.5)$, donc il
suffit de prouver que pour tout $n$, l'image de $t$ dans
$M/\mathfrak m^{n+1}M$ est nulle. Il suffit donc de prouver que la famille
déduite de $(u_\lambda)$ par le changement de base
$\operatorname{Spec}(B/\mathfrak m^{n+1}B)\to\operatorname{Spec}(B)$ est
séparante, ce qui signifie encore qu'on peut se borner au cas où $A$ est un
anneau local dont l'idéal maximal $\mathfrak m$ est nilpotent. Mais alors les
$\mathfrak m^k/\mathfrak m^{k+1}$ sont des $(A/\mathfrak m)$-modules libres,
et en vertu de l'hypothèse sur les $(u_\lambda)_s$, on est précisément dans les
conditions d'application de (11.9.15), d'où la conclusion annoncée.

Soient maintenant $h:S'\to S$ un morphisme changement de base, et
$(u'_\lambda)$ la famille obtenue à partir de $(u_\lambda)$ par le changement
de base $h':X\times_S S'\to X$~; prouvons que $(u'_\lambda)$ est aussi
séparante. Supposons d'abord que $h$ soit \emph{localement de type fini}~; il
en est alors de même de $h'$, donc $X'=X\times_S S'$ est localement
noethérien~; de plus, si $s'\in S'$ est au-dessus du point $s\in S$, il résulte
de (11.9.13) appliqué à $X_s$ et à
$X_{s'}=X_s\otimes_{k(s)}k(s')$ que, pour tout $s'\in S'$, la famille
$((u'_\lambda)_{s'})$ est séparante~; on peut par suite conclure de la première
partie de la démonstration que dans ce cas $(u'_\lambda)$ est séparante.

Enfin, si $h$ est quelconque, on peut évidemment se limiter au cas où
$S=\operatorname{Spec}(A)$ et $S'=\operatorname{Spec}(A')$ sont affines, et
considérer $A'$ comme limite inductive de ses sous-$A$-algèbres de type fini.
Comme les $\mathcal G_\lambda$ sont $S$-plats, il suffit d'appliquer ce qui
précède et le lemme (11.9.5) pour terminer la démonstration.
\end{proof}

\begin{proposition}[11.9.17]
\label{IV.11.9.17-fr}
Soient $f:X\to S$ un morphisme localement de présentation finie, $\mathcal F$
un $\mathcal O_X$-Module quasi-cohérent de présentation finie et $f$-plat,
$U$ un ouvert de $X$, $j:U\to X$ l'injection canonique,
$u:\mathcal F\to j_*(j^*(\mathcal F))$ l'homomorphisme canonique
$(\mathbf 0_{\mathrm I},4.4.3.2)$. Pour tout $s\in S$, soient $X_s$ la fibre
$X\times_S\operatorname{Spec}(k(s))$, $U_s$ l'ouvert $U\cap X_s$ de $X_s$,
$\mathcal F_s=\mathcal F\otimes_{\mathcal O_S}k(s)$,
$j_s:U_s\to X_s$ l'injection canonique,
$u_s:\mathcal F_s\to(j_s)_*((j_s)^*(\mathcal F_s))$ l'homomorphisme
canonique. Pour que $u$ soit universellement injectif relativement à $S$, il
faut et il suffit que $u_s$ soit injectif pour tout $s\in S$.
\end{proposition}

\begin{proof}
Il n'y a encore à prouver que la suffisance de la condition. Lorsque $X$ est
localement noethérien, la proposition est un corollaire immédiat de (11.9.16).
Nous allons nous ramener à ce cas en deux étapes, en nous bornant, comme on
peut évidemment le faire, au cas où $S=\operatorname{Spec}(A)$ et $X$ sont
affines.

\emph{A) Cas où $U$ est quasi-compact.} --- Nous utiliserons le lemme suivant~:

\begin{lemma}[11.9.17.1]
\label{IV.11.9.17.1-fr}
Sous les hypothèses générales de (11.9.17), et en supposant en outre $S$ et
$X$ affines et $U$ quasi-compact, l'ensemble $E$ des $s\in S$ tels que $u_s$
soit injectif est constructible.
\end{lemma}

\begin{proof}
En effet, les fibres $X_s$ sont des préschémas localement noethériens, donc
$E$ peut aussi, en vertu de (5.10.2), être défini comme l'ensemble des
$s\in S$ tels que $\operatorname{Ass}(\mathcal F_s)\subset U_s$. Notons en
outre que $U$, étant quasi-compact dans un schéma affine, est constructible.
Alors, la vérification de la condition (9.2.1, (i)) découle aussitôt de
(4.2.7)~; d'autre
\oldpage[IV-3]{169}
part, la vérification de (9.2.1, (ii)) se fait aisément en utilisant l'étude des
cycles premiers associés au voisinage de la fibre générique (9.8.3), ainsi que
(9.5.2) et (9.5.3).
\end{proof}

Ce lemme étant établi, on peut, en vertu de (8.9.1) et (8.2.11), supposer qu'il
existe un sous-anneau noethérien $A_0$ de $A$, un préschéma de type fini $X_0$
sur $S_0=\operatorname{Spec}(A_0)$, un ouvert $U_0$ de $X_0$ et un
$\mathcal O_{X_0}$-Module cohérent $\mathcal F_0$ tels que
$X=X_0\times_{S_0}S$ et que, si $p_0:X\to X_0$ est la projection canonique,
on ait $U=p_0^{-1}(U_0)$ et $\mathcal F=p_0^*(\mathcal F_0)$. Soit
$(A_\alpha)$ la famille filtrante des sous-anneaux de $A$ qui sont des
$A_0$-algèbres de type fini, et posons
$S_\alpha=\operatorname{Spec}(A_\alpha)$,
$X_\alpha=X_0\times_{S_0}S_\alpha$,
$U_\alpha=U_0\times_{S_0}S_\alpha$,
$\mathcal F_\alpha=\mathcal F_0\otimes_{S_0}S_\alpha$~; soit $u_\alpha$
l'homomorphisme canonique relatif à $\mathcal F_\alpha$ et $U_\alpha$, défini
comme dans (11.9.17). Pour tout $s\in S$, l'hypothèse que $u_s$ soit injectif
entraîne que $(u_\alpha)_{s_\alpha}$ l'est aussi (11.9.10, (i)), où
$s_\alpha=q_\alpha(s)$ est l'image de $s$ par le morphisme
$q_\alpha:S\to S_\alpha$. Si $E_\alpha$ est l'ensemble des
$s_\alpha\in S_\alpha$ tels que $(u_\alpha)_{s_\alpha}$ soit injectif, on a
donc $S=q_\alpha^{-1}(E_\alpha)$, et les $E_\alpha$ forment un système
projectif d'ensembles. Mais le lemme (11.9.17.1) appliqué à $u_\alpha$ montre
que $E_\alpha$ est constructible dans $S_\alpha$~; on déduit donc de (8.3.4)
qu'il existe un indice $\alpha$ tel que $E_\alpha=S_\alpha$. Mais alors, comme
$X_\alpha$ est noethérien, $u_\alpha$ est universellement injectif en vertu de
(11.9.16), donc il en est de même de $u$.

\emph{B) Cas général.} --- L'ouvert $U$ est réunion filtrante croissante
d'ouverts quasi-compacts $U_\lambda$~; si $j_\lambda:U_\lambda\to X$ est
l'injection canonique et
$u_\lambda:\mathcal F\to(j_\lambda)_*(j_\lambda^*(\mathcal F))$
l'homomorphisme correspondant, il résulte du lemme (11.9.17.1) que l'ensemble
$E_\lambda$ des $s\in S$ tels que $(u_\lambda)_s$ soit injectif est
constructible. D'autre part, pour un $s\in S$, dire que $u_s$ (resp.
$(u_\lambda)_s$) est injectif signifie que
$\operatorname{Ass}(\mathcal F_s)\subset U\cap X_s$ (resp.
$\operatorname{Ass}(\mathcal F_s)\subset U_\lambda\cap X_s$) (5.10.2).
Comme $\operatorname{Ass}(\mathcal F_s)$ est fini (3.1.6), dire que $u_s$ est
injectif signifie donc qu'il existe $\lambda$ tel que $s\in E_\lambda$~;
l'hypothèse signifie par suite que $S=\bigcup_\lambda E_\lambda$. En vertu de
(1.9.9), il existe un indice $\lambda$ tel que $S=E_\lambda$, d'où l'on
conclut par la première partie du raisonnement que $u_\lambda$ est
universellement injectif. Il résulte alors de la factorisation de $u_\lambda$~:
\[
\mathcal F\xrightarrow{u}j_*(\mathcal F|U)\longrightarrow
(j_\lambda)_*(\mathcal F|U_\lambda)
\]
que $u$ est aussi universellement injectif.
\end{proof}

\begin{remark}[11.9.18]
\label{IV.11.9.18-fr}
Soient $X$ un préschéma, $\mathcal F$, $\mathcal G$ deux
$\mathcal O_X$-Modules quasi-cohérents de présentation finie, $\mathcal G$
étant en outre supposé \emph{localement libre},
$u:\mathcal F\to\mathcal G$ un homomorphisme. Les conditions suivantes sont
équivalentes~:
\begin{enumerate}[label=\emph{\alph*)}]
  \item pour tout morphisme $g:X'\to X$, l'homomorphisme
  $g^*(u):g^*(\mathcal F)\to g^*(\mathcal G)$ est injectif~;
  \item pour tout $x\in X$, l'homomorphisme
  $u\otimes 1_{k(x)}:\mathcal F\otimes_{\mathcal O_X}k(x)
  \to\mathcal G\otimes_{\mathcal O_X}k(x)$ est injectif~;
  \item pour tout $x\in X$, il existe un voisinage ouvert $U$ de $x$ tel que
  $u|U:\mathcal F|U\to\mathcal G|U$ soit un isomorphisme de $\mathcal F|U$
  sur un facteur direct de $\mathcal G|U$.
\end{enumerate}

En effet, il est clair que c) entraîne a) et que a) entraîne b). Le fait que
b) entraîne c) résulte de $(\mathbf 0,19.1.12)$,
$(\mathbf 0_{\mathrm I},5.2.5)$ et $(\mathbf 0_{\mathrm I},5.5.5)$.

Lorsque les conditions équivalentes précédentes sont vérifiées, on dit que
$u$ est \emph{universellement injectif}.
\end{remark}
\oldpage[IV-3]{170}
\subsection{Familles schématiquement dominantes de morphismes et familles
schématiquement denses de sous-préschémas}
\label{subsection:IV.11.10-fr}

\begin{proposition}[11.10.1]
\label{IV.11.10.1-fr}
Soient $X$ un préschéma, $(Z_\lambda)_{\lambda\in L}$ une famille de
préschémas, et pour tout $\lambda\in L$, soit
$f_\lambda=(\psi_\lambda,\theta_\lambda):Z_\lambda\to X$ un morphisme. Les
conditions suivantes sont équivalentes~:
\begin{enumerate}[label=\emph{\alph*)}]
  \item La famille des homomorphismes
  $\theta_\lambda:\mathcal O_X\to(f_\lambda)_*(\mathcal O_{Z_\lambda})$ est
  séparante (autrement dit (11.9.1), l'intersection des noyaux des
  $\theta_\lambda$ est nulle).
  \item Pour tout ouvert $U$ de $X$, toute section $t$ de $\mathcal O_X$
  au-dessus de $U$ dont toutes les images par les homomorphismes canoniques
  \[
  (\theta_\lambda)_U:\Gamma(U,\mathcal O_X)
  \longrightarrow\Gamma(f_\lambda^{-1}(U),\mathcal O_{Z_\lambda})
  \tag{11.10.1.1}
  \label{IV.11.10.1.1-fr}
  \]
  sont nulles, est elle-même nulle.
  \item Pour tout ouvert $U$ de $X$, et tout sous-préschéma fermé $Y$ de $U$
  tel que pour tout $\lambda\in L$, il existe une factorisation
  \[
  f_\lambda^{-1}(U)\xrightarrow{g_\lambda}Y\xrightarrow{j}U
  \tag{11.10.1.2}
  \label{IV.11.10.1.2-fr}
  \]
  de la restriction $f_\lambda^{-1}(U)\to U$ de $f_\lambda$ (où $j$ est
  l'injection canonique), on a $Y=U$.
\end{enumerate}

Si de plus $X$ est un $S$-préschéma, ces conditions sont aussi équivalentes à
la suivante~:
\begin{enumerate}[label=\emph{\alph*)}]
\setcounter{enumi}{3}
  \item Pour tout morphisme séparé $g:X'\to S$ et tout couple de
  $S$-morphismes $u$, $v$ d'un ouvert $U$ de $X$ dans $X'$, tel que pour tout
  $\lambda$, les composés de $u$ et $v$ avec le morphisme
  $f_\lambda^{-1}(U)\to U$, restriction de $f_\lambda$, soient égaux, on a
  $u=v$.
\end{enumerate}
\end{proposition}

\begin{proof}
L'équivalence de a) et b) résulte des définitions. Pour voir que b) entraîne
c), il suffit de considérer l'Idéal quasi-cohérent $\mathcal J$ de
$\mathcal O_U$ définissant $Y$, et de noter que, pour tout ouvert $V\subset U$,
l'hypothèse entraîne que le morphisme $(\theta_\lambda)_V$ se factorise en
\[
\Gamma(V,\mathcal O_U)\longrightarrow
\Gamma(V,\mathcal O_U/\mathcal J)\longrightarrow
\Gamma(f_\lambda^{-1}(V),\mathcal O_{Z_\lambda}).
\]
On en conclut que toute section $t$ de $\mathcal J$ au-dessus de $V$ a pour
image 0 dans tous les $\Gamma(f_\lambda^{-1}(V),\mathcal O_{Z_\lambda})$,
donc, en vertu de b), $\mathcal J=0$ et $Y=U$. Inversement, si c) est vérifiée,
il suffit, pour prouver b), d'appliquer c) au sous-préschéma fermé $Y$ de $U$
défini par l'Idéal $t\mathcal O_U$~: l'hypothèse que les images de $t$ par les
$(\theta_\lambda)_U$ sont toutes nulles entraîne que l'on a des factorisations
(11.10.1.2) pour tout $\lambda$ ($\mathbf I$, 4.1.9). Pour prouver que c)
entraîne d), il suffit d'appliquer c) au sous-préschéma fermé $Y$ de $U$, image
réciproque de la diagonale de $X'\times_SX'$ par $(u,v)_S$ et d'utiliser
($\mathbf I$, 4.4.1). Inversement, on déduit b) de d) en considérant le
$S$-schéma $X'=S[T]=S\otimes_{\mathbf Z}\mathbf Z[T]$ ($T$ indéterminée) et en
se rappelant que les sections de $\mathcal O_U$ au-dessus de $U$ correspondent
biunivoquement aux $S$-morphismes $U\to S[T]$ ($\mathbf I$, 3.3.15)~; dire que
deux sections de $\mathcal O_U$ au-dessus de $U$ ont mêmes images par tous les
$(\theta_\lambda)_U$ équivaut à dire que les composés des deux morphismes
correspondants avec tous les morphismes $f_\lambda^{-1}(U)\to U$ sont égaux.
\end{proof}

\oldpage[IV-3]{171}
\begin{definition}[11.10.2]
\label{IV.11.10.2-fr}
Lorsque les conditions équivalentes de (11.10.1) sont vérifiées, on dit que la
famille $(f_\lambda)$ est \emph{schématiquement dominante}. Lorsque les
$f_\lambda$ sont les immersions canoniques dans $X$ d'une famille
$(Z_\lambda)$ de sous-préschémas de $X$, on dit aussi que la famille
$(Z_\lambda)$ est \emph{schématiquement dense}.
\end{definition}

\begin{remarks}[11.10.3]
\label{IV.11.10.3-fr}
\begin{enumerate}[label=(\roman*)]
  \item La notion de famille schématiquement dominante est \emph{locale} sur
  $X$, comme il résulte par exemple de la forme b) dans (11.10.1)~: si
  $(W_\alpha)$ est un recouvrement ouvert de $X$, la famille $(f_\lambda)$ est
  schématiquement dominante si et seulement s'il en est ainsi de chacune des
  familles formées des morphismes
  $f_\lambda^{-1}(W_\alpha)\to W_\alpha$, restrictions des $f_\lambda$.
  \item Si $Z$ est le préschéma somme des $Z_\lambda$, $f:Z\to X$ le morphisme
  coïncidant avec $f_\lambda$ dans chacun des $Z_\lambda$, il revient au même
  de dire que la famille $(f_\lambda)$ est schématiquement dominante, ou que
  la famille réduite au seul élément $f$ l'est.
  \item Soit $M$ un second ensemble d'indices et, pour tout $\lambda\in L$,
  soit $(g_{\lambda\mu})_{\mu\in M}$ une famille de morphismes
  $g_{\lambda\mu}:Z'_{\lambda\mu}\to Z_\lambda$~; si, pour chaque
  $\lambda\in L$, la famille $(g_{\lambda\mu})_{\mu\in M}$ est
  schématiquement dominante, alors, il revient au même de dire que la famille
  $(f_\lambda)_{\lambda\in L}$ est schématiquement dominante, ou que la
  famille
  $(f_\lambda\circ g_{\lambda\mu})_{(\lambda,\mu)\in L\times M}$ l'est
  (11.9.2).
  \item Soit $f:Z\to X$ un morphisme tel que $f_*(\mathcal O_Z)$ soit un
  $\mathcal O_X$-Module quasi-cohérent (par exemple un morphisme quasi-compact
  et quasi-séparé (1.7.4)). Alors, dire que $f$ est schématiquement dominant
  signifie que \emph{l'image fermée de $Z$ par $f$} ($\mathbf I$, 9.5.3) est
  \emph{identique à $X$}.
\end{enumerate}
\end{remarks}

\begin{proposition}[11.10.4]
\label{IV.11.10.4-fr}
Si la famille de morphismes $f_\lambda:Z_\lambda\to X$ est schématiquement
dominante, la réunion des $f_\lambda(Z_\lambda)$ est dense dans $X$.
Inversement, si cette réunion est dense dans $X$ et si de plus $X$ est réduit,
la famille $(f_\lambda)$ est schématiquement dominante.
\end{proposition}

\begin{proof}
La première assertion résulte aussitôt de (11.10.1, b)). D'autre part, si $X$
est réduit, et si la réunion des $f_\lambda(Z_\lambda)$ est dense dans $X$,
alors, si l'on a des factorisations (11.10.1.2) pour tout $\lambda\in L$, on a
aussi des factorisations
\[
(f_\lambda^{-1}(U))_{\mathrm{red}}
\xrightarrow{(g_\lambda)_{\mathrm{red}}}Y_{\mathrm{red}}
\xrightarrow{j_{\mathrm{red}}}U,
\]
et l'hypothèse entraîne que l'espace sous-jacent à $Y$ est identique à $U$,
donc $Y=Y_{\mathrm{red}}=U$ puisque $U$ est réduit.
\end{proof}

Les résultats du n$^\circ$~11.9 sur les familles séparantes se traduisent en
résultats sur les familles schématiquement dominantes~:

\begin{theorem}[11.10.5]
\label{IV.11.10.5-fr}
Soient $(f_\lambda)$ une famille de morphismes
$f_\lambda:Z_\lambda\to X$, $g:X'\to X$ un morphisme, et posons
$Z'_\lambda=Z_\lambda\times_XX'$,
$f'_\lambda=(f_\lambda)_{(X')}:Z'_\lambda\to X'$.
\begin{enumerate}[label=(\roman*)]
  \item Si $g$ est fidèlement plat et si la famille $(f'_\lambda)$ est
  schématiquement dominante, alors la famille $(f_\lambda)$ est
  schématiquement dominante.
  \item Supposons que $g$ soit un morphisme plat et en outre que l'une des
  conditions suivantes soit vérifiée~:
  \begin{enumerate}[label=\emph{\alph*)}]
    \item $L$ est fini et les $f_\lambda$ sont quasi-compacts.
    \item Le morphisme $g$ est localement de présentation finie.
  \end{enumerate}
  Alors, si la famille $(f_\lambda)$ est schématiquement dominante, il en est
  de même de la famille $(f'_\lambda)$.
\end{enumerate}
\end{theorem}

\oldpage[IV-3]{172}
\begin{proposition}[11.10.6]
\label{IV.11.10.6-fr}
Les notations étant celles de (11.10.5), supposons que $X$ soit un préschéma
sur un corps $k$, et qu'en posant $S=\operatorname{Spec}(k)$, on ait
$X'=X\times_SS'$, où $S'$ est un $k$-préschéma quelconque. Alors, si la famille
$(f_\lambda)$ est schématiquement dominante, il en est de même de
$(f'_\lambda)$.
\end{proposition}

\begin{corollary}[11.10.7]
\label{IV.11.10.7-fr}
Soient $X$ un préschéma sur un corps $k$, $(Z_\lambda)_{\lambda\in L}$ une
famille de $k$-préschémas géométriquement réduits sur $k$, et pour chaque
$\lambda\in L$, soit $f_\lambda:Z_\lambda\to X$ un $k$-morphisme. Désignons
par $Y$ le sous-préschéma réduit de $X$ ayant pour espace sous-jacent
l'adhérence de $\bigcup_{\lambda\in L}f_\lambda(Z_\lambda)$. Soient $k'$ une
extension de $k$, $X'$, $Z'_\lambda$,
$f'_\lambda:Z'_\lambda\to X'$ les préschémas et morphismes déduits de $X$,
$Z_\lambda$ et $f_\lambda$ par extension de la base à $k'$. Alors, si $Y'$ est
le sous-préschéma réduit de $X'$ ayant pour espace sous-jacent l'adhérence de
$\bigcup_{\lambda\in L}f'_\lambda(Z'_\lambda)$, on a
$Y'=Y\otimes_k k'$. En particulier, $Y$ est géométriquement réduit sur $k$.
\end{corollary}

\begin{proof}
Comme les $Z_\lambda$ sont réduits, les morphismes $f_\lambda$ se factorisent
en
\[
Z_\lambda\xrightarrow{g_\lambda}Y\xrightarrow{j}X,
\]
où $j$ est l'injection canonique ($\mathbf I$, 5.2.2). Il résulte alors de
(11.10.4) que $(g_\lambda)$ est une famille schématiquement dominante. Posons
$Y'_1=Y\otimes_k k'$, sous-préschéma fermé de $X'$, et soit $g'_\lambda$ le
morphisme déduit de $g_\lambda$ par extension de la base à $k'$, de sorte que
$f'_\lambda$ se factorise en
\[
Z'_\lambda\xrightarrow{g'_\lambda}Y'_1\xrightarrow{j'}X',
\]
où $j'$ est l'injection canonique. Il résulte de (11.10.6) que la famille
$(g'_\lambda)$ est schématiquement dominante. Mais par hypothèse les
$Z'_\lambda$ sont \emph{réduits}, donc ($\mathbf I$, 5.2.2) les
$g'_\lambda$ se factorisent en $Z'_\lambda\to Y'\to Y'_1$~; on conclut donc
de (11.10.2) que $Y'=Y'_1$ et $h'_\lambda=g'_\lambda$, ce qui établit le
corollaire.
\end{proof}

\begin{definition}[11.10.8]
\label{IV.11.10.8-fr}
Les notations étant celles de (11.10.5), supposons que $X$ soit un
$S$-préschéma. On dit que $(f_\lambda)$ est \emph{universellement
schématiquement dominante} (relativement à $S$) si, pour tout changement de
base $S'\to S$, la famille $(f'_\lambda)$ correspondant au changement de base
$X'=X\times_SS'\to X$, est schématiquement dominante.
\end{definition}

Lorsque $S$ est le spectre d'un corps, la prop. (11.10.6) signifie donc qu'une
famille schématiquement dominante est universellement schématiquement
dominante (relativement à $S$).

Lorsque les $f_\lambda$ sont des immersions canoniques $Z_\lambda\to X$, on
dira aussi que la famille des sous-préschémas $Z_\lambda$ est
\emph{universellement schématiquement dense dans $X$} (relativement à $S$) au
lieu de dire que la famille $(f_\lambda)$ est universellement schématiquement
dominante (relativement à $S$).

\begin{theorem}[11.10.9]
\label{IV.11.10.9-fr}
Les notations étant celles de (11.10.5), supposons que $X$ soit un
$S$-préschéma localement noethérien et que les $Z_\lambda$ soient tous
$S$-plats. Pour tout $s\in S$, soient
$X_s=X\times_S\operatorname{Spec}(k(s))$,
$(Z_\lambda)_s=Z_\lambda\times_S\operatorname{Spec}(k(s))$,
$(f_\lambda)_s=f_\lambda\times1:(Z_\lambda)_s\to X_s$. Pour que la famille
$(f_\lambda)$ soit universellement schématiquement dominante relativement à
$S$, il faut et il suffit que, pour tout $s\in S$, la famille
$((f_\lambda)_s)$ soit schématiquement dominante.
\end{theorem}

\begin{proposition}[11.10.10]
\label{IV.11.10.10-fr}
Soient $f:X\to S$ un morphisme plat localement de présentation finie, $U$ un
ouvert de $X$. Pour que $U$ soit universellement schématiquement dense dans
$X$ relativement à $S$, il faut et il suffit que, pour tout $s\in S$,
$U_s=U\cap X_s$ soit schématiquement dense dans $X_s$ (ou, ce qui revient au
même, que l'on ait $\operatorname{Ass}(\mathcal O_{X_s})\subset U_s$).
\end{proposition}
% IV3_B25_P172_END
